% ==============================================================================
%  Digital Design & ASIC/FPGA Engineering: 200 Master Interview Questions
%  LaTeX Source Document
%  Optimized for Instant Overleaf Compilation (Zero Timeout)
% ==============================================================================
\documentclass[11pt,a4paper]{article}

% -- Encoding & Typography -----------------------------------------------------
\usepackage[utf8]{inputenc}
\usepackage[T1]{fontenc}
\usepackage{lmodern}

% -- Page Geometry -------------------------------------------------------------
\usepackage[top=2.0cm, bottom=2.2cm, left=2.2cm, right=2.2cm, headheight=22pt]{geometry}

% -- Color Palette -------------------------------------------------------------
\usepackage[dvipsnames,table]{xcolor}
\definecolor{primaryblue}{HTML}{0D47A1}
\definecolor{secblue}{HTML}{1976D2}
\definecolor{ltblue}{HTML}{E3F2FD}
\definecolor{dkblue}{HTML}{0A2540}

\definecolor{advred}{HTML}{C62828}
\definecolor{ltred}{HTML}{FFEBEE}

\definecolor{modoran}{HTML}{EF6C00}
\definecolor{ltoran}{HTML}{FFF3E0}

\definecolor{easygreen}{HTML}{2E7D32}
\definecolor{ltgreen}{HTML}{E8F5E9}

\definecolor{ansgreen}{HTML}{1B5E20}
\definecolor{dkgreen}{HTML}{1B5E20}
\definecolor{ansbg}{HTML}{F1F8E9}

\definecolor{codebg}{HTML}{F8F9FA}
\definecolor{codefg}{HTML}{263238}
\definecolor{rulecol}{HTML}{90CAF9}

% -- Mathematics ---------------------------------------------------------------
\usepackage{amsmath,amssymb,mathtools}

% -- Tables --------------------------------------------------------------------
\usepackage{booktabs}
\usepackage{longtable}
\usepackage{tabularx}
\usepackage{array}
\usepackage{multirow}
\renewcommand{\arraystretch}{1.3}

% -- Lists & Spacing -----------------------------------------------------------
\usepackage{enumitem}
\usepackage{parskip}
\setlength{\parskip}{4pt}

% -- Code Listings -------------------------------------------------------------
\usepackage{listings}

\lstdefinestyle{verilogstyle}{
  language=Verilog,
  backgroundcolor=\color{codebg},
  basicstyle=\ttfamily\footnotesize\color{codefg},
  keywordstyle=\bfseries\color{primaryblue},
  commentstyle=\itshape\color{gray},
  stringstyle=\color{dkgreen},
  numberstyle=\tiny\color{gray},
  numbers=none,
  frame=single,
  rulecolor=\color{rulecol},
  breaklines=true,
  showstringspaces=false,
  tabsize=2,
  xleftmargin=1.0em,
  framexleftmargin=1.0em
}

\lstdefinestyle{tclstyle}{
  language=tcl,
  backgroundcolor=\color{codebg},
  basicstyle=\ttfamily\footnotesize\color{codefg},
  keywordstyle=\bfseries\color{Plum},
  commentstyle=\itshape\color{gray},
  stringstyle=\color{dkgreen},
  numbers=none,
  frame=single,
  rulecolor=\color{Plum!50},
  breaklines=true,
  showstringspaces=false,
  tabsize=2,
  xleftmargin=1.0em,
  framexleftmargin=1.0em,
  morekeywords={set,puts,foreach,if,proc,return,upvar,expr,array,open,
                close,read,gets,regexp,regsub,lindex,llength,lappend,
                get_ports,get_clocks,get_pins,get_cells,get_nets,
                foreach_in_collection,get_attribute,sizeof_collection,
                create_clock,create_generated_clock,set_clock_uncertainty,
                set_input_delay,set_output_delay,set_driving_cell,set_load,
                set_dont_touch,set_false_path,set_multicycle_path,
                set_clock_groups,set_max_delay,group_path,compile_ultra,
                set_clock_gating_style,set_propagated_clock,set_case_analysis}
}

% -- Fast Box Environments (mdframed default engine: 0.5s compile time) --------
\usepackage[framemethod=default]{mdframed}

% Fast Badges (Inline)
\newcommand{\badgeadv}{{\color{white}\colorbox{advred}{\bfseries\sffamily\tiny\ ADVANCED 60\%\ }}}
\newcommand{\badgemod}{{\color{white}\colorbox{modoran}{\bfseries\sffamily\tiny\ MODERATE 30\%\ }}}
\newcommand{\badgeeasy}{{\color{white}\colorbox{easygreen}{\bfseries\sffamily\tiny\ EASY 10\%\ }}}

% Fast Question Box Style
\mdfdefinestyle{qstyle}{
  backgroundcolor=ltblue!40,
  linecolor=primaryblue,
  linewidth=1.0pt,
  roundcorner=2pt,
  innertopmargin=5pt,
  innerbottommargin=5pt,
  innerleftmargin=7pt,
  innerrightmargin=7pt,
  skipabove=6pt,
  skipbelow=3pt
}

% Fast Answer Box Style
\mdfdefinestyle{ansstyle}{
  backgroundcolor=ansbg,
  linecolor=ansgreen,
  linewidth=0.9pt,
  roundcorner=2pt,
  innertopmargin=5pt,
  innerbottommargin=5pt,
  innerleftmargin=7pt,
  innerrightmargin=7pt,
  skipabove=3pt,
  skipbelow=7pt
}

% Fast Summary Box Style
\mdfdefinestyle{summstyle}{
  backgroundcolor=dkblue!5,
  linecolor=dkblue,
  linewidth=0.8pt,
  roundcorner=2pt,
  innertopmargin=5pt,
  innerbottommargin=5pt,
  innerleftmargin=7pt,
  innerrightmargin=7pt,
  skipabove=5pt,
  skipbelow=5pt
}

% Environment Definitions
\newenvironment{qbox}[2]{%
  \begin{mdframed}[style=qstyle]%
  {\bfseries\sffamily\color{primaryblue}#1 \hfill #2}\par\vspace{2pt}%
}{%
  \end{mdframed}%
}

\newenvironment{ansbox}{%
  \begin{mdframed}[style=ansstyle]%
  {\bfseries\sffamily\color{ansgreen}\small Detailed Solution \& Technical Explanation}\par\vspace{2pt}\small%
}{%
  \end{mdframed}%
}

\newenvironment{summarybox}{%
  \begin{mdframed}[style=summstyle]%
  {\bfseries\sffamily\color{dkblue}\small Core Concept Reference}\par\vspace{2pt}\small%
}{%
  \end{mdframed}%
}

% -- Headers & Footers ---------------------------------------------------------
\usepackage{fancyhdr}
\pagestyle{fancy}
\fancyhf{}
\fancyhead[L]{\small\sffamily\color{dkblue}\textbf{Digital Design Diploma} --- 200 Master Interview Questions}
\fancyhead[R]{\small\sffamily\color{dkblue}\leftmark}
\fancyfoot[C]{\small\sffamily\thepage}
\renewcommand{\headrulewidth}{0.6pt}
\renewcommand{\headrule}{\hbox to\headwidth{\color{primaryblue}\leaders\hrule height \headrulewidth\hfill}}

% -- Section Styling -----------------------------------------------------------
\usepackage{titlesec}
\titleformat{\section}
  {\normalfont\Large\bfseries\sffamily\color{primaryblue}}
  {\thesection}
  {1em}
  {}
  [\vspace{2pt}\color{primaryblue}\hrule height 1.5pt]
\titlespacing{\section}{0pt}{14pt}{8pt}

\titleformat{\subsection}
  {\normalfont\large\bfseries\sffamily\color{dkblue}}
  {\thesubsection}
  {0.8em}
  {}
\titlespacing{\subsection}{0pt}{10pt}{4pt}

% -- Hyperlinks ----------------------------------------------------------------
\usepackage[colorlinks=true, linkcolor=primaryblue,
            urlcolor=secblue, citecolor=primaryblue,
            bookmarks=true, bookmarksopen=true]{hyperref}

% ==============================================================================
\begin{document}

% -- Title Page ----------------------------------------------------------------
\begin{titlepage}
  \centering
  \vspace*{1.0cm}
  {\Huge\bfseries\sffamily\color{primaryblue} Digital Design \& ASIC/FPGA Engineering\par}
  \vspace{0.4cm}
  {\Large\bfseries\sffamily\color{secblue} 200 Master Technical Interview Questions \& Solutions\par}
  \vspace{0.8cm}
  
  \begin{mdframed}[style=qstyle]
    \centering\small\sffamily
    \textbf{Target Roles:} ASIC/FPGA Digital Design Engineer $\cdot$ RTL Verification Engineer $\cdot$ Synthesis \& STA Engineer $\cdot$ DFT Engineer $\cdot$ Physical Design (PnR) Engineer.\\[4pt]
    \textbf{Curriculum Coverage:} Verilog HDL $\cdot$ TCL $\cdot$ SpyGlass Lint/CDC $\cdot$ Design Compiler $\cdot$ PrimeTime STA $\cdot$ Low Power (UPF) $\cdot$ Asynchronous FIFO $\cdot$ Formal Verification (LEC/SVA) $\cdot$ DFT/ATPG/JTAG $\cdot$ GLS $\cdot$ PnR $\cdot$ FPGA.
  \end{mdframed}

  \vspace{0.8cm}

  \begin{table}[h!]
    \centering\small
    \begin{tabular}{p{3.2cm} c c p{6.8cm}}
      \toprule
      \textbf{Difficulty Level} & \textbf{Percentage} & \textbf{Questions} & \textbf{Target Competency} \\
      \midrule
      \textbf{Advanced} \badgeadv & \textbf{60\%} & 120 & Senior / Lead ASIC Architecture \& Signoff \\
      \textbf{Moderate} \badgemod & \textbf{30\%} & 60 & Mid-Level Engineering \& Tool Automation \\
      \textbf{Easy} \badgeeasy & \textbf{10\%} & 20 & Fundamental Digital Logic Principles \\
      \midrule
      \textbf{Total} & \textbf{100\%} & \textbf{200} & \textbf{Comprehensive Interview Preparation} \\
      \bottomrule
    \end{tabular}
  \end{table}

  \vfill
  {\small\sffamily\color{gray} Compiled for Top Semiconductor Industry Interviews (Apple, NVIDIA, Qualcomm, Intel, AMD, ARM, Synopsys, Cadence)\par}
  \vspace*{0.5cm}
\end{titlepage}

\newpage
\tableofcontents
\newpage

% ==============================================================================
\section{Verilog HDL \& Digital System Architecture}
% ==============================================================================

\begin{qbox}{Q1: Verilog Blocking vs Non-Blocking Assignment}{\badgeeasy}
In synthesizable Verilog RTL for synchronous sequential logic (flip-flops), which assignment operator should be used, and in which IEEE simulation event region is its LHS update scheduled?
\begin{enumerate}[label=\textbf{\Alph*.}]
  \item \texttt{=}, Active region
  \item \texttt{<=}, NBA (Non-Blocking Assignment) update region
  \item \texttt{=}, Postponed region
  \item \texttt{<=}, Observed region
\end{enumerate}
\end{qbox}
\begin{ansbox}
\textbf{Correct Answer: B}

In sequential \texttt{always @(posedge clk)} blocks, non-blocking assignments (\texttt{<=}) evaluate the RHS expression in the Active region and schedule the LHS variable update into the \textbf{NBA update region} of the IEEE 1364/1800 event queue. This guarantees deterministic behavior and eliminates race conditions between communicating sequential registers.
\end{ansbox}

\begin{qbox}{Q2: Register File Address Calculation}{\badgeeasy}
For a dual-port register file with parameterized depth $D$, what is the correct standard Verilog-2001 system function to calculate the exact address bus width?
\begin{enumerate}[label=\textbf{\Alph*.}]
  \item \texttt{input [log2(D)-1:0] addr;}
  \item \texttt{input [\$clog2(D)-1:0] addr;}
  \item \texttt{input [ln(D)/ln(2):0] addr;}
  \item \texttt{input [\$log(D):0] addr;}
\end{enumerate}
\end{qbox}
\begin{ansbox}
\textbf{Correct Answer: B}

\texttt{\$clog2} is a built-in standard system function supported in Verilog-2001 and SystemVerilog that computes the ceiling of $\log_2(x)$, giving the exact number of address bits required to uniquely address $D$ memory locations.
\end{ansbox}

\begin{qbox}{Q3: 3-Process FSM Registered Outputs}{\badgemod}
Why is a 3-process (or registered-output) Moore/Mealy FSM preferred in high-performance ASIC design over a 2-process combinational-output FSM?
\begin{enumerate}[label=\textbf{\Alph*.}]
  \item It uses fewer flip-flops and reduces dynamic power.
  \item It completely eliminates combinational output glitches and provides a full clock period of timing budget to downstream modules.
  \item It eliminates the need for state encoding.
  \item It reduces the total number of FSM states by half.
\end{enumerate}
\end{qbox}
\begin{ansbox}
\textbf{Correct Answer: B}

When FSM outputs are decoded combinationally from state registers, intermediate decoding hazards cause glitches that consume dynamic power and eat into the downstream setup timing budget. Registering the outputs ensures glitch-free signals and creates clean register-to-register paths.
\end{ansbox}

\begin{qbox}{Q4: Inferred Latches in Combinational Blocks}{\badgemod}
Consider the following combinational block:
\begin{lstlisting}[style=verilogstyle]
always @(*) begin
    case (mode)
        2'b00: y = a + b;
        2'b01: y = a - b;
        2'b10: y = a & b;
    endcase
end
\end{lstlisting}
What hardware structure will logic synthesis infer, and why?
\begin{enumerate}[label=\textbf{\Alph*.}]
  \item Pure 4-to-1 Multiplexer with unused input tied to ground.
  \item Combinational AND-OR tree with \texttt{y} defaulting to zero.
  \item A transparent D-Latch on \texttt{y} because branch \texttt{2'b11} is missing and \texttt{y} must retain its previous value.
  \item A D Flip-Flop clocked by the \texttt{mode} signal.
\end{enumerate}
\end{qbox}
\begin{ansbox}
\textbf{Correct Answer: C}

In an \texttt{always @(*)} block, if a variable is not assigned a value in all possible execution branches (e.g. missing \texttt{2'b11} without a \texttt{default:} statement), the synthesis tool infers a level-sensitive latch to preserve its state when \texttt{mode == 2'b11}.
\end{ansbox}

\begin{qbox}{Q5: Arithmetic vs Logical Shift Operators}{\badgemod}
What is the key functional difference between \texttt{>>>} (arithmetic right shift) and \texttt{>>} (logical right shift) in Verilog?
\begin{enumerate}[label=\textbf{\Alph*.}]
  \item \texttt{>>>} operates only on unsigned variables, filling MSBs with 0.
  \item \texttt{>>>} sign-extends the MSB if and only if the operand is explicitly declared as \texttt{signed}.
  \item \texttt{>>>} rotates bits through a carry register.
  \item Both operators always perform identical sign extension.
\end{enumerate}
\end{qbox}
\begin{ansbox}
\textbf{Correct Answer: B}

In Verilog-2001, \texttt{>>>} will only perform sign extension (filling empty MSB positions with the sign bit) if the operand being shifted is explicitly declared with the \texttt{signed} keyword. If the operand is unsigned, \texttt{>>>} behaves identically to \texttt{>>} by zero-filling.
\end{ansbox}

\begin{qbox}{Q6: LFSR Maximal Repetition Period}{\badgemod}
An $n$-bit Linear Feedback Shift Register (LFSR) uses a primitive polynomial for maximal-length pseudo-random sequence generation. What is its exact repetition period in clock cycles?
\begin{enumerate}[label=\textbf{\Alph*.}]
  \item $2^n$
  \item $2^n - 1$
  \item $2^{n-1}$
  \item $n^2 - 1$
\end{enumerate}
\end{qbox}
\begin{ansbox}
\textbf{Correct Answer: B}

A standard XOR-based LFSR has an illegal state of all zeros ($00\dots0$), which produces a lockup condition. Therefore, its maximum cycle length is $2^n - 1$.
\end{ansbox}

\begin{qbox}{Q7: 50\% Duty Cycle Odd Clock Divider}{\badgemod}
To generate a clean 50\% duty cycle clock divided by an \textbf{odd} integer $N$ (e.g., $N=3$ or $N=5$) from a reference clock without an internal PLL, what technique is required?
\begin{enumerate}[label=\textbf{\Alph*.}]
  \item Use an edge-triggered counter on positive edges only and toggle at $N/2$.
  \item Generate two intermediate clocks: one using positive clock edges and another using negative clock edges, each active for $(N-1)/2$ cycles, and OR (or AND) them together.
  \item Divide the clock by $2N$ and use a frequency multiplier.
  \item Use an asynchronous delay line with an inverter ring.
\end{enumerate}
\end{qbox}
\begin{ansbox}
\textbf{Correct Answer: B}

Because an odd number cannot be divided into integers on positive edges alone, standard 50\% duty cycle odd dividers generate two phase-shifted pulses (one triggered on \texttt{posedge ref\_clk} and the other on \texttt{negedge ref\_clk} with pulse width $(N-1)/2$) and combine them with an OR/AND gate.
\end{ansbox}

\begin{qbox}{Q8: Signed ALU Overflow Condition}{\badgemod}
In a 16-bit signed ALU performing two's complement addition $S = A + B$, which Boolean equation correctly flags signed arithmetic overflow?
\begin{enumerate}[label=\textbf{\Alph*.}]
  \item $V = A[15] \oplus B[15]$
  \item $V = (A[15] == B[15]) \ \&\ \&\ (S[15] \ne A[15])$
  \item $V = C_{out} \oplus A[15]$
  \item $V = S[15] \ \&\ (A[15] \mid B[15])$
\end{enumerate}
\end{qbox}
\begin{ansbox}
\textbf{Correct Answer: B}

In two's complement addition, overflow occurs if and only if adding two numbers of the same sign produces a result with the opposite sign (i.e. positive + positive = negative, or negative + negative = positive).
\end{ansbox}

\begin{qbox}{Q9: Verilog Event Queue Process Scheduling Race}{\badgeadv}
Examine the following two \texttt{always} blocks:
\begin{lstlisting}[style=verilogstyle]
// Block 1
always @(posedge clk) begin
    q1 = d;
end
// Block 2
always @(posedge clk) begin
    q2 = q1;
end
\end{lstlisting}
What is the value of \texttt{q2} after a clock edge where \texttt{d} changes from \texttt{0} to \texttt{1} (assuming initially \texttt{q1=0, q2=0}), and why is this non-deterministic?
\begin{enumerate}[label=\textbf{\Alph*.}]
  \item \texttt{q2} is guaranteed to be \texttt{0} because both blocks trigger simultaneously.
  \item \texttt{q2} is guaranteed to be \texttt{1} because Block 1 evaluates first in lexical order.
  \item Non-deterministic (Race condition): If Block 1 executes first, \texttt{q2} captures \texttt{1} (new value); if Block 2 executes first, \texttt{q2} captures \texttt{0} (old value).
  \item The simulator reports a compile-time syntax error.
\end{enumerate}
\end{qbox}
\begin{ansbox}
\textbf{Correct Answer: C}

Because blocking assignments (\texttt{=}) immediately update the LHS in the Active event region, the result depends entirely on the simulator's internal process scheduling order. Non-blocking assignments (\texttt{<=}) must always be used for sequential logic to prevent this race condition.
\end{ansbox}

\begin{qbox}{Q10: Multi-Driven Net Simulation vs Synthesis}{\badgeadv}
Consider two separate procedural blocks in the same module assigning to the same variable:
\begin{lstlisting}[style=verilogstyle]
always @(posedge clk1) data_out <= in1;
always @(posedge clk2) data_out <= in2;
\end{lstlisting}
How do the simulator and logic synthesis tool handle this construct?
\begin{enumerate}[label=\textbf{\Alph*.}]
  \item Simulator treats it as wired-OR; Synthesis infers a clean 2-to-1 MUX.
  \item Simulator resolves it according to the last executed block or drives \texttt{X}; Synthesis tool halts with a \textbf{Multi-Driven Net / Multi-Source Error} (\texttt{ELAB-900} / \texttt{W422}).
  \item Both simulator and synthesis tool automatically insert a clock domain arbiter.
  \item Simulator issues a fatal error; Synthesis infers a dual-clock flip-flop.
\end{enumerate}
\end{qbox}
\begin{ansbox}
\textbf{Correct Answer: B}

Multiple procedural blocks driving the same variable create multi-driver conflicts. In simulation, the last executed block overwrites the variable. Logic synthesis cannot map multiple unsynchronized clock drivers to a single physical register without multiplexing logic, triggering a fatal elaboration error.
\end{ansbox}

\begin{qbox}{Q11: Generate Loops and Scope Hierarchy}{\badgeadv}
In SystemVerilog/Verilog-2001, when instantiating parameterized processing elements inside a \texttt{genvar} loop:
\begin{lstlisting}[style=verilogstyle]
genvar i;
generate
    for (i = 0; i < 4; i = i + 1) begin : gen_pe
        PE u_pe (.clk(clk), .din(data[i]), .dout(out[i]));
    end
endgenerate
\end{lstlisting}
What is the full hierarchical path to the internal register \texttt{acc\_reg} inside the 3rd processing element (index 2)?
\begin{enumerate}[label=\textbf{\Alph*.}]
  \item \texttt{top.u\_pe[2].acc\_reg}
  \item \texttt{top.gen\_pe[2].u\_pe.acc\_reg}
  \item \texttt{top.gen\_pe.u\_pe[2].acc\_reg}
  \item \texttt{top.u\_pe.gen\_pe(2).acc\_reg}
\end{enumerate}
\end{qbox}
\begin{ansbox}
\textbf{Correct Answer: B}

Named \texttt{generate} blocks (\texttt{begin : gen\_pe}) create a generate scope array indexed by the \texttt{genvar} loop index. The correct hierarchy is \texttt{<module>.<generate\_block\_name>[<index>].<instance\_name>.<signal>}.
\end{ansbox}

\begin{qbox}{Q12: Hazards of full\_case and parallel\_case Pragmas}{\badgeadv}
Why are Synopsys synthesis directives \texttt{// synthesis full\_case} and \texttt{// synthesis parallel\_case} considered hazardous in modern ASIC design?
\begin{enumerate}[label=\textbf{\Alph*.}]
  \item They increase the cell count and silicon area by over 50\%.
  \item They disable all static timing analysis checks.
  \item They cause \textbf{Simulation-Synthesis Mismatches} because the RTL simulator ignores comments while the synthesis tool alters its logic optimization.
  \item They force the placement tool to use only high-Vt cells.
\end{enumerate}
\end{qbox}
\begin{ansbox}
\textbf{Correct Answer: C}

Synthesis pragmas inside comments (\texttt{// synthesis ...}) are invisible to standard RTL simulators (which evaluate the \texttt{case} statement strictly with priority or incomplete coverage). The synthesis tool synthesizes parallel multiplexers or don't-care states, resulting in silicon that behaves differently from the verified RTL simulation.
\end{ansbox}

\begin{qbox}{Q13: UART Receiver 16x Oversampling Timing}{\badgeadv}
In a UART receiver operating with a 16x oversampling clock (\texttt{baud\_16x\_clk}), why is the start bit sampled at the 8th tick after the falling edge of \texttt{RX} instead of the 1st tick?
\begin{enumerate}[label=\textbf{\Alph*.}]
  \item To allow the line to charge up its capacitance.
  \item To sample at the theoretical midpoint (center) of the start bit, maximizing tolerance against clock jitter and baud rate drift, and to filter false glitches.
  \item To calculate the parity bit in advance.
  \item Because UART hardware requires 8 bits of preamble.
\end{enumerate}
\end{qbox}
\begin{ansbox}
\textbf{Correct Answer: B}

16x oversampling checks for a valid falling edge at tick 0, verifies that \texttt{RX} remains low at tick 8 (the center of the start bit to reject high-frequency noise spikes), and subsequently samples each data bit at intervals of 16 ticks (the center of each data bit).
\end{ansbox}

\begin{qbox}{Q14: Parameterized Register File Write-Collision Handling}{\badgeadv}
Consider a parameterized register file with synchronous write and combinational read:
\begin{lstlisting}[style=verilogstyle]
always @(posedge clk) begin
    if (wr_en) mem[wr_addr] <= wr_data;
end
assign rd_data = mem[rd_addr];
\end{lstlisting}
What happens when \texttt{wr\_addr == rd\_addr} and \texttt{wr\_en == 1} simultaneously?
\begin{enumerate}[label=\textbf{\Alph*.}]
  \item \texttt{rd\_data} immediately reflects the new \texttt{wr\_data} combinationally in the same cycle (Write-Through / Read-First).
  \item \texttt{rd\_data} outputs \texttt{X} due to memory collision.
  \item \texttt{rd\_data} outputs the old value stored at \texttt{rd\_addr} for the current cycle.
  \item The register file enters a deadlock state.
\end{enumerate}
\end{qbox}
\begin{ansbox}
\textbf{Correct Answer: A}

Because the read port is a continuous assignment (\texttt{assign rd\_data = mem[rd\_addr]}), as soon as \texttt{mem[wr\_addr]} is updated or if bypassed, the combinational read reflects the memory contents. However, to guarantee exact write-forwarding without intra-cycle races, explicit forwarding logic \texttt{assign rd\_data = (wr\_en \&\& wr\_addr == rd\_addr) ? wr\_data : mem[rd\_addr];} is used.
\end{ansbox}

\begin{qbox}{Q15: One-Hot vs Binary FSM Encoding Trade-Offs}{\badgeadv}
When targeting high-frequency ASIC designs with a large number of states (e.g. 20 states), why is \textbf{One-Hot} state encoding often preferred over \textbf{Binary} encoding?
\begin{enumerate}[label=\textbf{\Alph*.}]
  \item One-hot encoding minimizes the number of state flip-flops required ($N = \lceil\log_2 S\rceil$).
  \item One-hot encoding simplifies next-state and output combinational decoding logic to simple shallow gates (1-2 logic levels), maximizing clock frequency at the expense of more flip-flops.
  \item One-hot encoding produces zero dynamic switching power.
  \item One-hot encoding eliminates the need for reset logic.
\end{enumerate}
\end{qbox}
\begin{ansbox}
\textbf{Correct Answer: B}

Binary encoding requires complex multi-level decoding trees ($O(\log_2 S)$ depth) to evaluate next-state logic. One-hot encoding dedicates 1 flip-flop per state, allowing each state transition to be decoded using a simple 1-level AND/OR gate, dramatically reducing critical path delay for timing closure.
\end{ansbox}

\begin{qbox}{Q16: Latch-Based Integrated Clock Gating (ICG) Architecture}{\badgeadv}
In an Integrated Clock Gating (ICG) cell for an active-high clock, why must the enable signal pass through a \textbf{negative-edge triggered latch} before feeding the AND gate?
\begin{enumerate}[label=\textbf{\Alph*.}]
  \item To synchronize the enable signal to an asynchronous domain.
  \item To hold the enable signal stable during the high phase of \texttt{CLK}, preventing positive clock glitches, hazard pulses, and shortened clock pulses.
  \item To invert the clock polarity for negative-edge flip-flops.
  \item To reduce leakage current during scan testing.
\end{enumerate}
\end{qbox}
\begin{ansbox}
\textbf{Correct Answer: B}

A simple AND gate without a latch produces clock glitches if \texttt{enable} toggles while \texttt{CLK} is high. By placing a latch that captures \texttt{enable} when \texttt{CLK} is low (transparent when \texttt{CLK=0}, latched when \texttt{CLK=1}), any transition on \texttt{enable} while \texttt{CLK=1} is blocked, guaranteeing a glitch-free gated clock output.
\end{ansbox}

\begin{qbox}{Q17: SystemVerilog unique case vs priority case}{\badgeadv}
What is the key verification and synthesis difference between \texttt{unique case} and \texttt{priority case} in SystemVerilog?
\begin{enumerate}[label=\textbf{\Alph*.}]
  \item \texttt{unique case} checks for full coverage only; \texttt{priority case} checks for overlap only.
  \item \texttt{unique case} asserts that exactly one condition is true at any time (flags a run-time warning on overlap or no match); \texttt{priority case} enforces evaluation order and only flags a warning if no condition matches.
  \item \texttt{priority case} synthesizes to a parallel MUX; \texttt{unique case} synthesizes to a priority encoder.
  \item They are strictly synonyms for backward compatibility with Verilog-95.
\end{enumerate}
\end{qbox}
\begin{ansbox}
\textbf{Correct Answer: B}

SystemVerilog \texttt{unique} instructs both simulator and synthesis tools that case arms are mutually exclusive (parallel) and fully specified; it produces run-time assertion warnings if multiple branches match or none match. \texttt{priority} enforces priority order and only checks that at least one branch matches.
\end{ansbox}

\begin{qbox}{Q18: Glitch Generation in Asynchronous Ripple Counters}{\badgeadv}
Why are asynchronous ripple counters (where the output of stage $N$ clocks stage $N+1$) strictly forbidden in high-reliability synchronous ASIC design?
\begin{enumerate}[label=\textbf{\Alph*.}]
  \item They require twice the silicon area of synchronous counters.
  \item Cumulative clock-to-Q propagation delays through the ripple chain create severe output skew and intermediate combinational decoding glitches, causing false triggers in downstream logic.
  \item They cannot count backwards.
  \item They are not supported by Verilog simulators.
\end{enumerate}
\end{qbox}
\begin{ansbox}
\textbf{Correct Answer: B}

In ripple counters, each flip-flop is clocked by the previous flip-flop, causing accumulated skew ($N \times t_{cq}$). When decoding multi-bit outputs (e.g., detecting state 4 \texttt{3'b100} from \texttt{3'b011}), temporary transition states like \texttt{3'b010} and \texttt{3'b000} appear briefly, generating massive glitches that corrupt synchronous clock tree analysis and DFT scan chains.
\end{ansbox}

\begin{qbox}{Q19: Gray Code Mathematical Conversion}{\badgeadv}
Given a 4-bit binary value $B = [B_3, B_2, B_1, B_0]$, what is the exact synthesizable expression to compute its corresponding Gray code $G = [G_3, G_2, G_1, G_0]$?
\begin{enumerate}[label=\textbf{\Alph*.}]
  \item \texttt{assign G = B \^{} (B >> 1);}
  \item \texttt{assign G = B \& (B << 1);}
  \item \texttt{assign G = B + 1'b1;}
  \item \texttt{assign G = \textasciitilde B \^{} (B >> 1);}
\end{enumerate}
\end{qbox}
\begin{ansbox}
\textbf{Correct Answer: A}

Gray code conversion operates by XORing the binary word with its right-shifted version: $G_i = B_i \oplus B_{i+1}$, with $G_{MSB} = B_{MSB}$. In Verilog: \texttt{G = B \^{} (B >> 1)}.
\end{ansbox}

\begin{qbox}{Q20: System Controller Master FSM Transaction Sequencing}{\badgeadv}
In the diploma's Master System Controller FSM, when a UART frame arrives requesting an ALU operation with operands, what is the correct state transition sequence?
\begin{enumerate}[label=\textbf{\Alph*.}]
  \item \texttt{IDLE -> WRITE\_REG -> EXEC\_ALU -> SEND\_UART -> IDLE}
  \item \texttt{IDLE -> READ\_CMD -> READ\_OPERANDS -> START\_ALU -> WAIT\_ALU\_VALID -> SEND\_UART\_TX -> IDLE}
  \item \texttt{IDLE -> ALU\_BUSY -> IDLE}
  \item \texttt{IDLE -> RESET -> EXEC -> IDLE}
\end{enumerate}
\end{qbox}
\begin{ansbox}
\textbf{Correct Answer: B}

The system controller FSM decodes the framing command, receives multi-byte operands via UART RX into temporary/register storage, asserts \texttt{ALU\_ENABLE}, waits for \texttt{ALU\_OUT\_VALID}, and subsequently pulses \texttt{UART\_TX\_START} to transmit the result packet back through the asynchronous FIFO/UART TX interface.
\end{ansbox}

\begin{qbox}{Q21: Synthesizable Verilog function vs task}{\badgeadv}
Which of the following statements correctly distinguishes a synthesizable Verilog \texttt{function} from a \texttt{task}?
\begin{enumerate}[label=\textbf{\Alph*.}]
  \item Functions can contain \texttt{\#} delay statements; tasks cannot.
  \item Functions execute in zero simulation time, return a single value directly, and cannot contain timing controls; tasks can have output/inout arguments and consume simulation time.
  \item Functions can invoke tasks, but tasks cannot invoke functions.
  \item Functions can drive global clock signals; tasks cannot.
\end{enumerate}
\end{qbox}
\begin{ansbox}
\textbf{Correct Answer: B}

In Verilog HDL, a \texttt{function} represents pure combinational logic: it executes in zero simulation time, cannot contain time-consuming statements (\texttt{\#}, \texttt{@}, \texttt{wait}), must have at least one input, and returns a single value. A \texttt{task} can consume time and return multiple values via \texttt{output}/\texttt{inout} arguments.
\end{ansbox}

\begin{qbox}{Q22: Structural Race Condition in Clock Multiplexing}{\badgeadv}
What is the fatal flaw in multiplexing two asynchronous clocks using a simple combinational 2-to-1 MUX (\texttt{assign clk\_out = sel ? clk\_fast : clk\_slow;})?
\begin{enumerate}[label=\textbf{\Alph*.}]
  \item The synthesis tool converts it to an inverter chain.
  \item Switching \texttt{sel} asynchronously produces narrow runt pulses, clock glitches, and severe setup/hold violations on all downstream flip-flops.
  \item The MUX consumes excessive leakage power.
  \item The clock frequency is halved.
\end{enumerate}
\end{qbox}
\begin{ansbox}
\textbf{Correct Answer: B}

Switching between unsynchronized clocks with a simple MUX creates sliver/runt clock pulses if \texttt{sel} toggles while either clock is high. A \textbf{Glitch-Free Clock Switch (GFCS)} circuit with cross-coupled negative-edge synchronizers and feedback suppression is required.
\end{ansbox}

% ==============================================================================
\section{TCL Scripting for EDA Automation}
% ==============================================================================

\begin{qbox}{Q23: TCL String \& Variable Substitution Rules}{\badgeeasy}
In TCL scripting for EDA tools, what is the exact output of \texttt{set a 5; set b \{\$a + 2\}; puts \$b}?
\begin{enumerate}[label=\textbf{\Alph*.}]
  \item \texttt{7}
  \item \texttt{5 + 2}
  \item \texttt{\$a + 2}
  \item \texttt{Error: invalid syntax}
\end{enumerate}
\end{qbox}
\begin{ansbox}
\textbf{Correct Answer: C}

Curly braces \texttt{\{\}} in TCL suppress all variable (\texttt{\$}) and command (\texttt{[]}) substitutions, treating the enclosed text as a literal string. Double quotes \texttt{""} or \texttt{expr} are required for evaluation.
\end{ansbox}

\begin{qbox}{Q24: EDA Tool Collections vs Standard TCL Lists}{\badgemod}
In Synopsys Design Compiler or PrimeTime, why will \texttt{foreach port [get\_ports *] \{ puts \$port \}} fail or produce unexpected results?
\begin{enumerate}[label=\textbf{\Alph*.}]
  \item The \texttt{puts} command does not work inside loops.
  \item \texttt{get\_ports} returns an opaque tool \textbf{Collection Pointer} (handle), not a standard TCL list; \texttt{foreach\_in\_collection} must be used.
  \item \texttt{get\_ports} only accepts single port names.
  \item Wildcards \texttt{*} are forbidden in Synopsys tools.
\end{enumerate}
\end{qbox}
\begin{ansbox}
\textbf{Correct Answer: B}

Synopsys tools use internal memory-efficient structures called \textbf{collections}. Standard TCL commands like \texttt{llength} and \texttt{foreach} cannot parse collection pointers directly. The script must use \texttt{foreach\_in\_collection} or convert the collection using \texttt{get\_object\_name [get\_ports *]}.
\end{ansbox}

\begin{qbox}{Q25: Querying Violating Timing Paths in PrimeTime}{\badgemod}
Which TCL command correctly queries all timing paths whose setup slack is less than \texttt{0.0} in PrimeTime?
\begin{enumerate}[label=\textbf{\Alph*.}]
  \item \texttt{get\_cells -filter "is\_sequential == true \&\& slack < 0.0"}
  \item \texttt{get\_timing\_paths -slack\_lesser\_than 0.0 -max\_paths 100}
  \item \texttt{filter\_collection [get\_pins */CLK] "slack < 0"}
  \item Both B and a properly filtered \texttt{get\_timing\_paths} are standard EDA methods.
\end{enumerate}
\end{qbox}
\begin{ansbox}
\textbf{Correct Answer: D}

In PrimeTime, \texttt{get\_timing\_paths -slack\_lesser\_than 0.0} directly queries violating paths, while \texttt{get\_cells -filter} queries cell properties. Both leverage EDA collection filtering.
\end{ansbox}

\begin{qbox}{Q26: Regular Expressions in TCL for Report Parsing}{\badgemod}
To extract the worst negative slack (WNS) value from a report line \texttt{Slack: -0.45 ns (VIOLATED)} into a TCL variable \texttt{wns\_val}, which command is correct?
\begin{enumerate}[label=\textbf{\Alph*.}]
  \item \texttt{regexp \{Slack:\textbackslash s+([-\textbackslash d\textbackslash .]+)\textbackslash s+ns\} \$line match wns\_val}
  \item \texttt{set wns\_val [string index \$line 8 12]}
  \item \texttt{regsub \{Slack\} \$line wns\_val}
  \item \texttt{scan \$line "Slack: \%s" wns\_val}
\end{enumerate}
\end{qbox}
\begin{ansbox}
\textbf{Correct Answer: A}

The regex \texttt{Slack:\textbackslash s+([-\textbackslash d\textbackslash .]+)\textbackslash s+ns} searches for the literal word "Slack:", skips whitespace, captures the negative floating-point number into submatch group 1 (\texttt{wns\_val}), and confirms trailing "ns".
\end{ansbox}

\begin{qbox}{Q27: Pass-By-Reference in TCL using upvar}{\badgemod}
What is the primary function of \texttt{upvar 1 var\_name local\_alias} inside a TCL procedure?
\begin{enumerate}[label=\textbf{\Alph*.}]
  \item It promotes a local variable to a global environment across all files.
  \item It binds a variable in the caller's stack frame to a local name, enabling true pass-by-reference.
  \item It increments the integer value of \texttt{var\_name} by 1.
  \item It imports environment variables from the Linux shell.
\end{enumerate}
\end{qbox}
\begin{ansbox}
\textbf{Correct Answer: B}

TCL procedures pass arguments by value. \texttt{upvar 1 <caller\_var> <local\_var>} creates an alias link to the caller's scope (1 level up the call stack), allowing the procedure to modify arrays or large data structures in place without copying.
\end{ansbox}

\begin{qbox}{Q28: Automated SDC Constraint Generator Budgeting}{\badgeadv}
Consider the following automated SDC constraint generator script:
\begin{lstlisting}[style=tclstyle]
set CLK_PERIOD 10.0
create_clock -name SYS_CLK -period $CLK_PERIOD [get_ports clk]
set_input_delay  [expr $CLK_PERIOD * 0.25] -clock SYS_CLK [get_ports -filter "direction == in && name != clk"]
set_output_delay [expr $CLK_PERIOD * 0.30] -clock SYS_CLK [get_ports -filter "direction == out"]
\end{lstlisting}
What timing budget is left inside the chip for internal combinational logic between input ports and output ports on a purely feedthrough path?
\begin{enumerate}[label=\textbf{\Alph*.}]
  \item $10.0\,\text{ns}$
  \item $5.5\,\text{ns}$
  \item $4.5\,\text{ns}$
  \item $7.5\,\text{ns}$
\end{enumerate}
\end{qbox}
\begin{ansbox}
\textbf{Correct Answer: C}

The total clock period is $10.0\,\text{ns}$. External input delay consumes $25\% = 2.5\,\text{ns}$. External downstream output delay consumes $30\% = 3.0\,\text{ns}$. The remaining allowable internal datapath budget is $10.0 - 2.5 - 3.0 = 4.5\,\text{ns}$ (assuming ideal clock edges).
\end{ansbox}

\begin{qbox}{Q29: Accumulating Pin Capacitances on an Internal Net}{\badgeadv}
In Synopsys Design Compiler, which TCL script snippet accurately calculates the total pin capacitance connected to an internal net \texttt{n\_data}?
\begin{enumerate}[label=\textbf{\Alph*.}]
  \item \texttt{set cap [get\_attribute [get\_nets n\_data] capacitance]}
  \item \texttt{set cap [get\_attribute [get\_pins -leaf -of\_objects [get\_nets n\_data]] capacitance]}
  \item \texttt{set cap 0; foreach\_in\_collection pin [get\_pins -of\_objects [get\_nets n\_data] -filter "direction == in"] \{ set cap [expr \$cap + [get\_attribute \$pin capacitance]] \}}
  \item \texttt{set cap [sizeof\_collection [get\_nets n\_data]]}
\end{enumerate}
\end{qbox}
\begin{ansbox}
\textbf{Correct Answer: C}

The total load on a net consists of the sum of the input capacitances of all driven load pins (plus wire load). Iterating through the input pins of the net and querying their \texttt{capacitance} attribute using \texttt{get\_attribute} accurately accumulates the pin load.
\end{ansbox}

\begin{qbox}{Q30: Safe Handling of Unconstrained Endpoints in Signoff}{\badgeadv}
You are writing a signoff checking script in PrimeTime. Which command correctly verifies that \textbf{no unconstrained timing endpoints} exist in the design?
\begin{enumerate}[label=\textbf{\Alph*.}]
  \item \texttt{report\_constraint -all\_violators}
  \item \texttt{check\_timing -include \{unconstrained\_endpoints no\_clock\}}
  \item \texttt{report\_timing -delay\_type max -path\_type full\_clock}
  \item \texttt{get\_timing\_paths -slack\_lesser\_than 0.0}
\end{enumerate}
\end{qbox}
\begin{ansbox}
\textbf{Correct Answer: B}

\texttt{report\_timing} and \texttt{report\_constraint} only evaluate constrained paths. If a register has a missing clock or unconstrained input/output, it is silently ignored by slack reports. \texttt{check\_timing} explicitly identifies missing clocks, unconstrained endpoints, and loops.
\end{ansbox}

\begin{qbox}{Q31: Traversing Hierarchical Leaf Cells in TCL}{\badgeadv}
When building an automated linting or cell-replacement script, how do you traverse all leaf-level instantiated library cells inside a hierarchical sub-block \texttt{u\_sys}?
\begin{enumerate}[label=\textbf{\Alph*.}]
  \item \texttt{get\_cells u\_sys/*}
  \item \texttt{get\_cells -hierarchical -filter "is\_hierarchical == false" u\_sys/*}
  \item \texttt{get\_objects -all u\_sys}
  \item \texttt{get\_pins -of\_objects u\_sys}
\end{enumerate}
\end{qbox}
\begin{ansbox}
\textbf{Correct Answer: B}

The \texttt{-hierarchical} flag recurses down all child levels, and filtering by \texttt{"is\_hierarchical == false"} excludes parent boundary blocks, returning only fundamental leaf standard cells (AND, Flip-Flops, MUXes).
\end{ansbox}

\begin{qbox}{Q32: Error Trapping in Batch Synthesis Scripts}{\badgeadv}
In an automated continuous-integration (CI) synthesis script, why is \texttt{if \{[catch \{compile\_ultra\} err\_msg]\} \{ puts "Synthesis Failed: \$err\_msg"; exit 1 \}} used?
\begin{enumerate}[label=\textbf{\Alph*.}]
  \item \texttt{catch} runs the command in a separate background Linux thread.
  \item \texttt{catch} intercepts fatal exceptions and errors from the EDA command, preventing abrupt script termination and allowing clean log capture, error formatting, and return codes.
  \item \texttt{catch} forces the tool to ignore timing violations.
  \item \texttt{catch} accelerates tool execution speed.
\end{enumerate}
\end{qbox}
\begin{ansbox}
\textbf{Correct Answer: B}

In TCL, unhandled command failures abort script execution immediately. The \texttt{catch} command executes the block, traps any abnormal return or syntax/runtime error, stores the error text into \texttt{err\_msg}, and returns a non-zero exit code for programmatic handling.
\end{ansbox}

\begin{qbox}{Q33: Storing Timing Metrics in TCL Dictionaries}{\badgeadv}
In PrimeTime, to construct a custom nested dictionary of startpoint-to-endpoint slacks: \texttt{dict set timing\_db \$sp \$ep \$slack}. What is the primary benefit of storing timing data in TCL 8.5+ \texttt{dict} structures versus standard arrays?
\begin{enumerate}[label=\textbf{\Alph*.}]
  \item Dictionaries are pass-by-value and can be nested, passed directly to procedures, and serialized easily without global scope collisions.
  \item Dictionaries run on GPU memory.
  \item Dictionaries automatically close setup timing violations.
  \item Dictionaries convert netlists to Verilog automatically.
\end{enumerate}
\end{qbox}
\begin{ansbox}
\textbf{Correct Answer: A}

TCL \texttt{dict} values are pure first-class values. Unlike classical TCL \texttt{array}s (which are pass-by-name structures that cannot be returned from procedures or nested easily), \texttt{dict} objects support multi-level hierarchical nesting and clean functional data manipulation.
\end{ansbox}

\begin{qbox}{Q34: Automated EDA Log Parsing Matrix}{\badgeadv}
In a regression synthesis script running multiple runs with varying clock periods, what is the best practice for capturing runtime performance metrics?
\begin{enumerate}[label=\textbf{\Alph*.}]
  \item Manually opening each GUI session.
  \item Using \texttt{cputime} and \texttt{mem} attributes via \texttt{get\_attribute [current\_design] ...} combined with TCL file I/O to output a consolidated \texttt{.csv} summary matrix.
  \item Rerunning the tool with \texttt{+define+DEBUG}.
  \item Relying on standard terminal printouts without file redirection.
\end{enumerate}
\end{qbox}
\begin{ansbox}
\textbf{Correct Answer: B}

Automated regression flows invoke batch runs, extract metrics (WNS, TNS, Area, Dynamic Power, Runtime, Peak Memory) into structured key-value maps, and dump consolidated \texttt{.csv} / \texttt{.json} tables for signoff tracking.
\end{ansbox}

% ==============================================================================
\section{SpyGlass RTL Linting \& RTL Quality}
% ==============================================================================

\begin{qbox}{Q35: Objective of RTL Linting in Design Flow}{\badgeeasy}
At which stage of the digital ASIC design flow is SpyGlass RTL Linting performed, and what is its primary objective?
\begin{enumerate}[label=\textbf{\Alph*.}]
  \item After Place-and-Route, to detect clock routing skew.
  \item Directly on RTL source code before synthesis, to detect syntax, structural bugs, latches, and simulation-synthesis mismatches.
  \item During fabrication, to measure silicon defect rates.
  \item Inside the FPGA bitstream programmer.
\end{enumerate}
\end{qbox}
\begin{ansbox}
\textbf{Correct Answer: B}

RTL Linting is static source code analysis performed before synthesis. It enforces coding rules, eliminates inferred latches, detects unconnected nets/ports, and identifies simulation-vs-synthesis divergences early in the flow.
\end{ansbox}

\begin{qbox}{Q36: SpyGlass Rule InferLatch (W188)}{\badgemod}
When SpyGlass triggers warning/error \texttt{InferLatch} (or \texttt{W188}) on variable \texttt{next\_state}, what is the root cause?
\begin{enumerate}[label=\textbf{\Alph*.}]
  \item The variable is declared as a \texttt{wire} instead of \texttt{reg}.
  \item The variable is assigned in a combinational block that lacks complete assignments across all conditional branches (missing \texttt{else} or incomplete \texttt{case}).
  \item The flip-flop has an asynchronous reset.
  \item The clock frequency exceeds technology limits.
\end{enumerate}
\end{qbox}
\begin{ansbox}
\textbf{Correct Answer: B}

\texttt{InferLatch} / \texttt{W188} indicates that a combinational always block fails to specify an output value for every possible input branch, forcing the synthesis tool to infer memory storage (a latch).
\end{ansbox}

\begin{qbox}{Q37: SpyGlass Rules W415 vs W416}{\badgemod}
What design violations do SpyGlass rules \texttt{W415} and \texttt{W416} flag?
\begin{enumerate}[label=\textbf{\Alph*.}]
  \item \texttt{W415}: Clock pin unconstrained; \texttt{W416}: Reset pin unconstrained.
  \item \texttt{W415}: Blocking assignment (\texttt{=}) used in a sequential block; \texttt{W416}: Non-blocking assignment (\texttt{<=}) used in a combinational block.
  \item \texttt{W415}: Latch inferred; \texttt{W416}: Combinational loop detected.
  \item \texttt{W415}: Unconnected input; \texttt{W416}: Floating output.
\end{enumerate}
\end{qbox}
\begin{ansbox}
\textbf{Correct Answer: B}

\texttt{W415} warns against blocking assignments in sequential \texttt{always @(posedge clk)} blocks (which cause race conditions). \texttt{W416} warns against non-blocking assignments in combinational \texttt{always @(*)} blocks (which cause simulation event degradation and delta-cycle overhead).
\end{ansbox}

\begin{qbox}{Q38: Incomplete Sensitivity List (W18)}{\badgemod}
Consider \texttt{always @(a or b) begin y = a \& b | c; end}. Why does SpyGlass flag this with rule \texttt{W18}?
\begin{enumerate}[label=\textbf{\Alph*.}]
  \item Signal \texttt{c} is missing from the sensitivity list, causing the RTL simulation output \texttt{y} to remain unchanged when \texttt{c} toggles alone, whereas synthesized gates will evaluate \texttt{c} continuously.
  \item Variables \texttt{a} and \texttt{b} must be separated by commas in Verilog-95.
  \item Bitwise OR \texttt{|} is not synthesizable.
  \item The sensitivity list has too many signals.
\end{enumerate}
\end{qbox}
\begin{ansbox}
\textbf{Correct Answer: A}

An incomplete sensitivity list causes a simulation-synthesis mismatch. The simulator only triggers when \texttt{a} or \texttt{b} changes, ignoring changes in \texttt{c}. Logic synthesis ignores sensitivity lists entirely and builds combinational gates that respond immediately to \texttt{c}.
\end{ansbox}

\begin{qbox}{Q39: Unloaded and Undriven Nets (W528 and W240)}{\badgemod}
Match the SpyGlass rule to its violation:
1. \texttt{W240} \quad 2. \texttt{W528}
\begin{enumerate}[label=\textbf{\Alph*.}]
  \item (1) Signal assigned but never read; (2) Input port declared but unused.
  \item (1) Input port/signal declared but not driven/connected; (2) Variable set/assigned but never read downstream.
  \item (1) Multi-driven net; (2) Combinational loop.
  \item (1) Latch inferred; (2) Setup timing failure.
\end{enumerate}
\end{qbox}
\begin{ansbox}
\textbf{Correct Answer: B}

\texttt{W240} flags undriven/floating inputs (floating nets can cause high gate leakage or unknown \texttt{X} states). \texttt{W528} flags dead code/signals that are driven but never consumed downstream, indicating dead logic or missing connectivity.
\end{ansbox}

\begin{qbox}{Q40: Combinational Feedback Loops (W120 / W121)}{\badgeadv}
Why are combinational feedback loops flagged with high-severity \textbf{Errors} in SpyGlass?
\begin{lstlisting}[style=verilogstyle]
assign a = b & c;
assign b = a ^ d;
\end{lstlisting}
\begin{enumerate}[label=\textbf{\Alph*.}]
  \item They increase power dissipation by a factor of $10\times$.
  \item They create astable ring oscillators or latching states where timing analysis cannot determine topological paths, causing tool hang-ups, glitch oscillations, and unpredictable circuit states.
  \item They cause physical routing DRC violations.
  \item They cannot be converted into standard cell netlists.
\end{enumerate}
\end{qbox}
\begin{ansbox}
\textbf{Correct Answer: B}

Combinational loops lack a synchronizing clock edge. They cause continuous high-frequency oscillations (glitch power) or unpredictable bi-stable memory storage. Furthermore, Static Timing Analysis (STA) tools cannot trace start-to-end paths through closed loops without breaking them artificially.
\end{ansbox}

\begin{qbox}{Q41: Bit-Width Mismatches and Truncation (W164a / W164b)}{\badgeadv}
Consider the following RTL assignment:
\begin{lstlisting}[style=verilogstyle]
wire [3:0] a;
wire [7:0] b = 8'hA5;
assign a = b;
\end{lstlisting}
What does SpyGlass rule \texttt{W164a} report, and what is the exact hardware consequence?
\begin{enumerate}[label=\textbf{\Alph*.}]
  \item Loss of 4 MSBs (\texttt{8'hA} is silently dropped, \texttt{a} gets \texttt{4'h5}), which may cause functional arithmetic truncation bugs if not intended.
  \item The lower 4 bits are dropped.
  \item Synthesis automatically widens \texttt{a} to 8 bits.
  \item Synthesis ties all bits of \texttt{a} to zero.
\end{enumerate}
\end{qbox}
\begin{ansbox}
\textbf{Correct Answer: A}

\texttt{W164a} flags LHS bit-width smaller than RHS (truncation). The upper 4 bits \texttt{b[7:4]} are dropped. If unintentional, this causes datapath corruption or overflow bugs.
\end{ansbox}

\begin{qbox}{Q42: Multi-Driven Bus Contention (W422 / W123)}{\badgeadv}
When multiple tri-state drivers drive a shared internal on-chip bus without an active-enable decoder, what hazard does SpyGlass rule \texttt{W422} prevent?
\begin{enumerate}[label=\textbf{\Alph*.}]
  \item Latch inference on bus lines.
  \item \textbf{Bus Contention} (multiple drivers simultaneously driving \texttt{0} and \texttt{1}, causing excessive short-circuit current and thermal damage) and \textbf{Floating Bus} (no driver active, causing floating gates to drift into high leakage).
  \item Setup timing violations on the clock port.
  \item False reset deassertion.
\end{enumerate}
\end{qbox}
\begin{ansbox}
\textbf{Correct Answer: B}

Tri-state internal buses can suffer from contention (short-circuit path between $V_{DD}$ and $V_{SS}$) or floating states. Modern ASIC guidelines ban internal tri-state buses in favor of explicit multiplexer trees (\texttt{MUX}).
\end{ansbox}

\begin{qbox}{Q43: SpyGlass Lint Waiver Methodology in Signoff}{\badgeadv}
When is it permissible to add a rule waiver to a SpyGlass \texttt{.awl} (Advanced Waiver Language) file during ASIC signoff?
\begin{enumerate}[label=\textbf{\Alph*.}]
  \item Whenever a rule generates more than 100 warnings to clean up the report.
  \item Only after an engineer reviews the specific hierarchy, file, and signal, documents the exact technical rationale (e.g. intentionally unused status bits from a standard IP), and gets formal Lead Signoff approval.
  \item Waivers are never allowed in ASIC design under any circumstances.
  \item Only for timing violations on clock paths.
\end{enumerate}
\end{qbox}
\begin{ansbox}
\textbf{Correct Answer: B}

Waivers must never be used blindly to hide violations. Every waiver must be scoped specifically (by rule name, file, and line/hierarchy) and accompanied by a documented technical justification explaining why the condition is safe.
\end{ansbox}

\begin{qbox}{Q44: STARC Guidelines for Reset Synchronizers}{\badgeadv}
According to STARC (Semiconductor Technology Academic Research Center) design guidelines checked by SpyGlass, how must an asynchronous reset signal be connected to flip-flops?
\begin{enumerate}[label=\textbf{\Alph*.}]
  \item Routed directly from the external chip pad to all registers.
  \item Passed through a Reset Synchronizer (Asynchronously Asserted, Synchronously Deasserted) before driving internal register reset pins.
  \item Connected through a low-pass analog filter.
  \item Converted to a synchronous data input using an XOR gate.
\end{enumerate}
\end{qbox}
\begin{ansbox}
\textbf{Correct Answer: B}

Direct external asynchronous reset deassertion can occur near a clock edge, causing recovery/removal violations and metastability across flip-flops. STARC guidelines require reset synchronizers so reset asserts instantly but releases synchronously with the clock.
\end{ansbox}

\begin{qbox}{Q45: Structuring SpyGlass Sequential Goals}{\badgeadv}
What is the recommended sequential goal execution order in a comprehensive SpyGlass verification flow?
\begin{enumerate}[label=\textbf{\Alph*.}]
  \item \texttt{lint\_rtl} $\rightarrow$ \texttt{cdc\_verify} $\rightarrow$ \texttt{constraints\_verify} $\rightarrow$ \texttt{power\_verify}
  \item \texttt{cdc\_verify} $\rightarrow$ \texttt{lint\_rtl} $\rightarrow$ \texttt{gls}
  \item \texttt{synthesis} $\rightarrow$ \texttt{lint\_rtl} $\rightarrow$ \texttt{pnr}
  \item \texttt{dft\_test} $\rightarrow$ \texttt{lint\_rtl}
\end{enumerate}
\end{qbox}
\begin{ansbox}
\textbf{Correct Answer: A}

Linting must always run first to guarantee clean syntax, types, and structure. Once RTL is lint-clean, CDC analysis checks multi-clock domains, followed by constraint verification and power intent validation.
\end{ansbox}

\begin{qbox}{Q46: Non-Synthesizable Constructs in RTL Code}{\badgeadv}
Which of the following Verilog code segments will be flagged as an un-synthesizable construct by SpyGlass Lint?
\begin{enumerate}[label=\textbf{\Alph*.}]
  \item \texttt{always @(posedge clk or negedge rst\_n)}
  \item \texttt{initial begin clk = 0; forever \#5 clk = \textasciitilde clk; end}
  \item \texttt{assign out = (sel) ? in1 : in0;}
  \item \texttt{always @(*) begin case (state) 2'b00: next = 2'b01; default: next = 2'b00; endcase end}
\end{enumerate}
\end{qbox}
\begin{ansbox}
\textbf{Correct Answer: B}

\texttt{initial} blocks with delays (\texttt{\#5}) and \texttt{forever} loops are simulation-only constructs used in testbenches; physical synthesis engines have no standard cell hardware equivalent for free-running procedural delays.
\end{ansbox}

\begin{qbox}{Q47: Multi-Bit Bus Re-convergence Lint Checks}{\badgeadv}
Why does SpyGlass flag an error when individual bits of a multi-bit data bus (e.g. \texttt{data\_bus[7:0]}) are synchronized across clock domains using independent 2-flip-flop synchronizers?
\begin{enumerate}[label=\textbf{\Alph*.}]
  \item It consumes too many flip-flops.
  \item Due to variable routing delays and cycle uncertainties (0 or 1 cycle latency variation), the synchronized bits can arrive in different destination clock cycles, creating corrupt intermediate data words (Data Incoherency).
  \item It causes clock jitter on the transmitter.
  \item It violates two's complement arithmetic rules.
\end{enumerate}
\end{qbox}
\begin{ansbox}
\textbf{Correct Answer: B}

Multi-bit signals synchronized independently suffer from bit skew. If one bit is captured in cycle $N$ and another in cycle $N+1$, the destination receives invalid garbage values. Multi-bit buses must use Gray coding, Async FIFOs, or MUX-data handshakes.
\end{ansbox}

\begin{qbox}{Q48: Unreachable FSM States Detection}{\badgeadv}
How does SpyGlass detect unreachable FSM states, and why is this critical for silicon area and safety?
\begin{enumerate}[label=\textbf{\Alph*.}]
  \item By checking power dissipation in sleep mode.
  \item By extracting the FSM state transition graph and performing formal reachability analysis from the reset state; unreachable states represent redundant logic that wastes silicon area or indicates a missing transition condition.
  \item By applying random vectors during gate-level simulation.
  \item By measuring clock tree skew.
\end{enumerate}
\end{qbox}
\begin{ansbox}
\textbf{Correct Answer: B}

SpyGlass FSM analysis tools build the complete directed state graph. If any state cannot be reached via any sequence of transitions from the reset state, it flags a violation. Unreachable states indicate logic bugs and waste flip-flops/gates.
\end{ansbox}

% ==============================================================================
\section{Logic Synthesis \& Design Compiler}
% ==============================================================================

\begin{qbox}{Q49: Logic Synthesis Transformation Stages}{\badgeeasy}
Logic synthesis with Synopsys Design Compiler transforms RTL code into a technology-dependent gate-level netlist. What is the correct sequence of internal synthesis steps?
\begin{enumerate}[label=\textbf{\Alph*.}]
  \item Placement $\rightarrow$ Routing $\rightarrow$ CTS
  \item Elaboration / Translation (GTECH) $\rightarrow$ Logic Optimization (Boolean structuring) $\rightarrow$ Technology Mapping (Target \texttt{.lib} cells)
  \item ATPG $\rightarrow$ Scan Insertion $\rightarrow$ DRC
  \item Floorplanning $\rightarrow$ Extraction $\rightarrow$ Signoff
\end{enumerate}
\end{qbox}
\begin{ansbox}
\textbf{Correct Answer: B}

Synthesis first parses and translates RTL into generic technology-independent gates (\texttt{GTECH}), optimizes Boolean equations to minimize logic depth and area, and then maps the optimized equations to specific standard cells from the target semiconductor library (\texttt{.lib}).
\end{ansbox}

\begin{qbox}{Q50: SDC create\_clock Command Syntax}{\badgeeasy}
Which SDC command correctly defines a $100\,\text{MHz}$ clock named \texttt{CORE\_CLK} with a 50\% duty cycle on input port \texttt{clk\_in}?
\begin{enumerate}[label=\textbf{\Alph*.}]
  \item \texttt{create\_clock -name CORE\_CLK -period 10.0 [get\_ports clk\_in]}
  \item \texttt{create\_clock -frequency 100 -duty 50 clk\_in}
  \item \texttt{set\_clock CORE\_CLK 10.0 [get\_ports clk\_in]}
  \item \texttt{create\_clock -period 100.0 -waveform \{0 50\} [get\_pins clk\_in]}
\end{enumerate}
\end{qbox}
\begin{ansbox}
\textbf{Correct Answer: A}

The time unit in standard SDC is nanoseconds ($100\,\text{MHz} \implies T = 10.0\,\text{ns}$). The default waveform is $\{0\ \text{Period}/2\}$ ($50\%$ duty cycle), so \texttt{-period 10.0} defines the $100\,\text{MHz}$ clock.
\end{ansbox}

\begin{qbox}{Q51: Wire Load Models (WLM) vs Topographical Mode}{\badgemod}
Why is standard synthesis using traditional Wire Load Models (WLM) inaccurate at deep sub-micron (DSM) technology nodes below $65\,\text{nm}$?
\begin{enumerate}[label=\textbf{\Alph*.}]
  \item WLMs overestimate gate leakage power.
  \item Interconnect resistance and capacitance dominate total path delay, and statistical WLMs fail to estimate actual physical wire lengths accurately, causing massive correlation discrepancies with post-PnR timing.
  \item WLMs do not support sequential cells.
  \item WLMs only work with FPGA architectures.
\end{enumerate}
\end{qbox}
\begin{ansbox}
\textbf{Correct Answer: B}

In DSM nodes ($45\,\text{nm}, 28\,\text{nm}, 7\,\text{nm}$), wire delay contributes $>60\%$ of total path delay. Statistical WLMs guess wire lengths based on fanout alone. Synopsys DC-Topographical (\texttt{compile\_ultra -spg}) uses physical floorplan and placement engines during synthesis to accurately predict real wire parasitics.
\end{ansbox}

\begin{qbox}{Q52: set\_driving\_cell vs set\_drive}{\badgemod}
Why is \texttt{set\_driving\_cell -lib\_cell INV\_X4 [get\_ports din]} preferred over \texttt{set\_drive 0 [get\_ports din]} in realistic SDC synthesis scripts?
\begin{enumerate}[label=\textbf{\Alph*.}]
  \item \texttt{set\_drive 0} sets the driving resistance to infinity.
  \item \texttt{set\_driving\_cell} accurately models the non-linear voltage-dependent drive resistance, output transition time (slew), and input capacitance from the exact standard cell library model.
  \item \texttt{set\_driving\_cell} automatically inserts a physical buffer on the PCB.
  \item \texttt{set\_drive} is not supported in Design Compiler.
\end{enumerate}
\end{qbox}
\begin{ansbox}
\textbf{Correct Answer: B}

Real input ports are driven by external chip pads or standard cells. \texttt{set\_driving\_cell} extracts the exact Non-Linear Delay Model (NLDM) or CCS timing tables of the driving cell to compute accurate input transition slews into the chip's first-stage logic.
\end{ansbox}

\begin{qbox}{Q53: Virtual Clocks in SDC Timing Constraints}{\badgemod}
What is a \textbf{Virtual Clock}, and when is it necessary in SDC constraints?
\begin{enumerate}[label=\textbf{\Alph*.}]
  \item A clock that has no physical source pin/port in the design, used solely as a timing reference for constraining input arrival times and output delay requirements.
  \item A simulated clock generated inside a testbench.
  \item A clock operating at zero frequency.
  \item A clock used only for scan chain shifting.
\end{enumerate}
\end{qbox}
\begin{ansbox}
\textbf{Correct Answer: A}

A virtual clock is defined using \texttt{create\_clock -name VCLK -period 10.0} without specifying any source pin/port. It models the external transmitting or receiving device's clock for accurate \texttt{set\_input\_delay} and \texttt{set\_output\_delay} constraints.
\end{ansbox}

\begin{qbox}{Q54: set\_clock\_uncertainty Components in Pre-Layout}{\badgemod}
During pre-layout logic synthesis, what physical effects does \texttt{set\_clock\_uncertainty} model?
\begin{enumerate}[label=\textbf{\Alph*.}]
  \item Power grid voltage drop only.
  \item Clock jitter (PLL phase noise) + estimated pre-CTS clock skew + clock tree synthesis margins.
  \item Leakage current variations.
  \item Thermal dissipation across package pins.
\end{enumerate}
\end{qbox}
\begin{ansbox}
\textbf{Correct Answer: B}

Before physical Clock Tree Synthesis (CTS), clock trees do not exist (clocks are ideal). \texttt{set\_clock\_uncertainty} acts as a timing margin representing both the physical clock jitter (source instability) and the estimated clock skew that will appear post-CTS.
\end{ansbox}

\begin{qbox}{Q55: Synthesis Optimization Goals Priority}{\badgemod}
In what order does Synopsys Design Compiler prioritize optimization constraints during \texttt{compile\_ultra}?
\begin{enumerate}[label=\textbf{\Alph*.}]
  \item Area $\rightarrow$ Power $\rightarrow$ Dynamic Timing $\rightarrow$ DRCs
  \item Design Rule Constraints (DRCs: Max Capacitance, Max Transition, Max Fanout) $\rightarrow$ Setup Timing (Delay) $\rightarrow$ Area $\rightarrow$ Dynamic/Leakage Power
  \item Dynamic Power $\rightarrow$ Area $\rightarrow$ Setup Timing $\rightarrow$ DRCs
  \item Minimum Delay (Hold) $\rightarrow$ Maximum Delay (Setup) $\rightarrow$ Area
\end{enumerate}
\end{qbox}
\begin{ansbox}
\textbf{Correct Answer: B}

Design Compiler strictly prioritizes electrical integrity: DRCs (Max Cap, Max Transition) must be satisfied first. Next is timing closure (WNS/TNS). Only after timing is met will the tool optimize for silicon area and power dissipation.
\end{ansbox}

\begin{qbox}{Q56: Operating Conditions \& Slow PVT Corner Definition}{\badgemod}
What does a "Worst-Case" (Slow) PVT corner represent for setup timing analysis in standard CMOS?
\begin{enumerate}[label=\textbf{\Alph*.}]
  \item High Voltage, Low Temperature, Fast Process
  \item Low Voltage (e.g. $0.9 \times V_{nom}$), High Temperature (e.g. $+125^\circ\text{C}$ in standard CMOS), Slow Process Corner (SS)
  \item Typical Voltage, $25^\circ\text{C}$, Typical Process (TT)
  \item Zero Voltage, $0^\circ\text{C}$, Fast Process (FF)
\end{enumerate}
\end{qbox}
\begin{ansbox}
\textbf{Correct Answer: B}

Transistor switching speed degrades at low supply voltage ($V_{DD}$), slow process silicon (SS), and high junction temperature ($125^\circ\text{C}$ in traditional CMOS due to reduced carrier mobility), representing the maximum path delay for setup analysis.
\end{ansbox}

\begin{qbox}{Q57: Generated Clock Inversion and Division}{\badgeadv}
A register \texttt{u\_div2/q\_reg} divides a $10\,\text{ns}$ master clock \texttt{CLK} by 2 with an inverted output. Which SDC command correctly defines this generated clock?
\begin{enumerate}[label=\textbf{\Alph*.}]
  \item \texttt{create\_clock -period 20.0 [get\_pins u\_div2/q\_reg/Q]}
  \item \texttt{create\_generated\_clock -name GEN\_CLK -source [get\_ports clk] -divide\_by 2 -invert [get\_pins u\_div2/q\_reg/Q]}
  \item \texttt{set\_clock\_latency -source 2.0 [get\_pins u\_div2/q\_reg/Q]}
  \item \texttt{create\_generated\_clock -period 20.0 [get\_pins u\_div2/q\_reg/Q]}
\end{enumerate}
\end{qbox}
\begin{ansbox}
\textbf{Correct Answer: B}

\texttt{create\_generated\_clock} maintains true phase and latency relationships with the master source clock. Specifying \texttt{-source [get\_ports clk]}, \texttt{-divide\_by 2}, and \texttt{-invert} accurately preserves clock tree path tracing and duty-cycle/polarity relationships in STA.
\end{ansbox}

\begin{qbox}{Q58: Clock Gating Insertion Transformation}{\badgeadv}
When \texttt{compile\_ultra -gate\_clock} is enabled in Design Compiler, what transformation is performed on wide registers?
\begin{enumerate}[label=\textbf{\Alph*.}]
  \item Multiplexer feedback loops on register D-inputs are converted into latch-based Integrated Clock Gating (ICG) cells on the clock pin, saving dynamic power.
  \item Registers are replaced with dynamic domino logic.
  \item Clock buffers are removed to save area.
  \item Registers are converted into dual-edge triggered flip-flops.
\end{enumerate}
\end{qbox}
\begin{ansbox}
\textbf{Correct Answer: A}

Standard RTL synchronous enable logic \texttt{if (en) q <= d;} synthesizes to a feedback MUX that keeps the clock toggling every cycle. Clock gating replaces the MUXes with an ICG cell that turns off the clock to the register bank whenever \texttt{en=0}, eliminating clock tree switching power.
\end{ansbox}

\begin{qbox}{Q59: Sequential Retiming Purpose}{\badgeadv}
What is the effect of sequential \textbf{Retiming} during synthesis?
\begin{enumerate}[label=\textbf{\Alph*.}]
  \item It removes all flip-flops from the netlist.
  \item It automatically pushes/pulls flip-flops across combinational logic gates without altering the design's input/output transfer latency, balancing slack across stages.
  \item It adjusts the clock frequency during runtime.
  \item It replaces asynchronous resets with synchronous resets.
\end{enumerate}
\end{qbox}
\begin{ansbox}
\textbf{Correct Answer: B}

Retiming relocates registers across logic gates to balance combinational delays between pipeline stages. This reduces the worst-case critical path delay ($T_{clk}$) and increases maximum operating frequency without changing the overall pipeline latency.
\end{ansbox}

\begin{qbox}{Q60: set\_max\_fanout vs set\_max\_capacitance DRCs}{\badgeadv}
Why is \texttt{set\_max\_capacitance} a true physical DRC constraint while \texttt{set\_max\_fanout} is considered a design rule guideline?
\begin{enumerate}[label=\textbf{\Alph*.}]
  \item Capacitance limits are determined by transistor physical limits from the standard cell library characterization (\texttt{.lib}), directly affecting slew and reliability; fanout is a numerical gate count that does not account for variable wire lengths.
  \item Fanout is only used in FPGA synthesis.
  \item \texttt{set\_max\_capacitance} applies only to clock pins.
  \item Design Compiler ignores \texttt{set\_max\_capacitance}.
\end{enumerate}
\end{qbox}
\begin{ansbox}
\textbf{Correct Answer: A}

Pin transition time and signal integrity depend on the physical capacitive load (farads). Fanout simply counts the number of load pins, ignoring wire capacitance and pin input capacitance differences. Standard cell \texttt{.lib} rules mandate \texttt{max\_capacitance} limits to prevent cell damage and slew degradation.
\end{ansbox}

\begin{qbox}{Q61: False Paths vs Multicycle Paths Decisions}{\badgeadv}
When should a path be constrained as a \textbf{Multicycle Path} rather than a \textbf{False Path}?
\begin{enumerate}[label=\textbf{\Alph*.}]
  \item When the path is physically impossible to trigger under any condition.
  \item When data is transferred across asynchronous clock domains without a handshake.
  \item When data requires multiple clock cycles (e.g. 2 or 3 cycles) to propagate through complex arithmetic logic, but is deliberately sampled only on cycle $N$ via an enable signal.
  \item When hold time violation cannot be fixed.
\end{enumerate}
\end{qbox}
\begin{ansbox}
\textbf{Correct Answer: C}

A multicycle path is a functionally valid, synchronous path where the destination flip-flop is designed to capture data after $N > 1$ clock cycles. A false path tells the STA tool to completely ignore the path for all timing checks. Constraining a multicycle path as a false path prevents hold and setup verification.
\end{ansbox}

\begin{qbox}{Q62: Enforcing No-Latch Synthesis Rules}{\badgeadv}
How does Design Compiler optimize an RTL block when the user sets \texttt{hdlin\_check\_no\_latch = true}?
\begin{enumerate}[label=\textbf{\Alph*.}]
  \item It replaces all latches with 2 flip-flops.
  \item It issues an immediate fatal compilation error if any latch is inferred during HDL translation, halting synthesis before mapping.
  \item It ties latch enable pins to $V_{DD}$.
  \item It converts latches into tri-state buffers.
\end{enumerate}
\end{qbox}
\begin{ansbox}
\textbf{Correct Answer: B}

\texttt{hdlin\_check\_no\_latch} is a strict synthesis directive that causes the parser to halt with a fatal error if any combinational block infers a latch, enforcing clean RTL standards before technology mapping.
\end{ansbox}

\begin{qbox}{Q63: Critical Path Area Recovery Mechanism}{\badgeadv}
During the final phase of \texttt{compile\_ultra}, how does Design Compiler perform \textbf{Area Recovery} on non-critical paths?
\begin{enumerate}[label=\textbf{\Alph*.}]
  \item By deleting redundant flip-flops.
  \item By down-sizing oversized high-drive standard cells (e.g. replacing \texttt{INV\_X16} with \texttt{INV\_X2} or \texttt{INV\_X1}) on paths with positive slack, reducing area and dynamic power without violating setup timing.
  \item By reducing the supply voltage.
  \item By disabling scan chains.
\end{enumerate}
\end{qbox}
\begin{ansbox}
\textbf{Correct Answer: B}

To close timing on the critical path, synthesis uses large, high-drive, low-$V_t$ cells. On non-critical paths with ample positive timing slack, the area recovery engine down-sizes cells to smaller sizes (and higher $V_t$) to minimize silicon area and leakage.
\end{ansbox}

\begin{qbox}{Q64: Boundary Optimization and Port Freezing}{\badgeadv}
What does the \texttt{set\_boundary\_optimization false [get\_designs u\_block]} command prevent Design Compiler from doing?
\begin{enumerate}[label=\textbf{\Alph*.}]
  \item It prevents the tool from synthesizing the internal logic of \texttt{u\_block}.
  \item It preserves module port boundaries, preventing constant propagation across ports, hierarchical pin inversion, and merging of unused input/output ports.
  \item It disables place and route boundaries.
  \item It forces all ports to be registered.
\end{enumerate}
\end{qbox}
\begin{ansbox}
\textbf{Correct Answer: B}

Boundary optimization allows the synthesis tool to optimize across hierarchical sub-module boundaries (e.g., inverting complementary ports, propagating fixed constants into sub-blocks). Disabling it freezes the module interface, which is essential for Formal Verification (LEC) and hierarchical re-use.
\end{ansbox}

\begin{qbox}{Q65: Multi-Corner Multi-Mode (MCMM) Synthesis Architecture}{\badgeadv}
Why is Multi-Corner Multi-Mode (MCMM) optimization essential in modern digital synthesis?
\begin{enumerate}[label=\textbf{\Alph*.}]
  \item To design chips that operate without clock signals.
  \item To simultaneously optimize the design across multiple operational functional modes (e.g. High-Performance Mode, Low-Power Sleep Mode, Scan Test Mode) and multiple PVT corners (Slow-Corner setup, Fast-Corner hold) in a single run.
  \item To combine FPGA and ASIC libraries together.
  \item To eliminate thermal dissipation.
\end{enumerate}
\end{qbox}
\begin{ansbox}
\textbf{Correct Answer: B}

Designs have different operating modes (Functional, Scan Shift, Scan Capture, BIST) with conflicting timing paths, operating at different voltages and temperatures. MCMM synthesizes and closes timing across all modes and PVT corners concurrently, preventing fix-induced regressions.
\end{ansbox}

\begin{qbox}{Q66: set\_dont\_touch Constraints on Debug Cells}{\badgeadv}
Why does an engineer apply \texttt{set\_dont\_touch [get\_cells u\_debug\_reg]} during synthesis?
\begin{enumerate}[label=\textbf{\Alph*.}]
  \item To instruct Design Compiler not to optimize, remove, gate, or resize the specified cell or net, preserving it for silicon probe points or ECOs.
  \item To prevent the cell from being placed in the floorplan.
  \item To disable timing checks on that cell.
  \item To make the cell immune to radiation.
\end{enumerate}
\end{qbox}
\begin{ansbox}
\textbf{Correct Answer: A}

Synthesis tools aggressively prune unused registers or constant logic. Setting \texttt{set\_dont\_touch} preserves critical debug registers, test structures, or precision analog-interface logic from being optimized away or modified.
\end{ansbox}

\begin{qbox}{Q67: Total Negative Slack (TNS) vs Worst Negative Slack (WNS)}{\badgeadv}
In a synthesis timing summary report: $\text{WNS} = -0.15\,\text{ns}$ and $\text{TNS} = -124.5\,\text{ns}$. What do these metrics indicate to the ASIC designer?
\begin{enumerate}[label=\textbf{\Alph*.}]
  \item The design has only 1 violating path.
  \item The worst failing path violates setup timing by $0.15\,\text{ns}$, and there are hundreds of paths violating timing across the chip whose accumulated negative slack sums to $-124.5\,\text{ns}$.
  \item The clock frequency must be reduced by $124.5\,\text{ns}$.
  \item The design passes all timing checks.
\end{enumerate}
\end{qbox}
\begin{ansbox}
\textbf{Correct Answer: B}

WNS (Worst Negative Slack) is the maximum single timing violation in the design. TNS (Total Negative Slack) is the sum of negative slacks across all violating endpoints. A small WNS with large TNS indicates widespread timing closure issues requiring architectural or global constraint adjustments.
\end{ansbox}

\begin{qbox}{Q68: High-Fanout Net Synthesis and Ideal Networks}{\badgeadv}
Why are clock and reset nets marked as \texttt{set\_ideal\_network} during logic synthesis in Design Compiler?
\begin{enumerate}[label=\textbf{\Alph*.}]
  \item To make synthesis finish faster by ignoring power consumption.
  \item Because high-fanout trees (clocks, resets) are physically synthesized and balanced later during Clock Tree Synthesis (CTS) in Physical Design; inserting unconstrained buffer trees during synthesis would distort real placement and waste area.
  \item Because ideal nets have infinite voltage.
  \item Because resets do not require routing wires.
\end{enumerate}
\end{qbox}
\begin{ansbox}
\textbf{Correct Answer: B}

During synthesis, physical placement locations of millions of standard cells are not finalized. If synthesis attempted to build buffer trees for clocks or global resets, it would create suboptimal, un-routable trees. They are treated as ideal (zero delay, infinite drive) and built accurately during CTS.
\end{ansbox}

% ==============================================================================
\section{Static Timing Analysis (STA) \& PrimeTime}
% ==============================================================================

\begin{qbox}{Q69: Fundamental Setup Timing Equation}{\badgeeasy}
For a register-to-register path with clock period $T_{clk}$, launch clock delay $T_{launch}$, capture clock delay $T_{capture}$, register clock-to-Q delay $T_{cq}$, combinational logic delay $T_{comb}$, and setup time $T_{setup}$, what is the condition for \textbf{Setup Timing} to be met?
\begin{enumerate}[label=\textbf{\Alph*.}]
  \item $T_{cq} + T_{comb} + T_{setup} \le T_{clk} + (T_{capture} - T_{launch})$
  \item $T_{cq} + T_{comb} \ge T_{hold}$
  \item $T_{clk} \le T_{setup} + T_{hold}$
  \item $T_{comb} = T_{clk}$
\end{enumerate}
\end{qbox}
\begin{ansbox}
\textbf{Correct Answer: A}

Data launched at cycle 0 must arrive and settle at the capture register before the next clock edge at cycle 1 ($T_{clk}$). The arrival time is $T_{launch} + T_{cq} + T_{comb}$. The required time is $T_{capture} + T_{clk} - T_{setup}$. Rearranging gives $T_{cq} + T_{comb} + T_{setup} \le T_{clk} + T_{skew}$, where $T_{skew} = T_{capture} - T_{launch}$.
\end{ansbox}

\begin{qbox}{Q70: Fundamental Hold Timing Equation}{\badgeeasy}
What is the condition for \textbf{Hold Timing} to be satisfied on a synchronous register-to-register path?
\begin{enumerate}[label=\textbf{\Alph*.}]
  \item $T_{launch} + T_{cq} + T_{comb(min)} \ge T_{capture} + T_{hold}$
  \item $T_{comb(max)} \le T_{clk} - T_{setup}$
  \item $T_{clk} \ge T_{hold}$
  \item $T_{skew} = T_{clk}$
\end{enumerate}
\end{qbox}
\begin{ansbox}
\textbf{Correct Answer: A}

Data launched at cycle 0 must not propagate so fast that it corrupts data being captured at cycle 0 at the capture flip-flop. Thus, minimum arrival time $T_{launch} + T_{cq(min)} + T_{comb(min)}$ must be greater than or equal to $T_{capture} + T_{hold}$. Notice that clock period $T_{clk}$ does not appear in the hold equation.
\end{ansbox}

\begin{qbox}{Q71: Clock Skew Impact on Setup and Hold}{\badgemod}
\textbf{Positive Clock Skew} ($T_{capture} > T_{launch}$, where the capture clock edge arrives later than the launch clock edge):
\begin{enumerate}[label=\textbf{\Alph*.}]
  \item Helps Setup time but hurts Hold time.
  \item Hurts Setup time but helps Hold time.
  \item Hurts both Setup and Hold time equally.
  \item Has zero impact on timing closure.
\end{enumerate}
\end{qbox}
\begin{ansbox}
\textbf{Correct Answer: A}

Positive skew provides extra time for data to propagate through combinational logic, thereby \textbf{increasing setup slack}. However, it increases the capture clock arrival time for the same clock edge, requiring data to remain stable longer, which \textbf{reduces hold slack} and can cause hold violations.
\end{ansbox}

\begin{qbox}{Q72: Setup vs Hold PVT Operating Corners}{\badgemod}
Which operating corners are used for Setup and Hold timing signoff?
\begin{enumerate}[label=\textbf{\Alph*.}]
  \item Setup: Fast-Process, High-Voltage, Low-Temp; Hold: Slow-Process, Low-Voltage, High-Temp
  \item Setup: Slow-Process, Low-Voltage, High-Temp (SS/LowV/HighT - Max Delay); Hold: Fast-Process, High-Voltage, Low-Temp (FF/HighV/LowT - Min Delay)
  \item Setup and Hold are always checked at Typical (TT) conditions.
  \item Setup is checked at $0\,\text{K}$; Hold is checked at $300\,\text{K}$.
\end{enumerate}
\end{qbox}
\begin{ansbox}
\textbf{Correct Answer: B}

Setup checks worst-case maximum datapath delay against the shortest effective period (slow silicon, low voltage, high temp). Hold checks best-case minimum datapath delay against fast clock transitions (fast silicon, high voltage, low temp), where fast data easily overwrites the previous state.
\end{ansbox}

\begin{qbox}{Q73: Clock Jitter vs Clock Skew Distinction}{\badgemod}
What is the fundamental difference between \textbf{Clock Skew} and \textbf{Clock Jitter} in STA?
\begin{enumerate}[label=\textbf{\Alph*.}]
  \item Skew is dynamic cycle-to-cycle variation; Jitter is static spatial delay difference.
  \item \textbf{Clock Skew} is the spatial difference in arrival times between two different physical clock pins; \textbf{Clock Jitter} is the temporal (cycle-to-cycle or phase) variation in the clock edge arrival time at a single clock pin over time.
  \item Skew can be eliminated with PLLs; Jitter cannot.
  \item Skew affects only hold time; Jitter affects only setup time.
\end{enumerate}
\end{qbox}
\begin{ansbox}
\textbf{Correct Answer: B}

Skew is spatial (difference in clock distribution routing latencies across the chip). Jitter is temporal (phase noise and cycle-to-cycle variation from the PLL and oscillator). Both reduce setup timing margins, but only skew directly impacts hold checks.
\end{ansbox}

\begin{qbox}{Q74: Recovery and Removal Asynchronous Checks}{\badgemod}
What are \textbf{Recovery} and \textbf{Removal} timing checks in PrimeTime?
\begin{enumerate}[label=\textbf{\Alph*.}]
  \item Checks for memory read and write cycles.
  \item \textbf{Recovery} is the minimum time an asynchronous control signal (reset/preset) must be deasserted \textit{before} the active clock edge (equivalent to Setup); \textbf{Removal} is the minimum time it must remain asserted \textit{after} the active clock edge (equivalent to Hold).
  \item Checks for power-down sequencing.
  \item Checks for scan chain shift frequency.
\end{enumerate}
\end{qbox}
\begin{ansbox}
\textbf{Correct Answer: B}

When releasing an asynchronous reset, if deassertion occurs too close to the active clock edge, flip-flops enter a metastable state or some flip-flops exit reset while others remain in reset. Recovery prevents setup-like violations, and Removal prevents hold-like violations on the reset pin.
\end{ansbox}

\begin{qbox}{Q75: Data-to-Data Timing Checks}{\badgemod}
When is a \texttt{set\_data\_check} constraint applied in PrimeTime?
\begin{enumerate}[label=\textbf{\Alph*.}]
  \item Between two flip-flop outputs on a scan chain.
  \item To define setup and hold requirements between two asynchronous or non-clocked data signals (e.g. write-enable setup to address lines in an asynchronous SRAM interface).
  \item To check data bus bit-width matching.
  \item To verify parity bits in UART.
\end{enumerate}
\end{qbox}
\begin{ansbox}
\textbf{Correct Answer: B}

Non-clocked interfaces (like asynchronous memory or custom handshakes) require setup and hold stability between two data signals (e.g. \texttt{data} valid before \texttt{write\_enable} strobe). \texttt{set\_data\_check} instructs STA to enforce timing checks between two pure data pins.
\end{ansbox}

\begin{qbox}{Q76: Integrated Clock Gating Setup and Hold Checks}{\badgemod}
For an active-high clock gating cell (AND gate driven by \texttt{CLK} and \texttt{EN}):
\begin{enumerate}[label=\textbf{\Alph*.}]
  \item \texttt{EN} must change only while \texttt{CLK} is high.
  \item \texttt{EN} must satisfy setup and hold timing relative to the \textbf{falling edge} of \texttt{CLK} before entering the latch, preventing glitches while \texttt{CLK} is high.
  \item No timing check is required for clock gating cells.
  \item \texttt{EN} must be tied to $V_{DD}$.
\end{enumerate}
\end{qbox}
\begin{ansbox}
\textbf{Correct Answer: B}

In an active-high integrated clock gating cell (ICG), the latch is transparent when \texttt{CLK=0} and captures \texttt{EN} on the falling edge of \texttt{CLK}. STA verifies that \texttt{EN} arrives before the falling clock edge (Clock Gating Setup) and does not change prematurely after it (Clock Gating Hold).
\end{ansbox}

\begin{qbox}{Q77: Quasi-Static False Path Declaration}{\badgemod}
Which of the following paths should be declared as a \textbf{False Path} (\texttt{set\_false\_path})?
\begin{enumerate}[label=\textbf{\Alph*.}]
  \item The critical datapath in a 64-bit floating point adder.
  \item Static configuration register bits written once during boot-up and never toggled during normal execution.
  \item The scan enable distribution tree.
  \item Memory read data bus during functional mode.
\end{enumerate}
\end{qbox}
\begin{ansbox}
\textbf{Correct Answer: B}

Quasi-static configuration signals (set during chip power-up and constant during normal operation) do not switch dynamically cycle-to-cycle. Constraining them as false paths prevents the STA tool from wasting runtime and inserting unnecessary buffers on paths that never transition in operation.
\end{ansbox}

\begin{qbox}{Q78: SDC Multicycle Path Setup and Hold Modifiers}{\badgemod}
A designer specifies a 2-cycle multicycle path:
\begin{lstlisting}[style=tclstyle]
set_multicycle_path 2 -setup -from [get_cells FF1] -to [get_cells FF2]
\end{lstlisting}
Without an explicit \texttt{-hold} command, what hold check does PrimeTime perform by default?
\begin{enumerate}[label=\textbf{\Alph*.}]
  \item Hold check at cycle 0 (the launch edge).
  \item Hold check at cycle 1 (one cycle before the new setup capture edge at cycle 2), which is overly restrictive and incorrect for a single-source clock.
  \item Hold checks are completely disabled.
  \item Hold check at cycle 2.
\end{enumerate}
\end{qbox}
\begin{ansbox}
\textbf{Correct Answer: B}

By default in SDC, setting \texttt{-setup N} shifts the hold check to $N-1$ cycles after the launch edge. For an actual multi-cycle datapath running on the same clock, the data is launched at cycle 0 and must not overwrite cycle 0 data, requiring: \texttt{set\_multicycle\_path 1 -hold -from [get\_cells FF1] -to [get\_cells FF2]}.
\end{ansbox}

\begin{qbox}{Q79: On-Chip Variation (OCV) Setup Slack Calculation}{\badgeadv}
In an OCV setup analysis with $T_{clk} = 10\,\text{ns}$, $T_{launch\_clk} = 2.0\,\text{ns}$, $T_{capture\_clk} = 2.0\,\text{ns}$, $T_{data} = 8.5\,\text{ns}$, $T_{setup} = 0.5\,\text{ns}$.  
If a \textbf{10\% late derate} ($1.10$) is applied to the launch path and a \textbf{10\% early derate} ($0.90$) is applied to the capture path:
What is the resulting setup slack?
\begin{enumerate}[label=\textbf{\Alph*.}]
  \item $+0.5\,\text{ns}$
  \item $-1.25\,\text{ns}$
  \item $-0.4\,\text{ns}$
  \item $-0.9\,\text{ns}$
\end{enumerate}
\end{qbox}
\begin{ansbox}
\textbf{Correct Answer: C}

\begin{itemize}[leftmargin=1.5em]
  \item Derated Launch Arrival Time = $(2.0 \times 1.10) + (8.5 \times 1.10) = 2.2 + 9.35 = 11.55\,\text{ns}$.
  \item Derated Capture Required Time = $10.0 + (2.0 \times 0.90) - 0.5 = 10.0 + 1.8 - 0.5 = 11.15\,\text{ns}$.
  \item Setup Slack = $\text{Required} - \text{Arrival} = 11.15 - 11.55 = -0.40\,\text{ns}$ (Violated).
\end{itemize}
\end{ansbox}

\begin{qbox}{Q80: Common Path Pessimism Removal (CPPR / CRPR)}{\badgeadv}
Why is Common Path Pessimism Removal (CPPR) essential during OCV Static Timing Analysis?
\begin{enumerate}[label=\textbf{\Alph*.}]
  \item To increase the clock frequency of the PLL.
  \item To remove artificial timing pessimism caused by applying conflicting derating factors (e.g. early derate and late derate simultaneously) to the identical, shared physical clock tree buffer segments feeding both launch and capture flip-flops.
  \item To compensate for board-level trace mismatches.
  \item To eliminate crosstalk glitches.
\end{enumerate}
\end{qbox}
\begin{ansbox}
\textbf{Correct Answer: B}

A single physical clock buffer in the common clock tree cannot be physically fast (early) and slow (late) at the exact same instant. CPPR calculates the delay difference between early and late derates on the shared clock path and adds it back as a credit to the slack calculation.
\end{ansbox}

\begin{qbox}{Q81: Advanced OCV (AOCV) and Parametric OCV (POCV)}{\badgeadv}
Why are AOCV and POCV significantly superior to flat global OCV derating at advanced FinFET nodes ($16\,\text{nm}, 7\,\text{nm}, 3\,\text{nm}$)?
\begin{enumerate}[label=\textbf{\Alph*.}]
  \item Flat OCV derates assume worst-case variations across all cells globally, causing severe over-design. AOCV scales derates based on physical path depth and bounding-box distance (statistical averaging), while POCV models individual gate delay variations as statistical Gaussian distributions ($\mu \pm 3\sigma$).
  \item Flat OCV does not support multi-voltage domains.
  \item AOCV and POCV eliminate the need for standard cell characterization.
  \item AOCV runs without a timing engine.
\end{enumerate}
\end{qbox}
\begin{ansbox}
\textbf{Correct Answer: A}

As logic depth increases, random independent process variations average out (Law of Large Numbers). Flat OCV applies a pessimistic uniform penalty (e.g. $15\%$). AOCV calculates lower derates for deeper paths and shorter distances, and POCV uses statistical cell-level sigma ($\sigma$) variations, eliminating excessive design margins.
\end{ansbox}

\begin{qbox}{Q82: Crosstalk Delay and Miller Coupling Capacitance}{\badgeadv}
When two parallel on-chip metal tracks (Aggressor and Victim) switch simultaneously:
What happens to the victim line delay when both switch in \textbf{opposite directions} versus \textbf{same direction}?
\begin{enumerate}[label=\textbf{\Alph*.}]
  \item Opposite: Delay decreases; Same: Delay increases.
  \item \textbf{Opposite directions}: Effective coupling capacitance doubles ($2\times C_c$ via Miller effect), significantly \textbf{increasing victim transition delay}; \textbf{Same direction}: Effective coupling capacitance drops towards zero, \textbf{accelerating signal propagation} (speed-up / hold hazard).
  \item No change in delay occurs in either case.
  \item Both signals enter high-impedance state.
\end{enumerate}
\end{qbox}
\begin{ansbox}
\textbf{Correct Answer: B}

Crosstalk Miller effect: Opposite-direction switching causes a $\Delta V = 2 V_{DD}$ transition across the coupling capacitor $C_c$, doubling dynamic charge transfer and slowing down the victim (causing setup violations). Same-direction switching eliminates the voltage delta, speeding up the victim (causing hold violations).
\end{ansbox}

\begin{qbox}{Q83: Resolving Hold Violations vs Setup Violations}{\badgeadv}
Why must hold timing violations be fixed exclusively by adding delay buffers on the datapath rather than increasing the clock period?
\begin{enumerate}[label=\textbf{\Alph*.}]
  \item Hold violations depend on the speed of light.
  \item The hold timing inequality ($T_{launch} + T_{cq} + T_{comb} \ge T_{capture} + T_{hold}$) is completely independent of the clock period $T_{clk}$; slowing down the clock does not fix hold violations, and uncorrected hold violations result in un-testable, permanently broken silicon.
  \item Setup violations can only be fixed with inverters.
  \item Hold violations only occur during fabrication testing.
\end{enumerate}
\end{qbox}
\begin{ansbox}
\textbf{Correct Answer: B}

Setup violations can always be resolved after tape-out by running the chip at a lower clock frequency (increasing $T_{clk}$). Hold violations are frequency-independent: if data races ahead of the same clock edge, the circuit fails at any clock frequency ($1\,\text{kHz}$ or $1\,\text{GHz}$). Hold must be closed clean in silicon.
\end{ansbox}

\begin{qbox}{Q84: Transparent Latches and Time Borrowing}{\badgeadv}
How does a two-phase non-overlapping clock scheme with transparent latches achieve \textbf{Time Borrowing}?
\begin{enumerate}[label=\textbf{\Alph*.}]
  \item By accelerating the PLL frequency.
  \item If combinational logic preceding Latch 1 completes late (after the opening clock edge), the data can still propagate through the transparent latch during its open window, borrowing unused time from the subsequent clock phase without causing a timing failure.
  \item By borrowing power from unused blocks.
  \item By converting negative slack to positive slack using software.
\end{enumerate}
\end{qbox}
\begin{ansbox}
\textbf{Correct Answer: B}

Unlike edge-triggered flip-flops that sample data strictly at a single instant, level-sensitive latches remain open and transparent for the entire active clock phase. If the input data arrives during the transparent phase, it passes through immediately, utilizing remaining slack in that phase.
\end{ansbox}

\begin{qbox}{Q85: Source-Synchronous Interface Advantages}{\badgeadv}
Why are high-speed DDR memory interfaces (e.g. LPDDR5, DDR4) designed as \textbf{Source-Synchronous} interfaces?
\begin{enumerate}[label=\textbf{\Alph*.}]
  \item They use a single global oscillator located on the motherboard.
  \item The transmitting chip sends both the Data bus and a matched Clock/Strobe signal along identical PCB trace lengths, eliminating board-level clock skew and allowing operating speeds beyond several gigabits per second.
  \item They eliminate the need for read/write controllers.
  \item They operate with zero power.
\end{enumerate}
\end{qbox}
\begin{ansbox}
\textbf{Correct Answer: B}

In common-clock (system-synchronous) interfaces, board-level clock skew between transmitter and receiver chips limits maximum frequency to $\sim 200\,\text{MHz}$. In source-synchronous interfaces, the strobe ($DQS$) travels with the data ($DQ$), tracking matched delays and PVT drifts.
\end{ansbox}

\begin{qbox}{Q86: Path Grouping and Optimization Effort}{\badgeadv}
Why does an STA engineer group timing paths into distinct path groups (\texttt{group\_path -name CLK\_GRP -critical\_range 1.0})?
\begin{enumerate}[label=\textbf{\Alph*.}]
  \item To randomize the order of timing checks.
  \item To prevent one dominant failing path group (e.g., a slow I/O interface) from starving other internal core register-to-register paths of optimization effort during synthesis and physical design.
  \item To disable timing checks on non-critical paths.
  \item To group registers by physical voltage domain.
\end{enumerate}
\end{qbox}
\begin{ansbox}
\textbf{Correct Answer: B}

Optimization engines focus primarily on the worst overall negative slack. If I/O paths have $-2.0\,\text{ns}$ WNS, the tool will ignore internal paths with $-0.2\,\text{ns}$ WNS. Grouping paths ensures each category receives dedicated optimization effort.
\end{ansbox}

\begin{qbox}{Q87: Max Transition (Slew) Violations Severity}{\badgeadv}
Why does PrimeTime report high severity for \textbf{Max Transition (Slew)} violations, even when setup and hold slacks are positive?
\begin{enumerate}[label=\textbf{\Alph*.}]
  \item Slow transitions cause large cell propagation delays, excessive dynamic short-circuit power ($I_{sc}$ direct path from $V_{DD}$ to $V_{SS}$ during switching), and extreme susceptibility to crosstalk noise and EM failure.
  \item Max transition violations cause immediate bitstream corruption.
  \item Slew rates only affect scan chain testing.
  \item Max transition violations force the chip into reset mode.
\end{enumerate}
\end{qbox}
\begin{ansbox}
\textbf{Correct Answer: A}

Slew violations mean signal edges rise/fall too slowly. During slow transitions, both NMOS and PMOS networks remain partially ON simultaneously, drawing massive short-circuit currents. Additionally, slow edges are easily perturbed by adjacent aggressor noise.
\end{ansbox}

\begin{qbox}{Q88: Pure Point-to-Point Min/Max Delay Constraints}{\badgeadv}
When are \texttt{set\_max\_delay} and \texttt{set\_min\_delay} constraints used instead of \texttt{create\_clock} and \texttt{set\_input\_delay}?
\begin{enumerate}[label=\textbf{\Alph*.}]
  \item For constraining pure asynchronous feedthroughs, custom analog IP interfaces, or point-to-point asynchronous crossing paths without reference clock periods.
  \item To define the master PLL frequency.
  \item To replace all setup and hold checks in a standard design.
  \item To fix scan chain hold violations automatically.
\end{enumerate}
\end{qbox}
\begin{ansbox}
\textbf{Correct Answer: A}

\texttt{set\_max\_delay} and \texttt{set\_min\_delay} enforce absolute maximum and minimum path delays between startpoints and endpoints without reference to clock edges or setup/hold relationships.
\end{ansbox}

\begin{qbox}{Q89: Half-Cycle Path Available Setup Budget}{\badgeadv}
A path launches on the \textbf{rising edge} of \texttt{CLK} and is captured on the \textbf{falling edge} of the same \texttt{CLK} ($T_{clk} = 8.0\,\text{ns}$, 50\% duty cycle).  
What is the available clock budget for setup analysis (assuming zero skew and $T_{setup} = 0.4\,\text{ns}$)?
\begin{enumerate}[label=\textbf{\Alph*.}]
  \item $8.0\,\text{ns}$
  \item $3.6\,\text{ns}$
  \item $4.0\,\text{ns}$
  \item $7.6\,\text{ns}$
\end{enumerate}
\end{qbox}
\begin{ansbox}
\textbf{Correct Answer: B}

In a rising-to-falling half-cycle path, the capture edge occurs at $T_{clk}/2 = 4.0\,\text{ns}$. The available budget for $T_{cq} + T_{comb}$ is $4.0\,\text{ns} - T_{setup} = 4.0 - 0.4 = 3.6\,\text{ns}$.
\end{ansbox}

\begin{qbox}{Q90: Source Latency vs Network Latency in Clock Distribution}{\badgeadv}
What is the distinction between \textbf{Source Latency} and \textbf{Network Latency} in PrimeTime?
\begin{enumerate}[label=\textbf{\Alph*.}]
  \item Source latency is on-chip clock tree delay; Network latency is PLL lock time.
  \item \textbf{Source Latency} is the delay from the primary clock generator/oscillator off-chip up to the chip's clock input port; \textbf{Network Latency} is the internal delay from the clock input port through the internal clock tree distribution network to the register clock pin.
  \item They are identical terms for wire resistance.
  \item Source latency applies only to reset signals.
\end{enumerate}
\end{qbox}
\begin{ansbox}
\textbf{Correct Answer: B}

Total clock latency = Source Latency (off-chip clock distribution) + Network Latency (internal on-chip clock tree insertion delay). Prior to CTS, both are set via \texttt{set\_clock\_latency}; after CTS, network latency is propagated from physical wires.
\end{ansbox}

\begin{qbox}{Q91: Transitioning to Propagated Clocks Post-CTS}{\badgeadv}
When transitioning from pre-layout synthesis to post-CTS physical timing analysis in PrimeTime, which command must be issued to enable real clock tree delay calculations?
\begin{enumerate}[label=\textbf{\Alph*.}]
  \item \texttt{set\_ideal\_network -all}
  \item \texttt{set\_propagated\_clock [all\_clocks]}
  \item \texttt{remove\_clock\_uncertainty -all}
  \item \texttt{set\_clock\_tree -enable}
\end{enumerate}
\end{qbox}
\begin{ansbox}
\textbf{Correct Answer: B}

In pre-CTS synthesis, clocks are ideal (zero insertion delay). Once the physical clock tree is synthesized in PnR, issuing \texttt{set\_propagated\_clock [all\_clocks]} instructs the timing engine to compute actual network insertion delays, cell delays, and physical clock skew from real wire parasitics.
\end{ansbox}

\begin{qbox}{Q92: Temperature Inversion Phenomenon at Advanced Nodes}{\badgeadv}
At advanced sub-40nm and FinFET process technologies, what is the \textbf{Temperature Inversion} phenomenon, and how does it impact STA signoff?
\begin{enumerate}[label=\textbf{\Alph*.}]
  \item Transistors run faster at high temperatures due to increased thermal energy.
  \item At lower supply voltages ($V_{DD} \approx V_{th}$), the reduction in threshold voltage ($V_{th}$) with decreasing temperature dominates over carrier mobility degradation, making standard cells \textbf{slower at cold temperatures (e.g. $-40^\circ\text{C}$)} than at hot temperatures ($+125^\circ\text{C}$).
  \item Standard cells stop functioning at room temperature.
  \item Temperature inversion only affects analog PLL circuits.
\end{enumerate}
\end{qbox}
\begin{ansbox}
\textbf{Correct Answer: B}

Traditionally, higher temperatures increased delay. In modern low-voltage nodes, mobility loss is outweighed by $V_{th}$ shifts. Consequently, slow-corner setup analysis must be performed at both $-40^\circ\text{C}$ (Cold) and $+125^\circ\text{C}$ (Hot) to catch the true worst-case delay.
\end{ansbox}

\begin{qbox}{Q93: Silicon Failures from Reset Recovery Violations}{\badgeadv}
What happens in silicon if a flip-flop violates its \textbf{Recovery Time} during asynchronous reset deassertion?
\begin{enumerate}[label=\textbf{\Alph*.}]
  \item The flip-flop burns out due to over-current.
  \item The flip-flop may enter a \textbf{metastable state} or experience a 1-cycle output resolution difference relative to other flip-flops, causing state machines to jump into invalid, corrupted states.
  \item The reset pin remains permanently locked to zero.
  \item The clock tree stops toggling.
\end{enumerate}
\end{qbox}
\begin{ansbox}
\textbf{Correct Answer: B}

Recovery time is the setup requirement of an asynchronous pin relative to the clock. Violating recovery time causes internal feedback race conditions in the master/slave stages, producing non-deterministic metastability or staggered state updates across the chip.
\end{ansbox}

\begin{qbox}{Q94: Case Analysis Mode Configuration}{\badgeadv}
What is the purpose of setting \texttt{set\_case\_analysis 0 [get\_pins u\_mux/sel]} in STA?
\begin{enumerate}[label=\textbf{\Alph*.}]
  \item It tests the circuit under all possible random conditions.
  \item It statically ties a pin/port to a constant logic state (\texttt{0} or \texttt{1}), disabling inactive functional paths through multiplexers and configuring the design for a specific operational mode (e.g. Functional Mode vs Scan Test Mode).
  \item It creates a high-impedance state on the bus.
  \item It checks for uninitialized variables.
\end{enumerate}
\end{qbox}
\begin{ansbox}
\textbf{Correct Answer: B}

Designs share logic between functional mode, scan shift, scan capture, and BIST modes. \texttt{set\_case\_analysis} freezes mode-select pins to constant logic values, preventing STA from reporting unrealistic timing paths across inactive modes.
\end{ansbox}

% ==============================================================================
\section{Low Power Design \& UPF}
% ==============================================================================

\begin{qbox}{Q95: CMOS Dynamic Power Formula}{\badgeeasy}
What is the equation for CMOS dynamic switching power dissipation?
\begin{enumerate}[label=\textbf{\Alph*.}]
  \item $P_{dyn} = I_{leak} \times V_{DD}$
  \item $P_{dyn} = \alpha \cdot C_L \cdot V_{DD}^2 \cdot f$
  \item $P_{dyn} = V_{DD} / R$
  \item $P_{dyn} = \alpha \cdot C_L \cdot V_{DD} \cdot f^2$
\end{enumerate}
\end{qbox}
\begin{ansbox}
\textbf{Correct Answer: B}

Dynamic power is consumed when capacitive loads charge and discharge: $P_{dyn} = \alpha \cdot C_L \cdot V_{DD}^2 \cdot f$, where $\alpha$ is switching activity factor, $C_L$ is load capacitance, $V_{DD}$ is supply voltage, and $f$ is operating frequency.
\end{ansbox}

\begin{qbox}{Q96: Static Subthreshold and Tunneling Leakage}{\badgeeasy}
What is the primary source of static (leakage) power in sub-micron CMOS circuits when no switching activity is occurring?
\begin{enumerate}[label=\textbf{\Alph*.}]
  \item Subthreshold leakage current ($I_{sub}$), gate oxide tunneling leakage ($I_{gate}$), and reverse-biased junction band-to-band tunneling ($I_{junc}$).
  \item Clock buffer switching capacitance.
  \item Resistor thermal Johnson noise.
  \item Package pin inductance.
\end{enumerate}
\end{qbox}
\begin{ansbox}
\textbf{Correct Answer: A}

Static leakage occurs even when the clock is stopped. Subthreshold conduction ($V_{gs} < V_{th}$), gate-dielectric quantum tunneling, and junction leakage drain current continuously from $V_{DD}$ to ground.
\end{ansbox}

\begin{qbox}{Q97: Multi-Vt Optimization Strategy}{\badgemod}
How are Multi-$V_t$ libraries (HVT, SVT, LVT) deployed across an ASIC design to minimize total power?
\begin{enumerate}[label=\textbf{\Alph*.}]
  \item LVT cells are placed everywhere to maximize speed.
  \item \textbf{LVT (Low-$V_t$)} cells (fastest switching, highest leakage) are used strictly on timing-critical paths to close setup time; \textbf{HVT (High-$V_t$)} cells (slowest switching, lowest leakage) are used on all non-critical paths with positive slack to suppress leakage power.
  \item HVT cells are used only for I/O pads.
  \item Multi-$V_t$ requires three different power supplies.
\end{enumerate}
\end{qbox}
\begin{ansbox}
\textbf{Correct Answer: B}

High-$V_t$ cells have orders-of-magnitude lower leakage current than Low-$V_t$ cells. Synthesis and PnR tools swap non-critical cells to HVT, maintaining target performance while drastically cutting standby leakage power.
\end{ansbox}

\begin{qbox}{Q98: Level Shifter Cell Insertion Requirements}{\badgemod}
When is a \textbf{Level Shifter} cell required in a Multi-Voltage design?
\begin{enumerate}[label=\textbf{\Alph*.}]
  \item Between two flip-flops on the same voltage domain.
  \item At the boundary where a signal crosses from a \textbf{Low-Voltage Domain ($V_{DDL}$)} to a \textbf{High-Voltage Domain ($V_{DDH}$)} to prevent PMOS transistors in the receiver from staying partially ON and drawing excessive static crowbar current.
  \item Only on clock signals.
  \item Inside the PLL circuit.
\end{enumerate}
\end{qbox}
\begin{ansbox}
\textbf{Correct Answer: B}

When a $0.8\,\text{V}$ signal drives a $1.2\,\text{V}$ gate, a high output ($0.8\,\text{V}$) cannot fully turn off the PMOS transistor (whose gate needs $1.2\,\text{V}$), creating a direct short-circuit path from $V_{DDH}$ to ground. A Low-to-High level shifter shifts the voltage swing to full rail.
\end{ansbox}

\begin{qbox}{Q99: Power Gating Isolation Cells}{\badgemod}
What is the role of an \textbf{Isolation Cell} at the interface of a Power-Gated domain?
\begin{enumerate}[label=\textbf{\Alph*.}]
  \item To step up the supply voltage.
  \item When a power domain is switched OFF ($0\,\text{V}$), its floating output signals are clamped by isolation cells to a known, safe logic state (\texttt{0} or \texttt{1}) before reaching an active "Always-On" domain, preventing floating gates, invalid states, and crowbar currents.
  \item To isolate clock domains from data domains.
  \item To filter electro-magnetic interference.
\end{enumerate}
\end{qbox}
\begin{ansbox}
\textbf{Correct Answer: B}

When power is cut to a power domain, its internal nodes drift into undefined intermediate floating voltages. Isolation cells clamp these boundary signals to constant safe values (\texttt{0} or \texttt{1}) so downstream active domains operate normally.
\end{ansbox}

\begin{qbox}{Q100: State Retention Power Gating (SRPG) Flip-Flops}{\badgemod}
What is the purpose of a \textbf{Retention Flip-Flop (State Retention Power Gating - SRPG)}?
\begin{enumerate}[label=\textbf{\Alph*.}]
  \item To increase the maximum clock frequency.
  \item It contains a secondary "shadow" latch powered by an Always-On ($V_{DD\_AON}$) supply that saves the register's state prior to main power shutdown and restores it immediately upon wake-up, bypassing lengthy system reboot/reconfiguration cycles.
  \item To retain data during scan testing.
  \item To prevent setup violations at low temperatures.
\end{enumerate}
\end{qbox}
\begin{ansbox}
\textbf{Correct Answer: B}

Retention flip-flops preserve register states during sleep mode. Before power-gating $V_{DD}$, a \texttt{SAVE} pulse copies the state into an always-on shadow latch. Upon power-up, a \texttt{RESTORE} pulse copies the state back, enabling near-instantaneous wake-up.
\end{ansbox}

\begin{qbox}{Q101: Unified Power Format (UPF / IEEE 1801) Role}{\badgemod}
What is the role of \textbf{UPF (IEEE 1801)} in the modern digital design flow?
\begin{enumerate}[label=\textbf{\Alph*.}]
  \item It is a programming language used to design analog PLLs.
  \item It provides a standardized TCL-based format to describe the low-power architectural intent (power domains, power switches, isolation strategies, level shifters, retention rules, power state tables) independently of the RTL HDL code.
  \item It replaces SDC timing constraints.
  \item It is used exclusively for packaging thermal simulations.
\end{enumerate}
\end{qbox}
\begin{ansbox}
\textbf{Correct Answer: B}

UPF allows architects to specify power gating, multi-voltage domains, isolation, and level shifting rules in a separate file. EDA tools (simulation, synthesis, formal verification, PnR) apply the UPF file to check and implement power management structures consistently.
\end{ansbox}

\begin{qbox}{Q102: Header (PMOS) vs Footer (NMOS) Power Switches}{\badgeadv}
What is the trade-off between \textbf{Header Switches (PMOS)} and \textbf{Footer Switches (NMOS)} in power-gating architectures?
\begin{enumerate}[label=\textbf{\Alph*.}]
  \item Header switches use less area than footer switches.
  \item \textbf{NMOS Footer switches} ($V_{SS}$ gating) have higher drive current and lower on-resistance ($R_{on}$) per unit area due to higher electron mobility, saving area; \textbf{PMOS Header switches} ($V_{DD}$ gating) preserve substrate ground reference, simplifying noise management and multi-voltage interfaces.
  \item Header switches can only be used with HVT libraries.
  \item Footer switches do not support retention latches.
\end{enumerate}
\end{qbox}
\begin{ansbox}
\textbf{Correct Answer: B}

NMOS electrons have $\approx 2-3\times$ higher mobility than PMOS holes, making NMOS footers physically smaller for the same $I_{on}$ drop. However, footers bounce the virtual ground ($V_{SS\_virtual}$), complicating substrate noise and level shifter design. PMOS headers keep true ground stable.
\end{ansbox}

\begin{qbox}{Q103: Inrush Current and Daisy-Chained Wakeup Sequences}{\badgeadv}
When waking up a massive power-gated domain containing millions of gates:
Why are the sleep transistors turned on in a staggered, \textbf{daisy-chained} sequence rather than simultaneously?
\begin{enumerate}[label=\textbf{\Alph*.}]
  \item To prevent excessive \textbf{Inrush Current} ($di/dt$) that causes severe power grid voltage drops (supply droop) on adjacent Always-On domains, which would otherwise induce false resets or timing failures.
  \item To allow clock trees to cool down thermally.
  \item Because the UPF standard limits power switches to 10 at a time.
  \item To give the PLL time to lock.
\end{enumerate}
\end{qbox}
\begin{ansbox}
\textbf{Correct Answer: A}

Turning on all sleep transistors at once causes a massive current surge to charge depleted power-rail capacitances ($i = C \cdot dv/dt$). This inrush current creates an $L \cdot di/dt + I \cdot R$ voltage droop across the chip's power grid, resetting or corrupting active adjacent blocks. Daisy-chaining buffers turn on switches gradually.
\end{ansbox}

\begin{qbox}{Q104: Dynamic Voltage and Frequency Scaling (DVFS) Power Math}{\badgeadv}
If a processor's supply voltage is reduced by $20\%$ ($V_{new} = 0.8\,V_{old}$) and frequency is reduced by $20\%$ ($f_{new} = 0.8\,f_{old}$), by what percentage is the dynamic power reduced?
\begin{enumerate}[label=\textbf{\Alph*.}]
  \item $20\%$
  \item $36\%$
  \item $48.8\%$
  \item $64\%$
\end{enumerate}
\end{qbox}
\begin{ansbox}
\textbf{Correct Answer: C}

$P_{new} = \alpha C (0.8 V)^2 (0.8 f) = 0.64 \times 0.8 \times P_{old} = 0.512 \, P_{old}$.  
Total Power reduction = $1 - 0.512 = 0.488 = 48.8\%$.
\end{ansbox}

\begin{qbox}{Q105: UPF Power State Table (PST) Verification}{\badgeadv}
Consider the following UPF snippet:
\begin{lstlisting}[style=tclstyle]
create_pst sys_pst -supplies {VDD_CORE VDD_MEM VDD_AON}
add_pst_state S0 -pst sys_pst -state {FULL_ON FULL_ON FULL_ON}
add_pst_state S1 -pst sys_pst -state {OFF     FULL_ON FULL_ON}
add_pst_state S2 -pst sys_pst -state {OFF     OFF     FULL_ON}
\end{lstlisting}
What is the function of the Power State Table (\texttt{sys\_pst}) during verification and synthesis?
\begin{enumerate}[label=\textbf{\Alph*.}]
  \item It defines the physical pad locations of power pins.
  \item It defines all valid combinations of power domain voltage states across operating modes, allowing automated tools to identify illegal voltage states and verify isolation/level-shifting rules at each boundary.
  \item It calculates the total static leakage in watts.
  \item It configures the CPU boot address.
\end{enumerate}
\end{qbox}
\begin{ansbox}
\textbf{Correct Answer: B}

The PST enumerates legal system power states. EDA tools use it to determine whether isolation cells, level shifters, or enable signals are required between domains for every reachable operational mode (e.g. Core Sleep, Memory Retention, Deep Standby).
\end{ansbox}

\begin{qbox}{Q106: Clock Gating Minimum Bit-Width Threshold Efficiency}{\badgeadv}
Why do synthesis tools enforce a minimum bit-width threshold (e.g. \texttt{set\_clock\_gating\_style -min\_bitwidth 4}) before inserting an Integrated Clock Gating (ICG) cell?
\begin{enumerate}[label=\textbf{\Alph*.}]
  \item An ICG cell has an intrinsic cell area and internal clock pin capacitive load; for 1 or 2-bit registers, the power saved by gating the flip-flops is smaller than the overhead power consumed by the ICG itself.
  \item ICG cells cannot drive more than 1 flip-flop.
  \item It violates Verilog language semantics for 1-bit registers.
  \item Narrow registers do not have clock pins.
\end{enumerate}
\end{qbox}
\begin{ansbox}
\textbf{Correct Answer: A}

Inserting an ICG adds silicon area and continuous clock pin switching capacitance. For wide data buses (e.g. 16, 32, 64 bits), gating the bus saves massive power. For 1-2 bits, the overhead of the ICG exceeds the savings.
\end{ansbox}

\begin{qbox}{Q107: Glitch Power Reduction in Deep Arithmetic Pipelines}{\badgeadv}
In deep combinational arithmetic structures (like carry-save multipliers or deep ripple adders), \textbf{glitch power} can account for up to $40\%$ of total dynamic power. What design technique effectively eliminates glitch power?
\begin{enumerate}[label=\textbf{\Alph*.}]
  \item Removing all reset logic.
  \item Pipelining the datapath with balanced register stages to equalize path delays and stop intermediate transitions from propagating down the logic tree.
  \item Operating at higher supply voltages.
  \item Using larger driving cells.
\end{enumerate}
\end{qbox}
\begin{ansbox}
\textbf{Correct Answer: B}

Glitches occur when inputs to a combinational gate arrive at different times due to unbalanced logic paths. Inserting balanced pipeline registers intercepts intermediate hazards, preventing glitches from propagating and charging downstream node capacitances.
\end{ansbox}

\begin{qbox}{Q108: Exponential Dependence of Subthreshold Leakage}{\badgeadv}
Subthreshold leakage current is modeled by $I_{sub} \propto \exp\left(\frac{V_{GS} - V_{th}}{n V_T}\right)$.  
If the threshold voltage $V_{th}$ of a standard cell is reduced by $100\,\text{mV}$ at room temperature ($n V_T \approx 30\,\text{mV}$):
By approximately what factor does the subthreshold leakage current increase?
\begin{enumerate}[label=\textbf{\Alph*.}]
  \item $2\times$
  \item $5\times$
  \item $\approx e^{100/30} = e^{3.33} \approx 28\times$
  \item $100\times$
\end{enumerate}
\end{qbox}
\begin{ansbox}
\textbf{Correct Answer: C}

Subthreshold leakage increases \textbf{exponentially} with decreases in $V_{th}$. A $100\,\text{mV}$ reduction in $V_{th}$ results in an increase of $e^{100/30} \approx 28\times$ in leakage power, illustrating the dramatic standby power penalty of Low-$V_t$ (LVT) cells.
\end{ansbox}

\begin{qbox}{Q109: Isolation Control Signal Sourcing from Always-On Domain}{\badgeadv}
When power-gating a domain \texttt{PD\_CPU}, where must the \texttt{iso\_en} (isolation enable) control signal be generated and powered from?
\begin{enumerate}[label=\textbf{\Alph*.}]
  \item From inside the \texttt{PD\_CPU} domain using its switched power rail.
  \item From an \textbf{Always-On Power Domain ($V_{DD\_AON}$)} so that the isolation control signal remains stable and asserted before, during, and after \texttt{PD\_CPU} is powered down.
  \item From the external JTAG clock.
  \item From the memory bitlines.
\end{enumerate}
\end{qbox}
\begin{ansbox}
\textbf{Correct Answer: B}

If the isolation enable signal originated from within the switched power domain, it would lose power and collapse when the domain is shut down, disabling isolation and allowing floating voltages to escape. Isolation controls must always originate from an Always-On domain.
\end{ansbox}

\begin{qbox}{Q110: Well Biasing (RBB vs FBB) Principles}{\badgeadv}
What is the objective of \textbf{Reverse Body Biasing (RBB)} in low-power standby modes?
\begin{enumerate}[label=\textbf{\Alph*.}]
  \item To increase transistor switching speed during functional mode.
  \item Applying a reverse bias between the transistor body/well and source increases the effective threshold voltage ($V_{th}$) via the body effect, suppressing subthreshold leakage current by up to $10\times$ during sleep.
  \item To reduce the gate oxide thickness.
  \item To eliminate wire resistance.
\end{enumerate}
\end{qbox}
\begin{ansbox}
\textbf{Correct Answer: B}

Reverse Body Biasing (RBB) raises $V_{th}$ when the chip is idle to cut leakage. Conversely, Forward Body Biasing (FBB) lowers $V_{th}$ dynamically during high-performance bursts to maximize clock frequency at low operating voltages.
\end{ansbox}

\begin{qbox}{Q111: Always-On Buffers in Physical Power Routing}{\badgeadv}
What is an \textbf{Always-On (AON) Buffer}, and when is it required during physical implementation?
\begin{enumerate}[label=\textbf{\Alph*.}]
  \item A buffer that switches continuously at $1\,\text{GHz}$.
  \item A special standard cell with a dedicated secondary $V_{DD\_AON}$ power pin, used when an always-on signal must be routed across a physically shut-down power domain floorplan tile.
  \item A buffer connected to an off-chip battery.
  \item A standard inverter used in clock trees.
\end{enumerate}
\end{qbox}
\begin{ansbox}
\textbf{Correct Answer: B}

When routing an active signal across a physical floorplan area whose main power rail may be turned off, standard buffers placed in that tile would lose power. AON buffers connect to a continuous un-switched backup rail, enabling signal feedthrough across sleep regions.
\end{ansbox}

\begin{qbox}{Q112: Power-Aware Static Timing Analysis of Level Shifters}{\badgeadv}
Why must Static Timing Analysis be power-aware when evaluating level-shifter paths?
\begin{enumerate}[label=\textbf{\Alph*.}]
  \item Level shifters have zero propagation delay.
  \item Level shifters introduce non-linear propagation delays that vary dramatically depending on the voltage delta between $V_{DD\_source}$ and $V_{DD\_dest}$; STA must evaluate the worst-case low voltage at source combined with high voltage at capture.
  \item Timing analysis is disabled in low-power domains.
  \item Level shifters eliminate setup checks.
\end{enumerate}
\end{qbox}
\begin{ansbox}
\textbf{Correct Answer: B}

Level shifters operating with low input voltage take significantly longer to pull up/down the internal cross-coupled PMOS latches. STA tools must analyze paths under asymmetric multi-voltage corners ($V_{DDL\_min} \rightarrow V_{DDH\_max}$) to verify timing closure.
\end{ansbox}

% ==============================================================================
\section{Clock Domain Crossing (CDC) \& Async FIFO}
% ==============================================================================

\begin{qbox}{Q113: Metastability Phenomenon in Digital Flip-Flops}{\badgeeasy}
What physical phenomenon occurs when an asynchronous signal violates setup or hold time at a destination flip-flop?
\begin{enumerate}[label=\textbf{\Alph*.}]
  \item The flip-flop immediately catches fire.
  \item The flip-flop output enters a non-deterministic \textbf{Metastable State}, hovering between legal logic levels \texttt{0} and \texttt{1} for an unbounded duration before settling randomly.
  \item The flip-flop acts as a frequency divider.
  \item The signal voltage doubles.
\end{enumerate}
\end{qbox}
\begin{ansbox}
\textbf{Correct Answer: B}

Violating setup or hold prevents internal feedback latches from resolving cleanly, leaving the internal nodes balanced at the threshold voltage. The output oscillates or hovers at an invalid intermediate voltage before thermal noise resolves it to \texttt{0} or \texttt{1}.
\end{ansbox}

\begin{qbox}{Q114: 2-Flip-Flop Synchronizer Role}{\badgeeasy}
What is the primary function of a standard \textbf{2-Flip-Flop (2-FF) Synchronizer}?
\begin{enumerate}[label=\textbf{\Alph*.}]
  \item To synchronize multi-bit parallel data buses.
  \item To provide a settling time of one full destination clock period ($T_{clk} - T_{setup}$) for a 1-bit asynchronous control signal to resolve metastability, reducing the probability of system failure.
  \item To double the clock frequency.
  \item To eliminate propagation delay completely.
\end{enumerate}
\end{qbox}
\begin{ansbox}
\textbf{Correct Answer: B}

If the first flip-flop enters metastability due to an asynchronous transition, the second flip-flop samples its output one full clock cycle later, giving the first stage time to resolve to a stable binary logic level.
\end{ansbox}

\begin{qbox}{Q115: Mean Time Between Failures (MTBF) Scaling}{\badgemod}
The MTBF for a synchronizer flip-flop is given by $\text{MTBF} = \frac{e^{t_{res} / \tau}}{T_0 \cdot f_{clk} \cdot f_{data}}$.  
To increase the MTBF from years to centuries:
\begin{enumerate}[label=\textbf{\Alph*.}]
  \item Increase the data frequency $f_{data}$.
  \item Increase the resolution settling time $t_{res}$ (e.g. by using a 3-FF synchronizer or slower clock) and use flip-flops with small technology time constant $\tau$.
  \item Decrease the clock period.
  \item Increase the operating temperature.
\end{enumerate}
\end{qbox}
\begin{ansbox}
\textbf{Correct Answer: B}

Because $t_{res}$ appears in the exponent ($e^{t_{res}/\tau}$), adding an extra synchronizer flip-flop (increasing resolution time by a full clock cycle) exponentially increases the MTBF, preventing system failures.
\end{ansbox}

\begin{qbox}{Q116: Fast-to-Slow Pulse Synchronizer Architecture}{\badgemod}
Why will a 1-clock-cycle wide pulse in a Fast Clock Domain ($1\,\text{GHz}$) be missed if sampled directly by a 2-FF synchronizer in a Slow Clock Domain ($50\,\text{MHz}$)?
\begin{enumerate}[label=\textbf{\Alph*.}]
  \item The fast clock frequency causes electromagnetic interference.
  \item The pulse width ($1\,\text{ns}$) is much shorter than the slow clock period ($20\,\text{ns}$); if the pulse rises and falls between two consecutive slow clock edges, the synchronizer flip-flops never capture the pulse.
  \item 2-FF synchronizers only work from slow to fast domains.
  \item Slow domains reject all external pulses.
\end{enumerate}
\end{qbox}
\begin{ansbox}
\textbf{Correct Answer: B}

For a 2-FF synchronizer to capture a signal, the input must remain stable for at least $1.5 \times T_{dest\_clk}$ (Open-Loop rule). A fast pulse must be converted into a toggle signal (using a Toggle Synchronizer) or held by a handshake before being passed to the slow domain.
\end{ansbox}

\begin{qbox}{Q117: Multi-Bit CDC Synchronization Failure}{\badgemod}
Why is passing a multi-bit binary counter across asynchronous clock domains using separate 2-FF synchronizers on each bit fatal?
\begin{enumerate}[label=\textbf{\Alph*.}]
  \item It consumes too much power.
  \item Due to wire routing skew and individual bit metastability resolution, when transitioning from \texttt{0111} ($7$) to \texttt{1000} ($8$), the destination domain may capture intermediate transient values like \texttt{1111} ($15$) or \texttt{0000} ($0$), corrupting system state.
  \item Binary counters cannot be synthesized.
  \item 2-FF synchronizers invert every odd bit.
\end{enumerate}
\end{qbox}
\begin{ansbox}
\textbf{Correct Answer: B}

Multiple bits changing simultaneously will arrive at slightly different times at the destination registers due to wire delay skews. The destination synchronizer will sample a mixture of old and new bits, creating severe false data words.
\end{ansbox}

\begin{qbox}{Q118: Gray Code Single-Bit Transition Property in CDC}{\badgemod}
Why is \textbf{Gray Code} universally used for FIFO pointer synchronization across asynchronous clock domains?
\begin{enumerate}[label=\textbf{\Alph*.}]
  \item Gray code uses fewer bits than binary.
  \item Exactly \textbf{one bit changes per increment/decrement transition}, ensuring that even with arbitrary clock edge alignment and routing skew, the synchronizer samples either the previous pointer or the new pointer, with zero possibility of corrupted intermediate values.
  \item Gray code operates without a clock signal.
  \item Gray code automatically clears FIFO full flags.
\end{enumerate}
\end{qbox}
\begin{ansbox}
\textbf{Correct Answer: B}

Because only a single bit toggles between adjacent Gray code numbers ($00 \rightarrow 01 \rightarrow 11 \rightarrow 10$), there are no intermediate transition states. The synchronizer either captures the old pointer value or the new one---both of which are safe and valid.
\end{ansbox}

\begin{qbox}{Q119: Reset Synchronizer Asynchronous Assert / Synchronous Deassert}{\badgemod}
Why does an industry-standard \textbf{Reset Synchronizer} assert reset asynchronously but deassert it synchronously?
\begin{enumerate}[label=\textbf{\Alph*.}]
  \item To save silicon area on the reset pad.
  \item Asynchronous assertion guarantees instant shutdown upon power-on or fault even if the clock is stopped; Synchronous deassertion guarantees that reset release satisfies flip-flop recovery/removal times relative to the active clock edge.
  \item To allow the CPU to run without a reset signal.
  \item To eliminate reset buffers.
\end{enumerate}
\end{qbox}
\begin{ansbox}
\textbf{Correct Answer: B}

Immediate asynchronous assertion protects the circuit without waiting for a clock. Releasing the reset through 2 flip-flops clocked by the target clock ensures the deassert edge is aligned with the clock, preventing recovery timing violations.
\end{ansbox}

\begin{qbox}{Q120: CDC Re-convergence Hazards}{\badgemod}
What is \textbf{CDC Re-convergence}, and why does SpyGlass CDC flag it as an error?
\begin{enumerate}[label=\textbf{\Alph*.}]
  \item When a clock tree divides into two branches.
  \item When two independently synchronized control signals originating from the same source domain are combined combinational-wise in the destination domain; cycle uncertainties ($0$ or $1$ cycle delay in each synchronizer) can cause the signals to arrive out-of-sync, corrupting the decoded state.
  \item When a reset signal drives a clock pin.
  \item When two PLLs are connected in series.
\end{enumerate}
\end{qbox}
\begin{ansbox}
\textbf{Correct Answer: B}

Even if two 1-bit signals are individually synchronized via 2-FF synchronizers, one may resolve in 1 cycle while the other resolves in 2 cycles. Combining them with logic in the destination domain produces invalid decoded states. They must be merged into a single state/handshake before crossing.
\end{ansbox}

\begin{qbox}{Q121: Asynchronous FIFO Full and Empty Domains}{\badgeadv}
In a dual-clock Asynchronous FIFO with synchronized Gray pointers:
In which clock domain are the \textbf{FIFO Empty} and \textbf{FIFO Full} flags evaluated?
\begin{enumerate}[label=\textbf{\Alph*.}]
  \item \texttt{Empty} in Write domain; \texttt{Full} in Read domain
  \item \textbf{\texttt{Empty} in Read Clock Domain} (comparing local read Gray pointer with synchronized write Gray pointer); \textbf{\texttt{Full} in Write Clock Domain} (comparing local write Gray pointer with synchronized read Gray pointer).
  \item Both flags are evaluated in a third asynchronous monitoring domain.
  \item Both flags are evaluated purely combinationally without clocks.
\end{enumerate}
\end{qbox}
\begin{ansbox}
\textbf{Correct Answer: B}

\texttt{Empty} controls reading, so it must be evaluated in the Read domain with zero latency to prevent reading an empty FIFO. \texttt{Full} controls writing, so it must be evaluated in the Write domain with zero latency to prevent overwriting a full FIFO.
\end{ansbox}

\begin{qbox}{Q122: Asynchronous FIFO Full Flag Boolean Condition}{\badgeadv}
For an Async FIFO of depth $2^N$ using $(N+1)$-bit Gray code pointers $W_{ptr}$ and $R_{ptr\_sync}$:  
What is the exact condition that indicates the FIFO is \textbf{Full}?
\begin{enumerate}[label=\textbf{\Alph*.}]
  \item $W_{ptr} == R_{ptr\_sync}$
  \item $W_{ptr}[N:N-1] == \sim R_{ptr\_sync}[N:N-1]$ and $W_{ptr}[N-2:0] == R_{ptr\_sync}[N-2:0]$ (MSB and 2nd MSB inverted, all remaining LSBs identical).
  \item $W_{ptr} - R_{ptr\_sync} == 2^N$
  \item $W_{ptr}[N] == R_{ptr\_sync}[N]$
\end{enumerate}
\end{qbox}
\begin{ansbox}
\textbf{Correct Answer: B}

Using an extra pointer bit ($(N+1)$ bits for depth $2^N$), when pointers wrap around (write has lapped read by exactly $2^N$), the MSB and 2nd MSB of the Gray code pointers are inverted ($\sim$), while the remaining lower $N-1$ bits match identically. (If all bits match, the FIFO is Empty).
\end{ansbox}

\begin{qbox}{Q123: Pessimistic Safe Behavior of Async FIFO Flags}{\badgeadv}
Because pointer synchronization across clock domains introduces a 2-cycle latency:
What is the operational safety implication on the \texttt{Full} and \texttt{Empty} flags in an Async FIFO?
\begin{enumerate}[label=\textbf{\Alph*.}]
  \item The FIFO may overflow or underflow by 2 words.
  \item The flags are \textbf{pessimistic (safe)}: \texttt{Full} may stay asserted for 2 extra write cycles after a read occurs (temporarily stalling writes when not strictly full), and \texttt{Empty} may stay asserted for 2 extra read cycles after a write occurs; but the FIFO will \textbf{never overflow or underflow}.
  \item The FIFO drops random data packets.
  \item The write pointer resets to zero automatically.
\end{enumerate}
\end{qbox}
\begin{ansbox}
\textbf{Correct Answer: B}

The write domain sees a delayed version of the read pointer, so it believes the FIFO has fewer free spaces than it actually does (safe). Similarly, the read domain sees a delayed write pointer, believing fewer words are available (safe).
\end{ansbox}

\begin{qbox}{Q124: Asynchronous FIFO Burst Depth Calculation}{\badgeadv}
Calculate the minimum Async FIFO depth for the following specifications:
\begin{itemize}[leftmargin=1.5em]
  \item \textbf{Write Clock}: $f_w = 200\,\text{MHz}$ ($T_w = 5\,\text{ns}$)
  \item \textbf{Read Clock}: $f_r = 50\,\text{MHz}$ ($T_r = 20\,\text{ns}$)
  \item \textbf{Burst Length}: $B = 80$ contiguous data words
  \item \textbf{Read Duty}: 1 read operation every clock cycle ($1\,\text{read} / 20\,\text{ns}$)
  \item \textbf{Write Duty}: 1 write operation every clock cycle ($1\,\text{write} / 5\,\text{ns}$)
\end{itemize}
What is the minimum FIFO depth required to prevent data loss during the burst?
\begin{enumerate}[label=\textbf{\Alph*.}]
  \item 20 words
  \item 40 words
  \item 60 words
  \item 80 words
\end{enumerate}
\end{qbox}
\begin{ansbox}
\textbf{Correct Answer: C}

\begin{enumerate}[leftmargin=1.5em]
  \item Time to write burst = $B \times T_w = 80 \times 5\,\text{ns} = 400\,\text{ns}$.
  \item Number of reads during burst = $\text{Time} / T_r = 400\,\text{ns} / 20\,\text{ns} = 20$ words read out.
  \item Minimum FIFO Depth = $\text{Burst} - \text{Reads} = 80 - 20 = 60$ words.  
  (Standard power-of-2 implementation uses Depth = 64).
\end{enumerate}
\end{ansbox}

\begin{qbox}{Q125: Non-Power-of-2 FIFO Depth Gray Code Pitfall}{\badgeadv}
Why is implementing an Asynchronous FIFO with a non-power-of-2 depth (e.g. Depth = 10) dangerous with standard Gray code pointers?
\begin{enumerate}[label=\textbf{\Alph*.}]
  \item Non-power-of-2 memories cannot be manufactured in silicon.
  \item Standard Gray code sequences only exhibit single-bit transitions when wrapping around from $2^N - 1$ back to $0$. Wrapping around at an arbitrary number (e.g. from index 9 back to 0) causes \textbf{multiple bits to toggle simultaneously}, destroying the single-bit CDC property and causing metastability bugs.
  \item Synthesis tools reject non-power-of-2 arrays.
  \item It requires 5-stage synchronizers.
\end{enumerate}
\end{qbox}
\begin{ansbox}
\textbf{Correct Answer: B}

Gray code symmetry is preserved only across $2^N$ counts. Wrapping from $9$ (\texttt{4'b1101}) to $0$ (\texttt{4'b0000}) changes 3 bits in a single step, re-introducing multi-bit CDC race conditions. Specialized even-count truncated Gray sequences are mandatory for non-power-of-2 depths.
\end{ansbox}

\begin{qbox}{Q126: 4-Phase vs 2-Phase CDC Handshake Protocols}{\badgeadv}
What is the difference between \textbf{4-Phase (Level-based)} and \textbf{2-Phase (Transition-based)} CDC handshakes?
\begin{enumerate}[label=\textbf{\Alph*.}]
  \item 4-phase handshakes require 4 data wires; 2-phase handshakes require 2 data wires.
  \item \textbf{4-Phase Handshake}: \texttt{REQ} asserts $\rightarrow$ \texttt{ACK} asserts $\rightarrow$ \texttt{REQ} deasserts $\rightarrow$ \texttt{ACK} deasserts (returns to zero, taking longer latency); \textbf{2-Phase Handshake}: Any toggle on \texttt{REQ} triggers a transfer, and any toggle on \texttt{ACK} confirms it (faster throughput, higher logic complexity).
  \item 4-phase is only for synchronous domains.
  \item 2-phase handshakes do not require synchronizers.
\end{enumerate}
\end{qbox}
\begin{ansbox}
\textbf{Correct Answer: B}

4-phase handshakes require a full return-to-zero cycle for both \texttt{REQ} and \texttt{ACK} (4 crossing transitions per transaction). 2-phase (Non-Return-to-Zero) handshakes treat each transition (0-to-1 or 1-to-0) as an active event, cutting transaction latency in half.
\end{ansbox}

\begin{qbox}{Q127: Fast-to-Slow Domain MUX-Data Handshake Synchronizer}{\badgeadv}
In a MUX-Data CDC synchronizer (Data-Load / Handshake Enable):
How is a 64-bit multi-bit datapath safely transferred from a fast to a slow domain without Gray coding?
\begin{enumerate}[label=\textbf{\Alph*.}]
  \item The data bus is routed directly to the destination registers, while an independent \texttt{data\_valid} pulse is stretched, synchronized through a 2-FF synchronizer to the slow domain, and used as the MUX/Clock-Enable to safely sample the stable data bus.
  \item All 64 bits are passed through 64 independent 2-FF synchronizers.
  \item By reducing the fast clock frequency to zero.
  \item By converting the 64-bit data to serial UART data.
\end{enumerate}
\end{qbox}
\begin{ansbox}
\textbf{Correct Answer: A}

The transmitter places data on the bus and asserts a control flag. The data bus remains completely static while the control flag is synchronized via 2 flip-flops. Once the synchronized enable arrives in the destination domain, the destination register samples the stable bus with zero metastability risk.
\end{ansbox}

\begin{qbox}{Q128: SpyGlass CDC Rules Ac\_unsync01 and Ac\_conv01}{\badgeadv}
Match the SpyGlass CDC violation rules:
1. \texttt{Ac\_unsync01} \quad 2. \texttt{Ac\_conv01}
\begin{enumerate}[label=\textbf{\Alph*.}]
  \item (1) Unsynchronized crossing from source to destination domain; (2) Re-convergence of synchronized signals in destination domain.
  \item (1) Latch inferred; (2) Reset glitch.
  \item (1) Clock domain mismatch; (2) Scan chain bypass.
  \item (1) Multi-driven net; (2) Combinational loop.
\end{enumerate}
\end{qbox}
\begin{ansbox}
\textbf{Correct Answer: A}

\texttt{Ac\_unsync01} flags direct structural CDC crossings lacking a valid synchronizer (2-FF, FIFO, handshake). \texttt{Ac\_conv01} flags multiple synchronized signals converging into combinational logic in the destination domain.
\end{ansbox}

\begin{qbox}{Q129: Reset Domain Crossing (RDC) Metastability}{\badgeadv}
What is a \textbf{Reset Domain Crossing (RDC)} bug, and how does it occur even when all flip-flops operate on the exact same synchronous clock?
\begin{enumerate}[label=\textbf{\Alph*.}]
  \item When a flip-flop clock is disconnected.
  \item When Flip-Flop A is reset by \texttt{rst\_a\_n} while communicating with Flip-Flop B which is reset by \texttt{rst\_b\_n}; when \texttt{rst\_a\_n} asserts asynchronously, the sudden transition on Flip-Flop A's output can violate setup/hold times of Flip-Flop B, inducing metastability.
  \item When a reset signal has high fanout.
  \item When synchronous resets are used instead of asynchronous resets.
\end{enumerate}
\end{qbox}
\begin{ansbox}
\textbf{Correct Answer: B}

Asynchronous reset assertions occur independently of the clock edge. If two registers share a clock but have independent reset domains, asserting reset on the transmitter acts as an asynchronous data transition into the receiver, causing RDC metastability.
\end{ansbox}

\begin{qbox}{Q130: SDC set\_clock\_groups -asynchronous Meaning}{\badgeadv}
What is the effect of issuing the following SDC constraint in PrimeTime:
\begin{lstlisting}[style=tclstyle]
set_clock_groups -asynchronous -group [get_clocks CLK_A] -group [get_clocks CLK_B]
\end{lstlisting}
\begin{enumerate}[label=\textbf{\Alph*.}]
  \item It aligns the phase of \texttt{CLK\_A} and \texttt{CLK\_B}.
  \item It instructs the Static Timing Analysis engine to completely disable all setup and hold timing checks between all paths crossing between domain \texttt{CLK\_A} and domain \texttt{CLK\_B}.
  \item It inserts 2-FF synchronizers automatically during synthesis.
  \item It generates a clock divider.
\end{enumerate}
\end{qbox}
\begin{ansbox}
\textbf{Correct Answer: B}

\texttt{set\_clock\_groups -asynchronous} informs the STA tool that no fixed phase relationship exists between \texttt{CLK\_A} and \texttt{CLK\_B}. STA cuts all inter-clock timing paths. (The designer must guarantee structural CDC synchronizers are in place via SpyGlass CDC).
\end{ansbox}

\begin{qbox}{Q131: 1-Deep Async FIFO Glitch and Throughput Collapse}{\badgeadv}
Why is an Asynchronous FIFO of Depth = 1 impossible to build reliably using standard Gray pointer synchronization?
\begin{enumerate}[label=\textbf{\Alph*.}]
  \item Depth 1 requires negative bits.
  \item Pointer synchronization requires at least 2 flip-flop delays ($2 \times T_{dest\_clk}$); in a 1-deep FIFO, the \texttt{Full} flag cannot clear in time to allow back-to-back writes, reducing throughput to near zero, and Gray pointers cannot distinguish empty from full with a 1-bit counter.
  \item Dual-port RAMs do not support 1 address location.
  \item Synthesis tools optimize Depth 1 into an inverter.
\end{enumerate}
\end{qbox}
\begin{ansbox}
\textbf{Correct Answer: B}

Gray code pointer comparison requires at least 2 addresses ($N+1$ pointer bits) to differentiate full from empty states. Furthermore, synchronizer latency requires a minimum depth (typically $\ge 4$ or 8 words) to sustain continuous pipelined throughput.
\end{ansbox}

\begin{qbox}{Q132: Combinational Glitches on Synchronizer Inputs}{\badgeadv}
If combinational logic is placed between the transmitting register and the input of a 2-FF synchronizer in an asynchronous clock domain crossing:
Why is this an extreme CDC violation?
\begin{enumerate}[label=\textbf{\Alph*.}]
  \item It increases cell area.
  \item Combinational logic generates dynamic hazards and intermediate glitches; if an asynchronous clock edge samples during an intermediate glitch, the synchronizer will capture and propagate a false pulse to the destination domain.
  \item The synthesis tool removes combinational logic on CDC paths.
  \item It inverts clock polarity.
\end{enumerate}
\end{qbox}
\begin{ansbox}
\textbf{Correct Answer: B}

Inputs to synchronizers must be \textbf{clean registered outputs} directly from flip-flop Q pins. Combinational hazards (glitches) that would normally settle harmlessly before a synchronous clock edge can be captured by an asynchronous clock edge, causing catastrophic false transitions.
\end{ansbox}

\begin{qbox}{Q133: Quasi-Static Signals Handling in CDC Verification}{\badgeadv}
What is a \textbf{Quasi-Static Signal} in CDC verification, and how is it handled?
\begin{enumerate}[label=\textbf{\Alph*.}]
  \item A signal that toggles at $10\,\text{GHz}$.
  \item A software configuration signal written during system initialization and held constant during normal operation; it does not require a 2-FF synchronizer, but must be formally constrained/waived in SpyGlass CDC as \texttt{quasi\_static}.
  \item A clock signal generated by an RC oscillator.
  \item A defective floating net.
\end{enumerate}
\end{qbox}
\begin{ansbox}
\textbf{Correct Answer: B}

Quasi-static signals (e.g., debug mode enable, baud rate configuration divisor) never switch during active datapath transactions. They do not need hardware synchronizers, but must be declared to CDC linting tools to suppress false violation reports.
\end{ansbox}

\begin{qbox}{Q134: Async FIFO Read-to-Write Pipelined Latency}{\badgeadv}
In an Async FIFO, when the read pointer is synchronized to the write clock domain, how many write clock cycles minimum does it take for the write domain to realize that a read has occurred?
\begin{enumerate}[label=\textbf{\Alph*.}]
  \item 0 cycles
  \item Exactly 1 cycle
  \item 2 to 3 write clock cycles (due to the 2-FF synchronizer stages plus clock phase alignment)
  \item 100 cycles
\end{enumerate}
\end{qbox}
\begin{ansbox}
\textbf{Correct Answer: C}

The Gray read pointer must pass through the 2-stage synchronizer clocked by \texttt{wr\_clk}. Depending on when the read pointer changes relative to the next \texttt{wr\_clk} active edge, the update is registered after 2 to 3 full \texttt{wr\_clk} periods.
\end{ansbox}

\begin{qbox}{Q135: Toggle Synchronizer Architecture Protocol}{\badgeadv}
How does a \textbf{Toggle Synchronizer} successfully convey a 1-cycle pulse from a fast clock domain into a slow clock domain?
\begin{enumerate}[label=\textbf{\Alph*.}]
  \item It amplifies the voltage of the pulse.
  \item In the fast domain, every incoming pulse inverts (toggles) the state of a T-flip-flop (\texttt{0 -> 1} or \texttt{1 -> 0}); the continuous level transition is safely synchronized via a 2-FF synchronizer in the slow domain, where an XOR edge detector recreates a 1-cycle pulse.
  \item It stores the pulse in an analog capacitor.
  \item It converts the pulse into a PWM waveform.
\end{enumerate}
\end{qbox}
\begin{ansbox}
\textbf{Correct Answer: B}

Converting a narrow transient pulse into a permanent level toggle converts the high-frequency event into a static DC level. The destination domain synchronizes the level and uses an XOR gate ($Q_{sync} \oplus Q_{sync\_delayed}$) to regenerate a clean 1-cycle destination pulse.
\end{ansbox}

\begin{qbox}{Q136: SDC set\_max\_delay -datapath\_only on CDC Nets}{\badgeadv}
Why is \texttt{set\_max\_delay -datapath\_only} specified on asynchronous CDC crossing paths instead of a standard \texttt{set\_max\_delay}?
\begin{enumerate}[label=\textbf{\Alph*.}]
  \item To force the tool to place all cells in the center of the chip.
  \item \texttt{-datapath\_only} ignores clock skew and source/capture clock latencies (which are meaningless between asynchronous clocks) and strictly constrains the maximum routing wire delay across the asynchronous boundary to prevent excessive bit skew.
  \item To disable power analysis.
  \item To convert flip-flops into latches.
\end{enumerate}
\end{qbox}
\begin{ansbox}
\textbf{Correct Answer: B}

Standard \texttt{set\_max\_delay} includes clock arrival times, which are undefined between asynchronous clocks. Specifying \texttt{-datapath\_only} instructs STA to bound purely the combinational and routing propagation delay between the domains, keeping bus skew within safe multi-bit margins.
\end{ansbox}

% ==============================================================================
\section{Formal Verification \& SystemVerilog Assertions (SVA)}
% ==============================================================================

\begin{qbox}{Q137: Logic Equivalence Checking (LEC) Goal}{\badgeeasy}
What is the primary objective of Logic Equivalence Checking (e.g. Synopsys Formality, Cadence Conformal)?
\begin{enumerate}[label=\textbf{\Alph*.}]
  \item To measure dynamic power dissipation during functional simulation.
  \item To mathematically prove that an \textbf{Implementation Netlist} (e.g. Synthesized or PnR netlist) has identical Boolean functionality to the \textbf{Reference Design} (e.g. Golden RTL) for all possible input combinations without applying test vectors.
  \item To calculate the physical silicon yield.
  \item To test for stuck-at physical fabrication defects.
\end{enumerate}
\end{qbox}
\begin{ansbox}
\textbf{Correct Answer: B}

Formal LEC is an exhaustive mathematical proof comparing two design representations (RTL vs Netlist, or Pre-PnR vs Post-PnR Netlist). It divides logic into combinatorial cones between register/port compare points and proves Boolean equivalence across all $2^N$ input states without running testbenches.
\end{ansbox}

\begin{qbox}{Q138: Formality Compare Points and Logic Cones}{\badgemod}
In Synopsys Formality, what structures are automatically extracted as \textbf{Compare Points} between Reference and Implementation designs?
\begin{enumerate}[label=\textbf{\Alph*.}]
  \item Power ground pads.
  \item Primary Inputs, Primary Outputs, Sequential Elements (Flip-Flops / Latches), and Black Boxes.
  \item Internal wire names.
  \item Clock buffers only.
\end{enumerate}
\end{qbox}
\begin{ansbox}
\textbf{Correct Answer: B}

LEC tools break complex circuits into independent logic cones bounded by \textbf{Compare Points}: Primary Inputs, Primary Outputs, Flip-Flops, Latches, and Black Box boundaries. Equivalence is proven independently for each logic cone.
\end{ansbox}

\begin{qbox}{Q139: Immediate vs Concurrent SystemVerilog Assertions}{\badgemod}
What is the key distinction between \textbf{Immediate Assertions} and \textbf{Concurrent Assertions} in SystemVerilog Assertions (SVA)?
\begin{enumerate}[label=\textbf{\Alph*.}]
  \item Immediate assertions consume time; Concurrent assertions do not.
  \item \textbf{Immediate assertions} (\texttt{assert (a == b)}) evaluate combinationally in the active region like procedural \texttt{if} statements and are prone to simulation glitches; \textbf{Concurrent assertions} (\texttt{assert property (@(posedge clk) ...)}) are clock-driven, multi-cycle temporal checks evaluated in the \textbf{Observed} simulation region.
  \item Concurrent assertions can only be used in C++ testbenches.
  \item Immediate assertions are only evaluated at time 0.
\end{enumerate}
\end{qbox}
\begin{ansbox}
\textbf{Correct Answer: B}

Immediate assertions execute procedurally within simulation blocks and can trigger false glitch warnings during intra-cycle transitions. Concurrent assertions sample signals deterministically in the Preponed region and evaluate temporal properties over multiple clock cycles in the Observed region.
\end{ansbox}

\begin{qbox}{Q140: SVA Implication Operators (|-> vs |=>)}{\badgemod}
In SystemVerilog Assertions, what is the operational difference between \texttt{req |-> ack} and \texttt{req |=> ack}?
\begin{enumerate}[label=\textbf{\Alph*.}]
  \item \texttt{|->} checks \texttt{ack} in the next clock cycle; \texttt{|=>} checks \texttt{ack} in the same clock cycle.
  \item \textbf{Overlapped (\texttt{|->})}: If \texttt{req} is true on clock cycle $N$, \texttt{ack} must be true on the \textbf{same clock cycle $N$}; \textbf{Non-Overlapped (\texttt{|=>})}: If \texttt{req} is true on clock cycle $N$, \texttt{ack} must be true on the \textbf{subsequent clock cycle $N+1$} (equivalent to \texttt{req \#\#1 ack}).
  \item Both operators are completely identical.
  \item \texttt{|=>} is an arithmetic assignment operator.
\end{enumerate}
\end{qbox}
\begin{ansbox}
\textbf{Correct Answer: B}

The overlapped implication operator \texttt{|->} evaluates the consequent in the same cycle that the antecedent succeeds. The non-overlapped operator \texttt{|=>} evaluates the consequent exactly one clock cycle later.
\end{ansbox}

\begin{qbox}{Q141: Resolving Unmapped Points in LEC with SVF Files}{\badgemod}
What causes \textbf{Unmapped Points} during a Formality LEC run, and what is the primary diagnosis step?
\begin{enumerate}[label=\textbf{\Alph*.}]
  \item All standard cells are missing from the technology library.
  \item Registers were renamed during synthesis (e.g. Register Retiming, FSM re-encoding, or unused register removal); the engineer must load a \texttt{.svf} (Synchronized Verification Format) guidance file from Design Compiler to provide name-mapping hints.
  \item The clock frequency was set too high.
  \item The netlist has too many power domains.
\end{enumerate}
\end{qbox}
\begin{ansbox}
\textbf{Correct Answer: B}

When synthesis performs aggressive optimizations (state encoding, register merging, boundary optimization, retiming), register names change. If Formality cannot automatically pair reference registers with implementation registers, it reports Unmapped Points. Loading the \texttt{.svf} file supplies transformation records that guide point alignment.
\end{ansbox}

\begin{qbox}{Q142: SVA \$past and \$rose Sampled Value Functions}{\badgeadv}
Consider the following SVA property:
\begin{lstlisting}[style=verilogstyle]
property p_check_ack;
    @(posedge clk) disable iff (!rst_n)
    $rose(ack) |-> ($past(req, 2) == 1'b1);
endproperty
assert property (p_check_ack);
\end{lstlisting}
What functional behavior does this assertion mathematically enforce?
\begin{enumerate}[label=\textbf{\Alph*.}]
  \item Whenever \texttt{ack} is low, \texttt{req} must be high 2 cycles later.
  \item Whenever \texttt{ack} transitions from \texttt{0} to \texttt{1} (rising edge), the signal \texttt{req} must have been high exactly 2 clock cycles prior.
  \item \texttt{ack} and \texttt{req} must toggle simultaneously every 2 cycles.
  \item \texttt{req} must remain high for 2 consecutive cycles after \texttt{ack} falls.
\end{enumerate}
\end{qbox}
\begin{ansbox}
\textbf{Correct Answer: B}

\texttt{\$rose(ack)} detects a 0-to-1 transition on \texttt{ack}. The antecedent triggers on this edge, requiring \texttt{\$past(req, 2)} (the value of \texttt{req} sampled 2 clock cycles in the past) to equal \texttt{1'b1}.
\end{ansbox}

\begin{qbox}{Q143: SVA Consecutive Repetition Operator [*n]}{\badgeadv}
What does the SVA sequence \texttt{a \#\#1 b [*3] \#\#1 c} verify?
\begin{enumerate}[label=\textbf{\Alph*.}]
  \item Signal \texttt{a} is high, followed 1 cycle later by signal \texttt{b} being high for \textbf{3 consecutive clock cycles}, followed 1 cycle later by signal \texttt{c} being high.
  \item Signals \texttt{a}, \texttt{b}, and \texttt{c} are multiplied together.
  \item Signal \texttt{b} toggles 3 times in 1 cycle.
  \item Signal \texttt{a} is repeated 3 times.
\end{enumerate}
\end{qbox}
\begin{ansbox}
\textbf{Correct Answer: A}

\texttt{b [*3]} is the consecutive repetition operator, denoting that condition \texttt{b} must hold true for 3 successive clock cycles: \texttt{b \#\#1 b \#\#1 b}.
\end{ansbox}

\begin{qbox}{Q144: SVA Non-Consecutive [=n] vs Goto [->n] Repetition}{\badgeadv}
What is the difference between Goto Repetition \texttt{a [->3]} and Non-Consecutive Repetition \texttt{a [=3]}?
\begin{enumerate}[label=\textbf{\Alph*.}]
  \item \texttt{[->3]} requires \texttt{a} to occur on consecutive cycles; \texttt{[=3]} does not.
  \item \textbf{Goto Repetition (\texttt{a [->3]})}: The sequence completes on the \textbf{exact clock cycle} where \texttt{a} is true for the 3rd time; \textbf{Non-Consecutive Repetition (\texttt{a [=3]})}: The sequence matches when \texttt{a} has occurred 3 times and can continue to match on subsequent cycles as long as \texttt{a} remains false before the next event.
  \item Both are identical in formal verification.
  \item \texttt{[=3]} is only used in testbench scoreboards.
\end{enumerate}
\end{qbox}
\begin{ansbox}
\textbf{Correct Answer: B}

\texttt{[->1]} (goto) matches at the precise cycle where the boolean becomes true (blocking further advancement until the match cycle). \texttt{[=1]} allows arbitrary empty cycles after the match has occurred before the next sequence element starts.
\end{ansbox}

\begin{qbox}{Q145: Bounded Model Checking (BMC) Bound Depth Meaning}{\badgeadv}
In Formal Property Verification, what does it mean when a tool reports that an assertion is \textbf{Proven up to Bound $K=50$ (BMC)}?
\begin{enumerate}[label=\textbf{\Alph*.}]
  \item The property is mathematically proven for infinite time ($K=\infty$).
  \item The tool mathematically proved that no bug can violate the assertion within 50 clock cycles from reset; however, a bug may still exist at cycle 51 or beyond unless a full mathematical induction proof (k-induction) is completed.
  \item The testbench ran 50 random vectors.
  \item The assertion was skipped.
\end{enumerate}
\end{qbox}
\begin{ansbox}
\textbf{Correct Answer: B}

Bounded Model Checking (BMC) unrolls the transition relation for $K$ steps. It guarantees no counterexample exists within depth $K$, but does not guarantee correctness for all time. A complete mathematical proof requires unbounded formal proof engines (e.g. Craig Interpolation, PDR/IC3, or Induction).
\end{ansbox}

\begin{qbox}{Q146: Formal Verification Vacuity and cover property Checks}{\badgeadv}
An engineer writes the assertion: \texttt{assert property (@(posedge clk) req |=> ack);}. During formal verification, the tool reports the assertion as \textbf{Passed}, but the design has a fatal bug where \texttt{req} is permanently tied to \texttt{0}. What happened, and how is it prevented?
\begin{enumerate}[label=\textbf{\Alph*.}]
  \item The formal tool has a bug.
  \item The assertion passed \textbf{vacuously} (because the antecedent \texttt{req} was never true, the implication \texttt{|=>} is automatically satisfied without ever checking \texttt{ack}); it must be accompanied by \texttt{cover property (@(posedge clk) req)} to prove that the antecedent was actively exercised.
  \item The formal tool automatically disables false assertions.
  \item Formal verification cannot detect tied signals.
\end{enumerate}
\end{qbox}
\begin{ansbox}
\textbf{Correct Answer: B}

In formal implication ($A \implies B$), if $A$ is false, the mathematical statement is true by definition (vacuous pass). To ensure assertions provide meaningful coverage, engineers write companion \texttt{cover property} checks to guarantee the antecedent triggers during verification.
\end{ansbox}

\begin{qbox}{Q147: assume property Over-Constraining Risks}{\badgeadv}
In Formal Property Verification, what is the danger of writing an over-constrained \texttt{assume property} on input signals?
\begin{enumerate}[label=\textbf{\Alph*.}]
  \item It increases simulation runtime by $10\times$.
  \item It can mask real design bugs (\textbf{Proof Vacuity / Over-Constraining}) by illegally restricting the formal engine from exploring valid real-world input stimulus scenarios.
  \item It causes synthesis to insert extra hardware.
  \item It turns the formal tool into an ATPG tool.
\end{enumerate}
\end{qbox}
\begin{ansbox}
\textbf{Correct Answer: B}

\texttt{assume} constraints limit the input state space explored by the formal solver. If an assumption is too restrictive (e.g. assuming an input can never arrive during a specific state), the formal tool will skip exploring those states, hiding genuine silicon bugs.
\end{ansbox}

\begin{qbox}{Q148: SVA disable iff Construct During System Reset}{\badgeadv}
Why must concurrent assertions always include \texttt{disable iff (!rst\_n)}?
\begin{enumerate}[label=\textbf{\Alph*.}]
  \item To save memory during compilation.
  \item During active reset, flip-flops are clearing and intermediate signals are invalid; \texttt{disable iff (!rst\_n)} immediately disables property evaluation while reset is asserted, preventing false assertion failures during system initialization.
  \item To allow assertions to run during scan shift mode.
  \item Because SVA does not support reset signals otherwise.
\end{enumerate}
\end{qbox}
\begin{ansbox}
\textbf{Correct Answer: B}

While reset is asserted, signals are indeterminate or in transition. Without \texttt{disable iff (!rst\_n)}, assertions would evaluate during reset, generating hundreds of false violation reports.
\end{ansbox}

\begin{qbox}{Q149: Formality Solver Aborts on Hardware Multipliers}{\badgeadv}
When Formality encounters a complex $64 \times 64$ hardware multiplier, the compare points often status as \textbf{ABORT / TIMEOUT}. Why does this happen, and what is the standard formal solution?
\begin{enumerate}[label=\textbf{\Alph*.}]
  \item Multipliers have high dynamic power.
  \item Multiplier Boolean equivalence creates an exponential explosion in Binary Decision Diagrams (BDD) and SAT-solver clauses; the solution is to Black-Box the multiplier architecture in both Reference and Implementation or verify the multiplier hierarchy independently.
  \item Formality cannot parse multiplication operators.
  \item The clock period is too short.
\end{enumerate}
\end{qbox}
\begin{ansbox}
\textbf{Correct Answer: B}

Formal equivalence of large non-linear arithmetic structures (multipliers, dividers) is NP-hard. BDD sizes grow exponentially with operand width ($O(2^N)$). Designers verify multipliers hierarchically or isolate them as black-boxes for system-level LEC.
\end{ansbox}

\begin{qbox}{Q150: SVA \$stable Handshake Data Preservation}{\badgeadv}
Which SVA assertion guarantees that once \texttt{valid} is asserted high, the 32-bit \texttt{data} bus must remain strictly unchanged until \texttt{ready} is asserted high?
\begin{enumerate}[label=\textbf{\Alph*.}]
  \item \texttt{assert property (@(posedge clk) (valid \&\& !ready) |=> \$stable(data));}
  \item \texttt{assert property (@(posedge clk) valid |-> \$changed(data));}
  \item \texttt{assert property (@(posedge clk) data == ready);}
  \item \texttt{assert property (@(posedge clk) ready |=> \$past(data));}
\end{enumerate}
\end{qbox}
\begin{ansbox}
\textbf{Correct Answer: A}

In standard AXI/Ready-Valid protocols, a sender cannot change data once offered until accepted. If \texttt{valid} is high and \texttt{ready} is low on cycle $N$, on the next cycle ($|=>$) \texttt{data} must satisfy \texttt{\$stable(data)} ($data_{current} == data_{previous}$).
\end{ansbox}

% ==============================================================================
\section{Design for Testability (DFT, ATPG \& JTAG)}
% ==============================================================================

\begin{qbox}{Q151: Primary Goal of DFT in ASIC Manufacturing}{\badgeeasy}
What is the primary objective of Design for Testability (DFT) in digital ASIC engineering?
\begin{enumerate}[label=\textbf{\Alph*.}]
  \item To verify that the RTL meets functional specifications.
  \item To detect physical manufacturing defects (e.g. Silicon shorts, opens, metal bridges, transistor defects) in post-fabrication packaged chips using automated test equipment (ATE).
  \item To eliminate power dissipation in standby mode.
  \item To measure the maximum operating frequency of the PLL.
\end{enumerate}
\end{qbox}
\begin{ansbox}
\textbf{Correct Answer: B}

DFT does not verify functional design intent; rather, it provides hardware structures (Scan Chains, BIST, Boundary Scan) that allow Automated Test Equipment (ATE) to detect physical semiconductor manufacturing flaws with high fault coverage.
\end{ansbox}

\begin{qbox}{Q152: Stuck-At-0 (SA0) Fault Physical Modeling}{\badgeeasy}
In structural fault modeling, what does a \textbf{Stuck-At-0 (SA0)} fault represent?
\begin{enumerate}[label=\textbf{\Alph*.}]
  \item A clock operating at zero frequency.
  \item A physical defect where a circuit node or transistor pin is permanently shorted to ground ($V_{SS}$ / logic \texttt{0}), regardless of input stimuli.
  \item A register that resets to 1.
  \item An open circuit on a power pin.
\end{enumerate}
\end{qbox}
\begin{ansbox}
\textbf{Correct Answer: B}

The classical Stuck-At Fault model abstracts physical manufacturing defects (such as bridging to ground or open gate connections) by modeling internal circuit nodes as permanently tied to logic \texttt{0} (SA0) or logic \texttt{1} (SA1).
\end{ansbox}

\begin{qbox}{Q153: Muxed-D Scan Flip-Flop Architecture}{\badgemod}
How does a standard \textbf{Muxed-D Scan Flip-Flop} modify a conventional D flip-flop?
\begin{enumerate}[label=\textbf{\Alph*.}]
  \item It adds an extra clock input for high-speed operation.
  \item It places a 2-to-1 multiplexer on the D input, controlled by a \texttt{Scan\_Enable (SE)} signal to select between functional data (\texttt{D}) when \texttt{SE=0} and scan shift test data (\texttt{SI}) when \texttt{SE=1}.
  \item It converts the flip-flop into a latch.
  \item It adds an integrated level shifter.
\end{enumerate}
\end{qbox}
\begin{ansbox}
\textbf{Correct Answer: B}

In Muxed-D scan architecture, setting \texttt{SE=1} connects all flip-flops into a serial shift register (\texttt{SI} to \texttt{SO}). Setting \texttt{SE=0} allows flip-flops to capture normal functional combinational logic outputs.
\end{ansbox}

\begin{qbox}{Q154: ATPG Scan Shift vs Capture Operational Sequence}{\badgemod}
What is the correct operational sequence during a standard ATPG scan test cycle on ATE?
\begin{enumerate}[label=\textbf{\Alph*.}]
  \item Capture $\rightarrow$ Shift $\rightarrow$ Reset
  \item \textbf{Scan Shift Mode (\texttt{SE=1})}: Shift test stimulus pattern into scan chains across $N$ clock cycles $\rightarrow$ \textbf{Capture Mode (\texttt{SE=0})}: Apply 1 functional clock pulse to capture combinational responses into flip-flops $\rightarrow$ \textbf{Scan Shift Mode (\texttt{SE=1})}: Shift out captured responses for signature comparison while shifting in the next pattern.
  \item Burn-in $\rightarrow$ Functional Execution $\rightarrow$ Power Off
  \item Shift $\rightarrow$ Continuous Free-Running Clocks for 1 hour
\end{enumerate}
\end{qbox}
\begin{ansbox}
\textbf{Correct Answer: B}

Scan testing converts the sequential testing problem into a simple combinational test: stimulus is shifted in serially (\texttt{SE=1}), captured through combinational logic in 1 cycle (\texttt{SE=0}), and shifted out for comparison (\texttt{SE=1}).
\end{ansbox}

\begin{qbox}{Q155: Fault Coverage vs Test Coverage Mathematical Calculation}{\badgemod}
Given a design with:
\begin{itemize}[leftmargin=1.5em]
  \item Total Faults = $10,000$
  \item Detected Faults = $9,500$
  \item Undetectable Faults (Tied/Redundant logic) = $500$
\end{itemize}
What are the \textbf{Fault Coverage} and \textbf{Test Coverage} percentages?
\begin{enumerate}[label=\textbf{\Alph*.}]
  \item Fault Coverage = $95\%$; Test Coverage = $100\%$
  \item Fault Coverage = $90\%$; Test Coverage = $95\%$
  \item Fault Coverage = $\frac{9500}{10000} = 95.0\%$; Test Coverage = $\frac{9500}{10000 - 500} = \frac{9500}{9500} = 100.0\%$
  \item Both are identically $95\%$.
\end{enumerate}
\end{qbox}
\begin{ansbox}
\textbf{Correct Answer: C}

\begin{itemize}[leftmargin=1.5em]
  \item $\text{Fault Coverage} = \frac{\text{Detected Faults}}{\text{Total Faults}} = \frac{9500}{10000} = 95.0\%$.
  \item $\text{Test Coverage} = \frac{\text{Detected Faults}}{\text{Total Faults} - \text{Undetectable Faults}} = \frac{9500}{9500} = 100.0\%$.
\end{itemize}
Test coverage measures the efficiency of the ATPG tool against physically testable faults.
\end{ansbox}

\begin{qbox}{Q156: Boundary Scan / JTAG (IEEE 1149.1) Dedicated Pins}{\badgemod}
What are the 4 mandatory dedicated test pins required by the IEEE 1149.1 JTAG Boundary Scan Standard?
\begin{enumerate}[label=\textbf{\Alph*.}]
  \item \texttt{CLK}, \texttt{RST}, \texttt{DIN}, \texttt{DOUT}
  \item \textbf{\texttt{TCK}} (Test Clock), \textbf{\texttt{TMS}} (Test Mode Select), \textbf{\texttt{TDI}} (Test Data In), \textbf{\texttt{TDO}} (Test Data Out) --- with optional \textbf{\texttt{TRST*}} (Test Reset).
  \item \texttt{VREF}, \texttt{VDD}, \texttt{VSS}, \texttt{SCAN\_EN}
  \item \texttt{SCL}, \texttt{SDA}, \texttt{INT}, \texttt{READY}
\end{enumerate}
\end{qbox}
\begin{ansbox}
\textbf{Correct Answer: B}

IEEE 1149.1 defines a 4-wire standard interface: \texttt{TCK} (clock), \texttt{TMS} (state machine control), \texttt{TDI} (serial data in), and \texttt{TDO} (serial data out), with an optional active-low asynchronous reset \texttt{TRST*}.
\end{ansbox}

\begin{qbox}{Q157: JTAG TAP Controller 16-State FSM Reset Mechanism}{\badgemod}
The JTAG TAP (Test Access Port) Controller state machine advances its states on which clock edge, and how is it reset to \texttt{Test-Logic-Reset} without using \texttt{TRST*}?
\begin{enumerate}[label=\textbf{\Alph*.}]
  \item Falling edge of \texttt{TCK}; by pulling \texttt{TDI} low for 5 cycles.
  \item \textbf{Rising edge of \texttt{TCK}}; by holding \textbf{\texttt{TMS = 1} for at least 5 consecutive \texttt{TCK} clock cycles}, guaranteeing return to \texttt{Test-Logic-Reset} from any of the 16 states.
  \item Both edges of \texttt{TCK}; by powering down the chip.
  \item High level on \texttt{TDO}.
\end{enumerate}
\end{qbox}
\begin{ansbox}
\textbf{Correct Answer: B}

The TAP controller state diagram is structured such that regardless of the starting state, holding \texttt{TMS=1} transitions the state machine back to \texttt{Test-Logic-Reset} within a maximum of 5 rising edges of \texttt{TCK}.
\end{ansbox}

\begin{qbox}{Q158: At-Speed Transition Delay Fault Testing (LOC vs LOS)}{\badgeadv}
What is the difference between \textbf{Launch-on-Shift (LOS)} and \textbf{Launch-on-Capture (LOC / Broadside)} in at-speed transition fault testing?
\begin{enumerate}[label=\textbf{\Alph*.}]
  \item LOS uses slow clock frequencies; LOC uses DC voltages.
  \item \textbf{LOS (Launch-on-Shift)}: The test transition is launched on the last shift pulse (\texttt{SE} must switch from 1 to 0 at rated gigahertz speed, requiring a timing-critical scan-enable tree); \textbf{LOC (Launch-on-Capture)}: The transition is launched by a functional clock pulse while \texttt{SE=0} (relaxed \texttt{SE} routing, but requires sequential ATPG).
  \item LOC does not support stuck-at faults.
  \item LOS requires analog test equipment.
\end{enumerate}
\end{qbox}
\begin{ansbox}
\textbf{Correct Answer: B}

At-speed transition testing requires two rated clock pulses: Launch and Capture. In LOS, the launch occurs on the final shift clock edge, requiring \texttt{Scan\_Enable} to deassert in less than one clock period ($T_{clk}$). LOC applies both pulses during functional capture mode (\texttt{SE=0}), eliminating the need for an at-speed \texttt{Scan\_Enable} network.
\end{ansbox}

\begin{qbox}{Q159: Scan Lockup Latches Across Clock Domains}{\badgeadv}
Why must a \textbf{Scan Lockup Latch (Active-Low Latch)} be inserted at the boundary where a scan chain crosses from Domain A (\texttt{CLK\_A}) to Domain B (\texttt{CLK\_B}) or between different clock edges?
\begin{enumerate}[label=\textbf{\Alph*.}]
  \item To invert the scan data.
  \item To eliminate \textbf{Hold Violations during Scan Shift Mode} caused by clock skew between the two domains; the negative-level latch holds data stable during the clock high phase, ensuring safe transfer even with severe inter-domain skew.
  \item To double the scan shift frequency.
  \item To compress ATPG test patterns.
\end{enumerate}
\end{qbox}
\begin{ansbox}
\textbf{Correct Answer: B}

During scan shift, flip-flops are connected in a direct shift chain with zero combinational logic ($T_{comb} = 0$). Any positive clock skew between registers causes immediate hold violations. An active-low lockup latch holds the output of the transmitter during the first half-cycle, providing a half-period of hold margin.
\end{ansbox}

\begin{qbox}{Q160: Memory BIST (MBIST) and March Algorithms}{\badgeadv}
Why are embedded SRAMs tested using \textbf{MBIST} engines running \textbf{March Algorithms (e.g. March C-)} rather than standard scan ATPG chains?
\begin{enumerate}[label=\textbf{\Alph*.}]
  \item Memory cells cannot be modeled with gate-level stuck-at models; March tests systematically write and read patterned sequences ($O(N)$ complexity) to detect complex memory-specific physical defects like transition faults, address decoder open faults, coupling faults between adjacent cells, and retention faults.
  \item Standard scan flip-flops cannot be fabricated inside SRAMs.
  \item MBIST requires zero silicon area.
  \item Both A and B are physically correct.
\end{enumerate}
\end{qbox}
\begin{ansbox}
\textbf{Correct Answer: D}

Embedded memories have dense transistor arrays that cannot accommodate internal scan chains. Dedicated hardware MBIST controllers run algorithmic March test patterns (e.g., $\Uparrow(w0), \Uparrow(r0,w1), \Downarrow(r1,w0)$) to detect cell coupling, address decoder, and pattern-sensitive defects.
\end{ansbox}

\begin{qbox}{Q161: Logic BIST (LBIST): PRPG and MISR Architecture}{\badgeadv}
In an autonomous Logic BIST (LBIST) architecture:
What are the roles of the \textbf{PRPG} and the \textbf{MISR}?
\begin{enumerate}[label=\textbf{\Alph*.}]
  \item PRPG measures temperature; MISR controls voltage.
  \item \textbf{PRPG (Pseudo-Random Pattern Generator)}: An LFSR that autonomously generates pseudo-random test stimulus vectors for scan chains; \textbf{MISR (Multiple-Input Signature Register)}: An LFSR/XOR compressor that compresses the captured output responses into a compact multi-bit signature for pass/fail comparison.
  \item PRPG is an external ATE pin; MISR is an internal clock divider.
  \item They replace the CPU instruction cache.
\end{enumerate}
\end{qbox}
\begin{ansbox}
\textbf{Correct Answer: B}

LBIST enables on-chip self-testing (critical for automotive ISO 26262 safety): a PRPG generates millions of pseudorandom patterns on-chip, and a MISR compresses output states into a final 32/64-bit signature compared against a golden ROM value.
\end{ansbox}

\begin{qbox}{Q162: Scan Compression (EDT / TestKompress)}{\badgeadv}
How does \textbf{Embedded Deterministic Test (Scan Compression)} achieve $50\times$ reduction in test time and ATE data volume?
\begin{enumerate}[label=\textbf{\Alph*.}]
  \item By dropping 50% of the test patterns.
  \item An on-chip \textbf{Decompressor} expands a small number of external ATE input channels into hundreds of internal short scan chains, while an on-chip \textbf{Compactor} compresses the hundreds of internal output chains back into a few ATE output pins, utilizing the fact that $>95\%$ of ATPG bits are don't-cares (\texttt{X}s).
  \item By increasing the ATE clock frequency to $100\,\text{GHz}$.
  \item By testing only 1 chip per wafer.
\end{enumerate}
\end{qbox}
\begin{ansbox}
\textbf{Correct Answer: B}

In standard ATPG, $>95\%$ of test bits are "don't cares" (\texttt{X}) needed only to propagate a few targeted faults. Scan compression uses linear feedback decompressors to broadcast deterministic values to internal scan chains, dramatically reducing test time and tester memory.
\end{ansbox}

\begin{qbox}{Q163: JTAG BYPASS Instruction Functionality}{\badgeadv}
What is the function of the mandatory JTAG instruction \textbf{\texttt{BYPASS}}?
\begin{enumerate}[label=\textbf{\Alph*.}]
  \item It disconnects the chip from the circuit board.
  \item It connects \texttt{TDI} directly to \texttt{TDO} through a single \textbf{1-bit Bypass Register} (with 1 cycle latency), allowing test data to bypass this IC quickly when testing other chips on the same PCB scan chain.
  \item It bypasses all setup timing checks in the chip.
  \item It bypasses the chip's internal voltage regulators.
\end{enumerate}
\end{qbox}
\begin{ansbox}
\textbf{Correct Answer: B}

When multiple chips share a JTAG chain on a PCB, testing a single target chip would be slowed down by shifting through hundreds of boundary scan cells in unrelated ICs. Setting unrelated chips to \texttt{BYPASS} shortens their internal shift path to a single 1-bit register.
\end{ansbox}

\begin{qbox}{Q164: Scan Chain Balancing Rationale}{\badgeadv}
Why must multiple internal scan chains be balanced to equal lengths (e.g. 10 chains of 500 flip-flops each, rather than 1 chain of 4500 and 9 of 50)?
\begin{enumerate}[label=\textbf{\Alph*.}]
  \item To balance the dynamic power across package pins.
  \item The total shift time for each test pattern is bounded by the \textbf{longest single scan chain}; unbalanced chains force shorter chains to be padded with dummy bits, wasting tester memory and dramatically increasing test time.
  \item Unbalanced chains cause clock skew during functional mode.
  \item Synthesis tools cannot route unbalanced chains.
\end{enumerate}
\end{qbox}
\begin{ansbox}
\textbf{Correct Answer: B}

All scan chains shift simultaneously during test. The ATE must supply shift cycles equal to the longest chain length ($L_{max}$). If one chain is 4500 bits and others are 50 bits, the ATE must execute 4500 shift cycles per pattern, wasting 90% of test bandwidth.
\end{ansbox}

\begin{qbox}{Q165: Undetectable Faults and Redundant Logic}{\badgeadv}
Consider the logic function $Y = A \cdot B + \overline{A} \cdot C + B \cdot C$.  
Why is the Stuck-At-0 fault on the input to term $B \cdot C$ classified by ATPG as an \textbf{Undetectable / Redundant Fault}?
\begin{enumerate}[label=\textbf{\Alph*.}]
  \item Because the cell is not connected to clock.
  \item The term $B \cdot C$ is a consensus term (algebraically redundant); no input stimulus vector can ever propagate the effect of a fault on this gate to primary output $Y$ without being masked by $A \cdot B$ or $\overline{A} \cdot C$.
  \item Because ATPG tools do not support 3-input equations.
  \item Because the output is always zero.
\end{enumerate}
\end{qbox}
\begin{ansbox}
\textbf{Correct Answer: B}

Redundant logic cannot be tested because its output can never uniquely control the circuit output independently of other paths. Redundant gates reduce test coverage metrics and waste area, though they are sometimes intentionally used to prevent glitches.
\end{ansbox}

\begin{qbox}{Q166: X-Bounding and X-Masking in Scan Compactors}{\badgeadv}
Why are uninitialized memory outputs, analog IP blocks, or asynchronous crossings hazardous to on-chip scan output compactors (MISRs), and how is this resolved?
\begin{enumerate}[label=\textbf{\Alph*.}]
  \item They consume excess static current.
  \item Unknown logic states (\texttt{X}s) corrupt the entire MISR signature calculation (since $X \oplus 0 = X$), destroying test validity; they must be blocked using \textbf{X-Bounding / X-Masking logic} before reaching the compactor.
  \item They cause physical short-circuits in the scan chains.
  \item They disable the JTAG TAP controller.
\end{enumerate}
\end{qbox}
\begin{ansbox}
\textbf{Correct Answer: B}

Compaction structures (MISRs or XOR trees) combine hundreds of bits. If even a single bit is \texttt{X} (unknown), the unknown state propagates through XOR gates and corrupts the entire signature. X-masking logic selectively forces known values onto suspect scan channels.
\end{ansbox}

\begin{qbox}{Q167: IDDQ Current Testing Principles}{\badgeadv}
What is the principle of \textbf{$I_{DDQ}$ Testing}, and what physical defects does it detect?
\begin{enumerate}[label=\textbf{\Alph*.}]
  \item It tests dynamic switching power at $5\,\text{GHz}$.
  \item It measures the quiescent steady-state supply current ($I_{DDQ}$) when all internal nodes are static; defective circuits with bridging shorts or gate-oxide breakdowns draw orders-of-magnitude higher quiescent current than nominal leakage.
  \item It tests radiation hardness in aerospace chips.
  \item It checks the resistance of bonding wires.
\end{enumerate}
\end{qbox}
\begin{ansbox}
\textbf{Correct Answer: B}

In a static CMOS circuit, steady-state current is purely leakage. If a bridging defect exists between two nodes driven to opposite values, a low-impedance short-circuit path from $V_{DD}$ to $V_{SS}$ draws massive current ($>100\,\mu\text{A}$), easily flagged by $I_{DDQ}$ monitoring.
\end{ansbox}

\begin{qbox}{Q168: IEEE 1500 Standard for Embedded Core Test}{\badgeadv}
What is the primary role of the \textbf{IEEE 1500 Standard} in modern multi-core SoC design?
\begin{enumerate}[label=\textbf{\Alph*.}]
  \item It provides high-speed PCIe communication.
  \item It defines an embedded modular test wrapper (Wrapper Instruction Register - WIR, Wrapper Boundary Cells - WBC) around individual reusable IP cores, enabling isolated modular testing of cores and SoC interconnects.
  \item It standardizes the USB protocol.
  \item It defines FPGA bitstream formats.
\end{enumerate}
\end{qbox}
\begin{ansbox}
\textbf{Correct Answer: B}

IEEE 1500 provides a standardized wrapper around complex modular IP cores (CPUs, GPUs, DSPs). It allows each core to be tested independently (Internal Test) and allows the surrounding interconnect between cores to be tested (External Test).
\end{ansbox}

% ==============================================================================
\section{Gate-Level Simulation (GLS) \& SDF}
% ==============================================================================

\begin{qbox}{Q169: Zero-Delay vs SDF-Annotated GLS}{\badgeeasy}
What is the difference between a \textbf{Zero-Delay Gate-Level Simulation} and a \textbf{Full-Timing SDF GLS}?
\begin{enumerate}[label=\textbf{\Alph*.}]
  \item Zero-delay runs on hardware emulators; SDF runs in SPICE.
  \item Zero-delay simulation evaluates logic gate functions with zero propagation delay across all gates and wires (used for initial functional netlist sanity checks); SDF GLS back-annotates precise cell delays, wire delays, and setup/hold timing checks from Standard Delay Format (\texttt{.sdf}) files.
  \item Zero-delay GLS includes crosstalk; SDF does not.
  \item There is no difference.
\end{enumerate}
\end{qbox}
\begin{ansbox}
\textbf{Correct Answer: B}

Zero-delay netlist simulation verifies that synthesis did not corrupt functional logic expressions. SDF back-annotated simulation uses physical timing data to verify dynamic operation, clock gating, asynchronous resets, and timing-dependent interfaces under real delays.
\end{ansbox}

\begin{qbox}{Q170: Verilog \$sdf\_annotate System Task}{\badgemod}
In a Verilog testbench, what does the following system task execute?
\begin{lstlisting}[style=verilogstyle]
$sdf_annotate("top_layout.sdf", u_dut, , "sdf.log", "MAXIMUM");
\end{lstlisting}
\begin{enumerate}[label=\textbf{\Alph*.}]
  \item It compiles the testbench into C++ code.
  \item It reads the timing delay file \texttt{top\_layout.sdf} and annotates the \textbf{Maximum} timing delays (slow-corner delays) onto the instantiated hierarchy \texttt{u\_dut}, logging errors to \texttt{sdf.log}.
  \item It disables all timing checks inside \texttt{u\_dut}.
  \item It measures dynamic power dissipation during simulation.
\end{enumerate}
\end{qbox}
\begin{ansbox}
\textbf{Correct Answer: B}

\texttt{\$sdf\_annotate} is the standard IEEE 1497 Verilog system task that parses \texttt{.sdf} files and back-annotates \texttt{IOPATH} cell delays, \texttt{INTERCONNECT} net delays, and \texttt{\$setup}/\texttt{\$hold} timing checks into the simulator's timing engine.
\end{ansbox}

\begin{qbox}{Q171: Simulator Response to \$setup and \$hold Violations}{\badgemod}
What occurs in an SDF-annotated gate-level simulation when a data transition violates the \texttt{\$setup} timing check specified inside a standard cell's \texttt{specify} block?
\begin{enumerate}[label=\textbf{\Alph*.}]
  \item The simulation immediately aborts with a segmentation fault.
  \item The simulator issues a timing violation warning in the console and drives the output pin \texttt{Q} of the affected flip-flop to an \textbf{unknown state (\texttt{X})} to model non-deterministic behavior.
  \item The simulator ignores the violation and continues normally.
  \item The clock frequency is automatically halved.
\end{enumerate}
\end{qbox}
\begin{ansbox}
\textbf{Correct Answer: B}

Standard cell Verilog models define specify blocks with \texttt{\$setup}, \texttt{\$hold}, \texttt{\$recovery}, and \texttt{\$removal} checks. When violated, the library cell model's notifier register toggles, causing the output \texttt{Q} to corrupt to \texttt{X} (unknown state).
\end{ansbox}

\begin{qbox}{Q172: Why GLS is Mandatory Alongside STA and LEC Signoff}{\badgemod}
If Static Timing Analysis (STA) exhaustively proves timing closure and Formality (LEC) proves functional equivalence, why is Gate-Level Simulation (GLS) still required before tapeout?
\begin{enumerate}[label=\textbf{\Alph*.}]
  \item To check that standard cell colors look correct.
  \item To verify asynchronous initialization sequences, reset deassertion recovery/removal behavior, dynamic clock switch behaviors, PLL lock sequences, and power-up/down UPF isolation sequences that cannot be fully verified by static analysis.
  \item Because STA tools cannot measure setup slack.
  \item To replace physical DRC checks.
\end{enumerate}
\end{qbox}
\begin{ansbox}
\textbf{Correct Answer: B}

STA is static and relies on design constraints (which may contain incorrect false path or case analysis waivers). GLS provides dynamic verification of multi-clock synchronization, power-up reset release, and complex mixed-signal handshakes in real time.
\end{ansbox}

\begin{qbox}{Q173: Simulation Arguments: +notimingchecks vs +nospecify}{\badgemod}
What is the effect of running a gate-level simulator with the argument \texttt{+notimingchecks}?
\begin{enumerate}[label=\textbf{\Alph*.}]
  \item It deletes the gate-level netlist.
  \item It executes the simulation using the cell and interconnect propagation delays, but disables all \texttt{\$setup}, \texttt{\$hold}, \texttt{\$recovery}, and \texttt{\$width} timing checks, preventing flip-flops from corrupting to \texttt{X}.
  \item It converts all clocks to zero delay.
  \item It bypasses all assertion checks.
\end{enumerate}
\end{qbox}
\begin{ansbox}
\textbf{Correct Answer: B}

\texttt{+notimingchecks} is commonly used during initial reset power-up phases to prevent uninitialized asynchronous reset assertions from generating cascades of false \texttt{X} states across the chip before clocks stabilize.
\end{ansbox}

\begin{qbox}{Q174: X-Pessimism in Combinational MUX Gates}{\badgeadv}
Consider a 2-to-1 MUX implemented with basic gates: \texttt{assign Y = (A \& S) | (B \& \textasciitilde S);}.  
If \texttt{A = 1}, \texttt{B = 1}, and select line \texttt{S = X} (unknown):
What does the gate-level simulator output for \texttt{Y}, and why is this \textbf{X-Pessimism}?
\begin{enumerate}[label=\textbf{\Alph*.}]
  \item Simulator outputs \texttt{Y = 1} (correct physical behavior).
  \item Simulator outputs \textbf{\texttt{Y = X}} (because $1 \& X = X$, $1 \& \sim X = X$, and $X | X = X$), even though in real silicon $Y$ is physically guaranteed to be \texttt{1} regardless of $S$.
  \item Simulator outputs \texttt{Y = 0}.
  \item Simulator crashes.
\end{enumerate}
\end{qbox}
\begin{ansbox}
\textbf{Correct Answer: B}

4-state HDL simulators evaluate gates individually without symbolic correlation. In reality, whether $S=0$ or $S=1$, $Y=1$. The simulator evaluates each branch as $X$ and their OR as $X$. This "X-Pessimism" causes false simulation lockups in GLS that do not exist in silicon.
\end{ansbox}

\begin{qbox}{Q175: Resolving X-Pessimism in GLS via Xprop}{\badgeadv}
How do verification teams overcome catastrophic X-Pessimism during gate-level simulation?
\begin{enumerate}[label=\textbf{\Alph*.}]
  \item By deleting all multiplexers.
  \item Using specialized simulator flags (e.g. Synopsys VCS \texttt{+vcs+xprop} or Cadence X-Optimism modes), modifying standard cell simulation models to handle correlated conditions, and initializing memories/registers with random deterministic values during reset.
  \item By forcing all signals to 1.
  \item By disabling the clock tree.
\end{enumerate}
\end{qbox}
\begin{ansbox}
\textbf{Correct Answer: B}

EDA tools provide X-propagation modes (\texttt{Xprop}) that analyze logic cones to determine if boolean outputs are independent of the unknown input, preventing artificial \texttt{X} propagation while preserving true uninitialized state tracking.
\end{ansbox}

\begin{qbox}{Q176: Glitch Filtering: Inertial vs Transport Delay Models}{\badgeadv}
In gate-level SDF simulation, what is the difference between \textbf{Inertial Delay} and \textbf{Transport Delay}?
\begin{enumerate}[label=\textbf{\Alph*.}]
  \item Inertial delay applies only to wires; Transport delay applies only to gates.
  \item \textbf{Inertial Delay} (default in Verilog gates): Pulses shorter than the gate's propagation delay are rejected and filtered out as glitches; \textbf{Transport Delay} (\texttt{+transport\_path\_delays}): Propagates all pulses through wires and interconnects regardless of duration, accurately modeling narrow glitch propagation across physical routes.
  \item Transport delay doubles the pulse width.
  \item Inertial delay converts pulses into sine waves.
\end{enumerate}
\end{qbox}
\begin{ansbox}
\textbf{Correct Answer: B}

Standard gate models use inertial delay, which suppresses pulses narrower than the cell propagation delay. Physical metal interconnects behave with transport delay (pure delay line). Running GLS with \texttt{+transport\_path\_delays} ensures narrow glitches on routing tracks reach downstream gates.
\end{ansbox}

\begin{qbox}{Q177: SDF IOPATH vs INTERCONNECT Delays}{\badgeadv}
In an SDF file, what is the distinction between \texttt{(IOPATH in out (0.15))} and \texttt{(INTERCONNECT u1/out u2/in (0.35))}?
\begin{enumerate}[label=\textbf{\Alph*.}]
  \item \texttt{IOPATH} models wire delay; \texttt{INTERCONNECT} models transistor delay.
  \item \textbf{\texttt{IOPATH}} specifies the internal cell delay from an input pin to an output pin of a specific standard cell; \textbf{\texttt{INTERCONNECT}} specifies the net/wire propagation delay from a driving cell output pin to a load cell input pin across physical metal routes.
  \item \texttt{INTERCONNECT} is only used for analog pins.
  \item \texttt{IOPATH} is deprecated in IEEE 1497.
\end{enumerate}
\end{qbox}
\begin{ansbox}
\textbf{Correct Answer: B}

SDF separates cell delay from interconnect delay. \texttt{IOPATH} records the internal switching delay of the standard cell arc, while \texttt{INTERCONNECT} records the distributed RC wire delay extracted from physical layout parasitics (SPEF).
\end{ansbox}

\begin{qbox}{Q178: Negative Setup and Hold Times in Characterized Cells}{\badgeadv}
When an SDF file annotates a \textbf{Negative Hold Time} (e.g. \texttt{(HOLD D (posedge CK) (-0.05))}) on a flip-flop:
What does this physically mean, and how does the simulator handle it?
\begin{enumerate}[label=\textbf{\Alph*.}]
  \item The flip-flop operates backwards in time.
  \item An internal delay on the data path inside the standard cell exceeds the delay on the internal clock path; data can physically change $0.05\,\text{ns}$ \textit{before} the external clock pin transition without corrupting the captured state.
  \item The flip-flop has a manufacturing defect.
  \item Negative numbers are illegal in SDF and cause simulation aborts.
\end{enumerate}
\end{qbox}
\begin{ansbox}
\textbf{Correct Answer: B}

Standard cells are characterized at their external boundary pins. If internal buffering delays the data line more than the clock line, data does not need to be held until the external clock edge---it can change slightly before it. Simulators handle negative values using negative timing check algorithms (\texttt{\$setuphold}).
\end{ansbox}

\begin{qbox}{Q179: Pulse Rejection and Error Limits (+pulse\_r / +pulse\_e)}{\badgeadv}
What is the purpose of specifying simulator arguments \texttt{+pulse\_r/0} and \texttt{+pulse\_e/100} during SDF gate-level simulation?
\begin{enumerate}[label=\textbf{\Alph*.}]
  \item To control the simulator's CPU thread count.
  \item To define pulse rejection thresholds: Any pulse with width $< 0\%$ of cell delay is rejected (ignored), while any pulse with width between $0\%$ and $100\%$ of cell delay is flagged as a hazard and driven to \textbf{\texttt{X} (Error)} to expose potential glitch corruption.
  \item To scale the clock period.
  \item To measure power dissipation.
\end{enumerate}
\end{qbox}
\begin{ansbox}
\textbf{Correct Answer: B}

Pulse control options handle narrow transition spikes. Setting \texttt{pulse\_e/100} converts any partial pulse (glitch narrower than cell delay) into an unknown \texttt{X} state, highlighting potential race conditions in simulation waveforms.
\end{ansbox}

\begin{qbox}{Q180: Diagnostic Workflow for GLS Reset Glitches}{\badgeadv}
During full-timing SDF GLS of a multi-million gate SoC, the entire chip corrupts to \texttt{X} within the first 100ns of simulation. What is the standard diagnostic workflow to identify the root cause?
\begin{enumerate}[label=\textbf{\Alph*.}]
  \item Immediately re-synthesize the design with higher optimization.
  \item Trace back the origin of the first \texttt{X} transition in the waveform viewer (e.g. Verdi / DVE), inspect timing violation logs at time $T_0$ to check for uninitialized flip-flops or recovery/removal violations on the asynchronous reset tree, and verify whether \texttt{+notimingchecks} was enabled during the reset asserting window.
  \item Increase the supply voltage in the testbench.
  \item Disable all clock generators.
\end{enumerate}
\end{qbox}
\begin{ansbox}
\textbf{Correct Answer: B}

In GLS, a single flip-flop entering \texttt{X} rapidly cascades across the entire chip. Tracing back through the schematic/waveform to the very first signal that went to \texttt{X} reveals whether it was triggered by a real timing violation (e.g. reset recovery failure) or an uninitialized memory array.
\end{ansbox}

% ==============================================================================
\section{Physical Design \& Place \& Route (PnR)}
% ==============================================================================

\begin{qbox}{Q181: Floorplanning Core Utilization Calculation}{\badgeeasy}
If a chip has a total standard cell area of $2.0\,\text{mm}^2$ and the designer defines a core area of $4.0\,\text{mm}^2$:
What is the \textbf{Core Utilization} factor?
\begin{enumerate}[label=\textbf{\Alph*.}]
  \item $25\%$
  \item $50\%$
  \item $75\%$
  \item $200\%$
\end{enumerate}
\end{qbox}
\begin{ansbox}
\textbf{Correct Answer: B}

$\text{Core Utilization} = \frac{\text{Total Cell Area}}{\text{Total Core Area}} = \frac{2.0\,\text{mm}^2}{4.0\,\text{mm}^2} = 0.50 = 50\%$.
\end{ansbox}

\begin{qbox}{Q182: Sequential Stages of the Physical Design Flow}{\badgeeasy}
What is the standard sequential execution flow of Physical Design (PnR)?
\begin{enumerate}[label=\textbf{\Alph*.}]
  \item Routing $\rightarrow$ Placement $\rightarrow$ Floorplanning $\rightarrow$ Tapeout
  \item Floorplanning / Power Planning $\rightarrow$ Standard Cell Placement $\rightarrow$ Clock Tree Synthesis (CTS) $\rightarrow$ Detail Routing $\rightarrow$ Physical Verification (DRC/LVS) \& Signoff Extraction
  \item CTS $\rightarrow$ ATPG $\rightarrow$ Synthesis $\rightarrow$ Fabrication
  \item Packaging $\rightarrow$ Floorplanning $\rightarrow$ Extraction
\end{enumerate}
\end{qbox}
\begin{ansbox}
\textbf{Correct Answer: B}

Physical design begins with die/core floorplanning and power grid construction, places standard cells, synthesizes balanced clock distribution trees (CTS), routes all signal nets, and finishes with signoff extraction (SPEF) and physical DRC/LVS verification.
\end{ansbox}

\begin{qbox}{Q183: Macro Placement Halos and Keepout Margins}{\badgemod}
Why are \textbf{Placement Halos (Keepout Margins)} placed around large hard IP macros (e.g. SRAMs, PLLs) during floorplanning?
\begin{enumerate}[label=\textbf{\Alph*.}]
  \item To prevent the macro from consuming static power.
  \item To reserve physical routing and placement channels around the macro perimeter, preventing standard cells from being placed too close to macro pins, avoiding pin access congestion and noise coupling.
  \item To isolate analog grounds from digital grounds.
  \item To allow heat dissipation only.
\end{enumerate}
\end{qbox}
\begin{ansbox}
\textbf{Correct Answer: B}

Placement halos create an exclusion boundary around memory blocks where no standard cells can be placed. This ensures adequate space for power routing rings and provides routing channels for the high density of pins entering and exiting the macro.
\end{ansbox}

\begin{qbox}{Q184: Static vs Dynamic IR Drop on Power Grids}{\badgemod}
What is the distinction between \textbf{Static IR Drop} and \textbf{Dynamic IR Drop} on the on-chip power distribution network (PDN)?
\begin{enumerate}[label=\textbf{\Alph*.}]
  \item Static IR drop occurs in sleep mode; Dynamic occurs in test mode.
  \item \textbf{Static IR Drop} is the steady-state average DC voltage drop ($\Delta V = I_{avg} \times R_{grid}$) caused by wire resistance; \textbf{Dynamic IR Drop} is the localized transient voltage drop caused by simultaneous high-frequency switching current surges ($\Delta V = I(t) \cdot R + L \cdot \frac{di}{dt}$) on active clock edges.
  \item Static IR drop only affects I/O pads.
  \item Dynamic IR drop can be eliminated by removing decoupling capacitors.
\end{enumerate}
\end{qbox}
\begin{ansbox}
\textbf{Correct Answer: B}

Static IR drop depends on average DC current consumption and metal rail resistance. Dynamic IR drop is a high-frequency transient drop caused by thousands of flip-flops switching simultaneously on a clock edge ($L \cdot di/dt$).
\end{ansbox}

\begin{qbox}{Q185: Decoupling Capacitors (Decap Cells) Role}{\badgemod}
How do \textbf{Decoupling Capacitor (Decap)} cells mitigate dynamic IR drop in high-speed digital designs?
\begin{enumerate}[label=\textbf{\Alph*.}]
  \item They disconnect the power supply during switching.
  \item They act as local high-frequency charge reservoirs placed close to switching logic cells; when large current surges occur on clock edges, Decap cells supply instantaneous charge locally, preventing power grid voltage droop.
  \item They increase wire resistance.
  \item They eliminate static leakage current.
\end{enumerate}
\end{qbox}
\begin{ansbox}
\textbf{Correct Answer: B}

Standard power supply pads are too far away (high inductance) to supply instantaneous sub-nanosecond switching currents. Decap cells placed in empty floorplan sites store local charge and supply transient currents directly to neighboring gates.
\end{ansbox}

\begin{qbox}{Q186: Clock Tree Synthesis (CTS) Optimization Targets}{\badgemod}
What are the primary goals of \textbf{Clock Tree Synthesis (CTS)}?
\begin{enumerate}[label=\textbf{\Alph*.}]
  \item To minimize cell area and remove all buffers.
  \item To build a balanced distribution network (using H-Tree, Clock Mesh, or Balanced Tree structures) that distributes the clock signal to millions of sequential elements with \textbf{Minimum Skew}, \textbf{Controlled Latency (Insertion Delay)}, and \textbf{Symmetrical Transition Slews}.
  \item To convert flip-flops into scan chains.
  \item To eliminate dynamic power completely.
\end{enumerate}
\end{qbox}
\begin{ansbox}
\textbf{Correct Answer: B}

CTS inserts balanced buffer/inverter trees to deliver clock edges to all flip-flops simultaneously, minimizing spatial clock skew ($T_{capture} - T_{launch}$), maintaining sharp rise/fall slews, and controlling total insertion delay.
\end{ansbox}

\begin{qbox}{Q187: Antenna Effect and Diode Shunt Insertion}{\badgeadv}
During semiconductor fabrication, what is the \textbf{Antenna Effect}, and how is it resolved in Physical Design?
\begin{enumerate}[label=\textbf{\Alph*.}]
  \item Metal wires picking up 5G radio signals.
  \item Long exposed metal tracks act as antennas during plasma etching, collecting electrostatic charge; the accumulated charge discharges through thin MOSFET gate oxides, permanently destroying the transistor dielectric.  
  \textbf{Fix}: Metal layer jumping (routing up to higher metals) or inserting \textbf{Antenna Diodes} near the gate to safely shunt charge to the substrate.
  \item Wires vibrating due to thermal heat.
  \item Clock trees radiating electromagnetic noise.
\end{enumerate}
\end{qbox}
\begin{ansbox}
\textbf{Correct Answer: B}

During plasma etching, long metal routes accumulate charge. If connected to a transistor gate without a source/drain diffusion path to bleed charge, the voltage exceeds gate oxide breakdown limits. Inserting reverse-biased antenna diodes provides a low-impedance discharge path to ground.
\end{ansbox}

\begin{qbox}{Q188: Electromigration (EM) Physics and Multi-Cut Vias}{\badgeadv}
What causes \textbf{Electromigration (EM)} in metal interconnects, and what are its physical consequences?
\begin{enumerate}[label=\textbf{\Alph*.}]
  \item High voltage causing dielectric breakdown.
  \item High directional direct-current density ($J > J_{max}$) causes moving electrons to physically transfer momentum to metal atoms ("electron wind"), moving metal ions over time to create \textbf{Voids (open circuits)} and \textbf{Hillocks/Whiskers (short circuits to adjacent tracks)}.  
  \textbf{Fix}: Widen metal wires and insert multi-cut vias.
  \item Thermal expansion cracking the silicon substrate.
  \item Chemical corrosion from packaging materials.
\end{enumerate}
\end{qbox}
\begin{ansbox}
\textbf{Correct Answer: B}

Electromigration is the physical transport of material in a conductor due to the momentum transfer between electrons and metal ions. Over time, high current densities create voids (leading to open circuits) or hillocks (leading to shorts). Widening wires reduces current density $J = I/A$.
\end{ansbox}

\begin{qbox}{Q189: Global Routing vs Detail Routing in PnR Tools}{\badgeadv}
In PnR tools (Innovus / ICC2), what is the difference between \textbf{Global Routing} and \textbf{Detail Routing}?
\begin{enumerate}[label=\textbf{\Alph*.}]
  \item Global routing routes clocks; Detail routing routes resets.
  \item \textbf{Global Routing}: Partitions the core into global routing cells (G-cells) and plans coarse capacity/demand usage to identify congestion hotspots without placing actual metal shapes; \textbf{Detail Routing}: Instantiates physical metal tracks, assigns exact routing layers, and places vias to meet all foundry DRC rules.
  \item Detail routing runs before placement.
  \item Global routing is performed in SPICE.
\end{enumerate}
\end{qbox}
\begin{ansbox}
\textbf{Correct Answer: B}

Global routing models the chip as a grid of bins (G-cells) to optimize wire topology and avoid congestion. Detailed routing follows the global routing guide to layout actual polygon metal paths, resolving design rule constraints (DRC) like spacing and width.
\end{ansbox}

\begin{qbox}{Q190: Non-Default Routing Rules (NDR) for Clock Nets}{\badgeadv}
Why are \textbf{Non-Default Routing Rules (NDR: 2W2S - Double Width, Double Spacing)} applied to clock distribution nets during CTS and routing?
\begin{enumerate}[label=\textbf{\Alph*.}]
  \item To make clock wires look visually distinct in GUI.
  \item \textbf{Double Width (2W)} reduces wire resistance, minimizing clock insertion delay; \textbf{Double Spacing (2S)} dramatically cuts coupling capacitance ($C_c$) to adjacent tracks, eliminating crosstalk delay jitter and noise glitches.
  \item To double the clock operating frequency.
  \item To prevent scan chain shifting.
\end{enumerate}
\end{qbox}
\begin{ansbox}
\textbf{Correct Answer: B}

Clock nets toggle every single cycle and are highly sensitive to skew and jitter. Doubling wire width reduces resistance ($R$), while doubling wire spacing reduces coupling capacitance to neighboring aggressors, isolating the clock from crosstalk.
\end{ansbox}

\begin{qbox}{Q191: Well Tap Cells vs Standard Filler Cells}{\badgeadv}
What are the distinct physical functions of \textbf{Well Tap Cells} and \textbf{Standard Filler Cells}?
\begin{enumerate}[label=\textbf{\Alph*.}]
  \item Filler cells provide power; Tap cells provide clocks.
  \item \textbf{Well Tap Cells}: Connect $V_{DD}$ to the N-well and $V_{SS}$ to the P-substrate at regular maximum distance intervals to prevent CMOS \textbf{Latch-Up}; \textbf{Filler Cells}: Inserted into empty spaces in placement rows to ensure continuous N-well/P-well diffusion continuity and layer density for DRC compliance.
  \item Both are functional logic gates.
  \item Tap cells are only used for scan chains.
\end{enumerate}
\end{qbox}
\begin{ansbox}
\textbf{Correct Answer: B}

Tapless standard cell libraries require separate Well Tap cells placed at maximum pitch intervals (e.g. Every $30\,\mu\text{m}$) to bias wells and prevent parasitic PNPN SCR latch-up. Filler cells contain no transistors; they maintain layer density and well continuity in empty floorplan gaps.
\end{ansbox}

\begin{qbox}{Q192: Physical DRC vs LVS Signoff Verification}{\badgeadv}
What is the difference between \textbf{DRC (Design Rule Checking)} and \textbf{LVS (Layout Versus Schematic)} signoff checks?
\begin{enumerate}[label=\textbf{\Alph*.}]
  \item DRC checks timing; LVS checks power.
  \item \textbf{DRC}: Verifies that the physical layout polygons satisfy all foundry geometric manufacturing rules (minimum metal widths, spacings, enclosure, density); \textbf{LVS}: Extracts the transistor-level circuit from physical layout polygons and proves 1-to-1 schematic/netlist connectivity and device sizing equivalence.
  \item LVS is run only on PCBs.
  \item DRC can be skipped if timing passes.
\end{enumerate}
\end{qbox}
\begin{ansbox}
\textbf{Correct Answer: B}

DRC guarantees the layout can be manufactured without fabrication defects (shorts, opens, lithography errors). LVS extracts devices and nets from the physical layout and compares them against the post-synthesis gate-level netlist to prove correct physical connectivity.
\end{ansbox}

\begin{qbox}{Q193: Parasitic Extraction and SPEF Format}{\badgeadv}
What is \textbf{SPEF (Standard Parasitic Exchange Format)}, and how is it used in the design flow?
\begin{enumerate}[label=\textbf{\Alph*.}]
  \item A format for describing testbenches.
  \item An IEEE standard ASCII format that contains extracted parasitic wire resistances ($R$), grounded capacitances ($C_g$), and cross-coupling capacitances ($C_c$) for every routed net in the layout, used by PrimeTime to perform accurate signoff STA.
  \item A floorplan description file.
  \item An FPGA programming bitstream.
\end{enumerate}
\end{qbox}
\begin{ansbox}
\textbf{Correct Answer: B}

After detailed routing, extraction tools (e.g. StarRC) calculate exact wire parasitics based on physical metal shapes and dielectric layers, formatting them as a \texttt{.spef} file for back-annotation into PrimeTime for final timing signoff.
\end{ansbox}

\begin{qbox}{Q194: Engineering Change Orders (ECO) Methodology}{\badgeadv}
Late in the physical design cycle, a minor RTL logic bug is discovered after full routing:  
Why is an \textbf{ECO (Engineering Change Order)} flow preferred over restarting full synthesis and PnR?
\begin{enumerate}[label=\textbf{\Alph*.}]
  \item Re-running synthesis and PnR is illegal under foundry rules.
  \item An ECO modifies only a few gates using pre-placed \textbf{Spare Cells} or minimal local cell additions, preserving the fully closed timing, clock trees, and layout routing of 99.9% of the chip, saving weeks of engineering effort.
  \item ECOs run without an SDC file.
  \item ECOs eliminate DRC rules.
\end{enumerate}
\end{qbox}
\begin{ansbox}
\textbf{Correct Answer: B}

Full re-implementation causes chaotic disruption to placed cells and clock trees ("timing closure roulette"). An ECO netlist patch modifies only the affected logic cone and routes minimal local wires, preserving all existing closed timing.
\end{ansbox}

% ==============================================================================
\section{FPGA Architecture \& ASIC vs FPGA}
% ==============================================================================

\begin{qbox}{Q195: Look-Up Table (LUT) Combinational Architecture}{\badgeeasy}
In modern FPGA architectures (e.g. AMD/Xilinx UltraScale+, Intel Stratix), what basic hardware structure implements arbitrary combinational logic functions?
\begin{enumerate}[label=\textbf{\Alph*.}]
  \item Standard cell NAND gates.
  \item \textbf{SRAM-based Look-Up Tables (LUTs)} (e.g. 6-input LUTs that can implement any arbitrary 6-input Boolean function as a $64 \times 1$-bit memory with a 64-to-1 MUX).
  \item Hardwired microprocessors.
  \item Dynamic RAM capacitors.
\end{enumerate}
\end{qbox}
\begin{ansbox}
\textbf{Correct Answer: B}

FPGAs do not have individual AND/OR gates. Combinational logic is implemented using SRAM-based Look-Up Tables (LUTs). A 6-input LUT contains 64 SRAM bits initialized with the truth table of the Boolean function, addressed by the 6 inputs.
\end{ansbox}

\begin{qbox}{Q196: FPGA vs ASIC Commercial and Technical Trade-Offs}{\badgeeasy}
What are the primary commercial and technical trade-offs of an \textbf{FPGA} compared to a custom \textbf{ASIC}?
\begin{enumerate}[label=\textbf{\Alph*.}]
  \item FPGA has higher unit cost at high volume, lower maximum frequency, and higher power dissipation, but zero Non-Recurring Engineering (NRE) mask costs and instant time-to-market with re-programmability.
  \item ASIC has zero upfront design cost.
  \item FPGA is faster and uses less silicon area than ASIC.
  \item FPGAs cannot implement sequential circuits.
\end{enumerate}
\end{qbox}
\begin{ansbox}
\textbf{Correct Answer: A}

ASICs require millions of dollars in upfront mask/NRE costs but have low per-unit silicon cost, maximum performance, and minimal power at high volume. FPGAs have zero mask cost and are reconfigurable, but have higher per-unit cost and lower power/area efficiency.
\end{ansbox}

\begin{qbox}{Q197: Hard IP vs Soft IP in FPGAs}{\badgemod}
Which of the following is an example of \textbf{Hard IP} in a modern FPGA?
\begin{enumerate}[label=\textbf{\Alph*.}]
  \item A synthesizable UART RTL module written in Verilog.
  \item Dedicated silicon blocks fabricated directly on the die, such as \textbf{DSP48 Slices}, \textbf{Block RAMs (BRAM / UltraRAM)}, \textbf{PCIe Hard Controllers}, and \textbf{Multi-Gigabit SerDes Transceivers}.
  \item A software C-program running on a host PC.
  \item An unconstrained clock net.
\end{enumerate}
\end{qbox}
\begin{ansbox}
\textbf{Correct Answer: B}

Soft IP consists of synthesizable RTL code mapped into generic LUTs and flip-flops. Hard IP refers to custom, dedicated silicon blocks (BRAMs, DSP arithmetic units, high-speed transceivers) embedded in the FPGA fabric for optimal performance.
\end{ansbox}

\begin{qbox}{Q198: Dedicated Global Clock Buffers (BUFG)}{\badgemod}
Why must all high-fanout clock signals in an FPGA design be routed through dedicated \textbf{Global Clock Buffers (e.g. BUFG, BUFGCE)}?
\begin{enumerate}[label=\textbf{\Alph*.}]
  \item General-purpose FPGA interconnect routing has high resistance and unpredictable segment delays, which would cause massive clock skew and hold violations; BUFG buffers drive dedicated, low-skew global clock spines.
  \item General routing fabric cannot carry electrical voltage.
  \item BUFGs convert digital clocks to analog waveforms.
  \item BUFGs eliminate the need for setup timing checks.
\end{enumerate}
\end{qbox}
\begin{ansbox}
\textbf{Correct Answer: A}

General FPGA routing fabric consists of segmented pass-transistor routing switches that introduce large, non-uniform delays. Dedicated clock buffers (\texttt{BUFG}) drive specialized, low-impedance clock distribution networks with minimal skew across the die.
\end{ansbox}

\begin{qbox}{Q199: Synchronous vs Asynchronous Reset Strategy in FPGAs}{\badgemod}
Why do FPGA vendors (e.g. AMD/Xilinx) strongly recommend \textbf{Synchronous Resets} over Asynchronous Resets for FPGA design, unlike standard ASIC practice?
\begin{enumerate}[label=\textbf{\Alph*.}]
  \item Asynchronous resets do not work in FPGAs.
  \item FPGA logic slices (CLBs) have dedicated synchronous control set/reset pins; using synchronous resets allows the synthesis tool to pack control logic efficiently into dedicated flip-flop control sets and LUT inputs without wasting routing resources.
  \item Synchronous resets eliminate all setup checks.
  \item Asynchronous resets double dynamic power.
\end{enumerate}
\end{qbox}
\begin{ansbox}
\textbf{Correct Answer: B}

FPGA flip-flops inside Configurable Logic Blocks (CLBs) share common control sets (clock enable, reset). Asynchronous resets prevent optimal register packing and restrict the synthesis tool from merging reset conditions into LUT logic.
\end{ansbox}

\begin{qbox}{Q200: Hazardous Nature of Inferred Latches in FPGAs}{\badgeadv}
Why is an accidental inferred latch far more damaging in an \textbf{FPGA} implementation than in an ASIC?
\begin{enumerate}[label=\textbf{\Alph*.}]
  \item FPGAs do not possess physical latch cells in their standard slices; the synthesis tool must either emulate the latch using a slice flip-flop with feedback or construct a combinational loop across multiple LUTs, consuming excessive resources, creating severe timing hazards, and causing static timing analysis tools to fail.
  \item Latches cause the FPGA bitstream to erase itself.
  \item Latches reverse the polarity of the power supply.
  \item Latches disable the JTAG programming interface.
\end{enumerate}
\end{qbox}
\begin{ansbox}
\textbf{Correct Answer: A}

FPGA slices are optimized for synchronous edge-triggered D flip-flops and LUTs. Inferred latches require either using flip-flops in non-standard latch modes (which disables surrounding registers in the slice) or creating combinational loops through LUT routing, which leads to race conditions and timing closure failures.
\end{ansbox}

\end{document}
